% Calcul intégral — Mathématiques Tle D — Cours signé OnBuch+
% Programme officiel MINESEC. Compiler avec : tectonic M08-calcul-integral.tex
\newcommand{\DOCMATIERE}{Mathématiques}
\newcommand{\DOCNIVEAU}{Tle D}
\newcommand{\DOCMODULE}{Relations et opérations fondamentales dans l'ensemble des nombres réels}
\newcommand{\DOCLECON}{08}
\newcommand{\DOCTITRE}{Calcul intégral}
% OnBuch+ — préambule « cours signés » ENRICHI (compilé avec tectonic / XeLaTeX)
% Système visuel « Pop » d'OnBuch+ : fond crème (jamais blanc), bordures encre
% épaisses, ombres « dures » décalées, titres Archivo Black en capitales.
% Version MATHÉMATIQUES : graphes (pgfplots/tikz), boîtes pédagogiques
% (définition, propriété, méthode, attention, le savais-tu, exercice résolu,
% exercice, TP, bilan), page de garde et signature OnBuch+ ancrée en bas.
%
% Macros attendues dans main.tex AVANT \input{preamble} :
%   \DOCMATIERE  \DOCNIVEAU  \DOCMODULE  \DOCLECON (numéro)  \DOCTITRE
\documentclass[11pt]{article}

\usepackage{fontspec}
\usepackage[french]{babel}
\usepackage[a4paper,margin=2cm,top=3.4cm,bottom=2.8cm,headheight=1.4cm,footskip=1.3cm]{geometry}
\usepackage{amsmath,amssymb,amsfonts}
\usepackage{mathtools} % \xmapsto, \coloneqq... souvent utilisés par les modèles
\usepackage{cancel} % simplifications barrées
% \abs{z} (module, valeur absolue) et \norm{u} (norme), utilisés par les modèles
\providecommand{\abs}{}\let\abs\relax\DeclarePairedDelimiter{\abs}{\lvert}{\rvert}
\providecommand{\norm}{}\let\norm\relax\DeclarePairedDelimiter{\norm}{\lVert}{\rVert}
\usepackage[table]{xcolor} % [table] : fournit aussi \rowcolors (alternance de lignes)
\usepackage{pagecolor}
\usepackage{tikz}
\usetikzlibrary{arrows.meta,positioning,calc,shapes.geometric,shapes.misc,shapes.arrows,decorations.pathmorphing,decorations.markings,patterns,shadows,mindmap,trees,fit,backgrounds,matrix,intersections}
\usepackage{pgfplots}
\usepackage{pgf-pie} % diagrammes circulaires \pie (statistiques)
\pgfplotsset{compat=1.18}
\usepgfplotslibrary{fillbetween,groupplots} % aires entre courbes (calcul intégral)
\usepackage{enumitem}
\usepackage{fancyhdr}
% [most] charge toutes les bibliothèques tcolorbox (listings, minted, raster,
% documentation...) et gonfle inutilement la mémoire TeX sur un document long
% (~300 boîtes) : on ne charge que ce qui est réellement utilisé.
\usepackage[skins,breakable]{tcolorbox}
\usepackage{graphicx}
\usepackage{colortbl}
\usepackage{booktabs}
\usepackage{tabularx}
\usepackage{diagbox} % case d'en-tête diagonale des tableaux à double entrée
% Colonnes extensibles centrée / gauche / droite (C, L, R), souvent utilisées par les modèles
\newcolumntype{C}{>{\centering\arraybackslash}X}
\newcolumntype{L}{>{\raggedright\arraybackslash}X}
\newcolumntype{R}{>{\raggedleft\arraybackslash}X}
\usepackage{array}
\usepackage{multicol}
\usepackage{siunitx}
% parse-numbers=false : accepte aussi \SI{2,5 \times 10^{-3}}{...}
\sisetup{locale=FR,output-decimal-marker={,},group-minimum-digits=4,parse-numbers=false}
\DeclareSIUnit\ohm{\ensuremath{\symup{\Omega}}} % U+2126 absent de la police
\usepackage{qrcode}
% \degree : commande fréquemment utilisée par les modèles pour un angle ou une
% température en dehors de \SI{}{} (non fournie par les paquets chargés).
\providecommand{\degree}{^{\circ}}
% \qed (fin de démonstration), utilisé par les modèles sans amsthm
\providecommand{\qed}{\hfill$\square$}
% \barre : double barre des valeurs interdites dans un tableau de variations (tkz-tab)
\providecommand{\barre}{\ensuremath{\|}}
% environnement proof (amsthm non chargé) utilisé par les modèles
\ifdefined\proof\else\newenvironment{proof}[1][Démonstration]{\par\noindent\textit{#1.}\ }{\qed\par}\fi
\usepackage{lastpage}
\usepackage{hyperref}
\hypersetup{hidelinks}

% ============================================================
% Palette « Pop » OnBuch+ — mêmes hex que la classe P (plus_ui.dart)
% ============================================================
\definecolor{poporange}{HTML}{F59321}
\definecolor{popdark}{HTML}{E07A0C}
\definecolor{popcream}{HTML}{FFF6E8}
\definecolor{popink}{HTML}{1C1714}
\definecolor{poppurple}{HTML}{7A5AE0}
\definecolor{popgreen}{HTML}{2E9E6A}
\definecolor{poppink}{HTML}{E0457B}
\definecolor{popblue}{HTML}{2F7DE1}
\definecolor{popgold}{HTML}{C98A12}
\definecolor{popmuted}{HTML}{8A7F76}
\definecolor{popline}{HTML}{EADFD2}
\definecolor{popgray}{HTML}{8A7F76}
\colorlet{popgrey}{popgray}
% Alias tolérants (le modèle nomme parfois les couleurs différemment)
\colorlet{popwhite}{white}
\colorlet{popblack}{popink}
\colorlet{popred}{poppink}
\colorlet{popyellow}{popgold}
% Teintes claires pour fonds de tableaux / zones
\colorlet{popblueL}{popblue!10!white}
\colorlet{poporangeL}{poporange!14!white}
\colorlet{popgreenL}{popgreen!12!white}
\colorlet{poppurpleL}{poppurple!12!white}
\colorlet{poppinkL}{poppink!12!white}
\colorlet{popgoldL}{popgold!14!white}

\pagecolor{popcream}

% ============================================================
% Typographie
% ============================================================
\defaultfontfeatures{Ligatures=TeX}
\setmainfont{PlusJakartaSans-Medium.ttf}[Path=../../../fonts/,BoldFont=PlusJakartaSans-Bold.ttf]
% Maths en sans-serif (Fira Math) avec chiffres et lettres droites de Plus
% Jakarta, pour que les formules se fondent dans le texte.
\usepackage[math-style=ISO,bold-style=ISO]{unicode-math}
\setmathfont{FiraMath-Regular.otf}[Path=../../../fonts/]
\setmathfont{PlusJakartaSans-Medium.ttf}[Path=../../../fonts/,range={up/num,up/{Latin,latin},\mathup}]
\usepackage{icomma} % virgule décimale sans espace en mode math
% Caractères mathématiques Unicode absents des polices de texte (Plus Jakarta,
% Archivo Black) : rendus par leur équivalent mathématique, en texte comme en
% formule — sinon ils s'affichent en carré vide (ex. « Arithmétique dans ℤ »).
\usepackage{newunicodechar}
\newunicodechar{ℝ}{\ensuremath{\mathbb{R}}}
\newunicodechar{ℤ}{\ensuremath{\mathbb{Z}}}
\newunicodechar{ℂ}{\ensuremath{\mathbb{C}}}
\newunicodechar{ℕ}{\ensuremath{\mathbb{N}}}
\newunicodechar{ℚ}{\ensuremath{\mathbb{Q}}}
\newunicodechar{𝔻}{\ensuremath{\mathbb{D}}}
\newunicodechar{α}{\ensuremath{\alpha}}
\newunicodechar{β}{\ensuremath{\beta}}
\newunicodechar{γ}{\ensuremath{\gamma}}
\newunicodechar{δ}{\ensuremath{\delta}}
\newunicodechar{ε}{\ensuremath{\varepsilon}}
\newunicodechar{θ}{\ensuremath{\theta}}
\newunicodechar{λ}{\ensuremath{\lambda}}
\newunicodechar{μ}{\ensuremath{\mu}}
\newunicodechar{σ}{\ensuremath{\sigma}}
\newunicodechar{τ}{\ensuremath{\tau}}
\newunicodechar{φ}{\ensuremath{\varphi}}
\newunicodechar{ψ}{\ensuremath{\psi}}
\newunicodechar{ω}{\ensuremath{\omega}}
\newunicodechar{Σ}{\ensuremath{\Sigma}}
% majuscules produites par \MakeUppercase dans les titres (α-aminés -> Α)
\newunicodechar{Α}{\ensuremath{\alpha}}
\newunicodechar{Β}{\ensuremath{\beta}}
\newunicodechar{Γ}{\ensuremath{\Gamma}}
\newunicodechar{Δ}{\ensuremath{\Delta}}
\newunicodechar{Ε}{\ensuremath{\varepsilon}}
\newunicodechar{Θ}{\ensuremath{\Theta}}
\newunicodechar{Λ}{\ensuremath{\Lambda}}
\newunicodechar{Μ}{\ensuremath{\mu}}
\newunicodechar{Τ}{\ensuremath{\tau}}
\newunicodechar{Φ}{\ensuremath{\Phi}}
\newunicodechar{Ψ}{\ensuremath{\Psi}}
\newunicodechar{Ω}{\ensuremath{\Omega}}
\newunicodechar{↦}{\ensuremath{\mapsto}}
\newunicodechar{∈}{\ensuremath{\in}}
\newunicodechar{∉}{\ensuremath{\notin}}
\newunicodechar{∧}{\ensuremath{\wedge}}
\newunicodechar{⇔}{\ensuremath{\Leftrightarrow}}
\newunicodechar{⟺}{\ensuremath{\Longleftrightarrow}}
\newunicodechar{⇒}{\ensuremath{\Rightarrow}}
\newunicodechar{⟹}{\ensuremath{\Longrightarrow}}
\newunicodechar{∘}{\ensuremath{\circ}}
\newunicodechar{∪}{\ensuremath{\cup}}
\newunicodechar{∩}{\ensuremath{\cap}}
\newunicodechar{≡}{\ensuremath{\equiv}}
\newunicodechar{⊂}{\ensuremath{\subset}}
\newunicodechar{⊥}{\ensuremath{\perp}}
\newunicodechar{∃}{\ensuremath{\exists}}
\newunicodechar{∀}{\ensuremath{\forall}}
\newunicodechar{∛}{\ensuremath{\sqrt[3]{\,}}}
\newunicodechar{∜}{\ensuremath{\sqrt[4]{\,}}}
\newunicodechar{⃗}{\ensuremath{{}^{\to}}}
\newunicodechar{ⁿ}{\ifmmode^{n}\else\textsuperscript{\ensuremath{n}}\fi}
\newunicodechar{ˣ}{\ifmmode^{x}\else\textsuperscript{\ensuremath{x}}\fi}
\newunicodechar{ᵇ}{\ifmmode^{b}\else\textsuperscript{\ensuremath{b}}\fi}
\newunicodechar{ʳ}{\ifmmode^{r}\else\textsuperscript{\ensuremath{r}}\fi}
\newunicodechar{ᵘ}{\ifmmode^{u}\else\textsuperscript{\ensuremath{u}}\fi}
\newunicodechar{ʸ}{\ifmmode^{y}\else\textsuperscript{\ensuremath{y}}\fi}
\newunicodechar{ᵃ}{\ifmmode^{a}\else\textsuperscript{\ensuremath{a}}\fi}
\newunicodechar{ᵉ}{\ifmmode^{e}\else\textsuperscript{\ensuremath{e}}\fi}
\newunicodechar{ᵏ}{\ifmmode^{k}\else\textsuperscript{\ensuremath{k}}\fi}
\newunicodechar{ᵖ}{\ifmmode^{p}\else\textsuperscript{\ensuremath{p}}\fi}
\newunicodechar{ᴺ}{\ifmmode^{N}\else\textsuperscript{\ensuremath{N}}\fi}
\newunicodechar{ᶜ}{\ifmmode^{c}\else\textsuperscript{\ensuremath{c}}\fi}
\newunicodechar{ʰ}{\ifmmode^{h}\else\textsuperscript{\ensuremath{h}}\fi}
\newunicodechar{ᵠ}{\ifmmode^{q}\else\textsuperscript{\ensuremath{q}}\fi}
\newunicodechar{ₙ}{\ifmmode_{n}\else\textsubscript{\ensuremath{n}}\fi}
\newunicodechar{ₐ}{\ifmmode_{a}\else\textsubscript{\ensuremath{a}}\fi}
\newunicodechar{ᵢ}{\ifmmode_{i}\else\textsubscript{\ensuremath{i}}\fi}
\newunicodechar{ᵣ}{\ifmmode_{r}\else\textsubscript{\ensuremath{r}}\fi}
\newunicodechar{ₖ}{\ifmmode_{k}\else\textsubscript{\ensuremath{k}}\fi}
\newunicodechar{ᵦ}{\ifmmode_{\beta}\else\textsubscript{\ensuremath{\beta}}\fi}
\newunicodechar{ₓ}{\ifmmode_{x}\else\textsubscript{\ensuremath{x}}\fi}
\newfontfamily\popdisplay{ArchivoBlack-Regular.ttf}[Path=../../../fonts/]
\newfontfamily\poplogo{SpaceGrotesk-Bold.ttf}[Path=../../../fonts/]
\newfontfamily\popsemi{PlusJakartaSans-SemiBold.ttf}[Path=../../../fonts/]
\newfontfamily\popxbold{PlusJakartaSans-ExtraBold.ttf}[Path=../../../fonts/]
\newfontfamily\popxboldls{PlusJakartaSans-ExtraBold.ttf}[Path=../../../fonts/,LetterSpace=3.0]

\setlist[enumerate]{leftmargin=1.7em,itemsep=5pt,topsep=4pt}
\setlist[itemize]{leftmargin=1.7em,itemsep=4pt,topsep=4pt,label={\color{poporange}\textbullet}}
\setlength{\parindent}{0pt}
\setlength{\parskip}{5pt}
\renewcommand{\arraystretch}{1.25}

% pgfplots : style Pop par défaut
\pgfplotsset{
  every axis/.append style={axis line style={popink,line width=1pt},tick style={popink},
    grid style={popline,line width=0.5pt},label style={font=\small},tick label style={font=\footnotesize}},
  popcurve/.style={poporange,line width=1.8pt,no markers,smooth},
}
% Même style disponible dans un \draw TikZ simple (hors pgfplots)
\tikzset{popcurve/.style={poporange,line width=1.8pt}}

% ============================================================
% Pill / squiggle
% ============================================================
\newcommand{\poppill}[3]{%
  \tikz[baseline=-0.55ex]\node[draw=popink,line width=1.2pt,rounded corners=2.6pt,%
    fill=#2,inner xsep=6.5pt,inner ysep=3pt]{%
    \popxboldls\fontsize{7}{7}\selectfont\color{#3}\MakeUppercase{#1}};%
}
\newcommand{\popsquiggle}[1]{%
  \tikz[baseline=-0.4ex]{
    \draw[#1,line width=2.6pt,line cap=round]
      plot [smooth] coordinates {(0,0.09) (0.34,0.01) (0.68,0.21) (1.02,0.01) (1.36,0.10)};
  }%
}
% Mot-clé surligné (marqueur orange sous le texte)
\newcommand{\cle}[1]{{\popxbold\color{popdark}#1}}

% ============================================================
% En-tête / pied
% ============================================================
\pagestyle{fancy}
\fancyhf{}
\renewcommand{\headrulewidth}{0pt}
\renewcommand{\footrulewidth}{0pt}
\lhead{\raisebox{-0.15ex}{\poplogo\fontsize{13}{13}\selectfont\color{popink}OnBuch\color{poporange}+}}
\rhead{\poppill{\DOCMATIERE\ \textbullet\ \DOCNIVEAU\ \textbullet\ Leçon \DOCLECON}{white}{popink}}
\lfoot{\popxbold\fontsize{7.6}{7.6}\selectfont\color{popmuted}\MakeUppercase{Cours signé OnBuch+}\ \textbullet\ {\popsemi\DOCTITRE}}
\rfoot{\tikz[baseline=-0.6ex]\node[rounded corners=8.5pt,fill=popink,minimum height=17pt,minimum width=17pt,inner xsep=5pt,inner sep=0]{\popxbold\fontsize{8}{8}\selectfont\color{popcream}\thepage\,/\,\pageref*{LastPage}};}

% ============================================================
% Titres
% ============================================================
\newcommand{\chaptitre}[2]{%
  \begin{tcolorbox}[enhanced,colback=poporange,colframe=popink,boxrule=2.6pt,arc=11pt,
      drop shadow=popink,left=15pt,right=15pt,top=12pt,bottom=11pt,before skip=3pt,after skip=18pt]
    \poppill{#2}{popink}{popcream}\par\vspace{9pt}
    {\raggedright\hyphenpenalty=10000\exhyphenpenalty=10000\popdisplay\fontsize{22}{24}\selectfont\color{popcream}\MakeUppercase{#1}\par}
  \end{tcolorbox}
}
% Section numérotée : pastille orange avec le numéro + titre Archivo
\newcounter{popsec}
\newcommand{\coursec}[1]{%
  \stepcounter{popsec}\setcounter{popsub}{0}%
  \par\vspace{16pt}\noindent
  \begin{minipage}{\linewidth}
  \tikz[baseline=-0.7ex]\node[draw=popink,line width=1.6pt,fill=poporange,rounded corners=4pt,minimum size=22pt,inner sep=0]{\popdisplay\fontsize{12}{12}\selectfont\color{popcream}\thepopsec};%
  \hspace{8pt}{\popdisplay\fontsize{15}{17}\selectfont\color{popink}\MakeUppercase{#1}}\par
  \vspace{4pt}\noindent{\color{poporange}\rule{36pt}{3.6pt}}
  \end{minipage}\par\nopagebreak\vspace{8pt}
}
\newcounter{popsub}
\newcommand{\courssub}[1]{%
  \stepcounter{popsub}%
  \par\vspace{9pt}\noindent{\popxbold\fontsize{11.5}{13}\selectfont\color{popink}{\color{poporange}\thepopsec.\thepopsub}\hspace{6pt}#1}\par\nopagebreak\vspace{4pt}
}

% ============================================================
% Boîtes pédagogiques — même recette : fond blanc, bordure épaisse colorée,
% ombre dure encre, badge plein. Titre optionnel [#1].
% ============================================================
\tcbset{popbase/.style={breakable,enhanced,colback=white,boxrule=2.3pt,arc=9pt,
  shadow={3pt}{-3pt}{0pt}{popink},left=14pt,right=14pt,top=11pt,bottom=11pt,before skip=11pt,after skip=15pt,
  fontupper=\popsemi\fontsize{10.6}{14.5}\selectfont\color{popink},
  fontlower=\popsemi\fontsize{10.6}{14.5}\selectfont\color{popink}}}
% Coche dessinée (le glyphe ✓ n'existe pas dans les polices du document)
\newcommand{\popcheck}[1]{\tikz[baseline=-0.5ex]\draw[#1,line width=1.6pt,line cap=round,line join=round](-0.13,0.0)--(-0.04,-0.09)--(0.13,0.1);}
\newcommand{\popboxhead}[3]{\poppill{#1}{#2}{white}\ifx\relax#3\relax\else\hspace{6pt}{\popxbold\fontsize{10.6}{12}\selectfont\color{popink}#3}\fi\par\vspace{7pt}}

\newtcolorbox{aretenir}[1][]{popbase,colframe=popgreen,before upper={\popboxhead{À retenir}{popgreen}{#1}}}
\newtcolorbox{exemplebox}[1][]{popbase,colframe=poporange,before upper={\popboxhead{Exemple}{poporange}{#1}}}
\newtcolorbox{definition}[1][]{popbase,colframe=popblue,before upper={\popboxhead{Définition}{popblue}{#1}}}
\newtcolorbox{propriete}[1][]{popbase,colframe=poppurple,before upper={\popboxhead{Propriété}{poppurple}{#1}}}
\newtcolorbox{methode}[1][]{popbase,colframe=popgold,before upper={\popboxhead{Méthode}{popgold}{#1}}}
\newtcolorbox{attention}[1][]{popbase,colframe=poppink,before upper={\popboxhead{Attention}{poppink}{#1}}}
\newtcolorbox{experience}[1][]{popbase,colframe=popblue,colback=popblueL,before upper={\popboxhead{Expérience}{popblue}{#1}}}
% « Le savais-tu ? » : carte violette pleine (contexte, culture, Cameroun)
\newtcolorbox{savaistu}[1][]{popbase,colback=poppurple,colframe=popink,
  fontupper=\popsemi\fontsize{10.4}{14.2}\selectfont\color{white},
  before upper={\poppill{Le savais-tu\,?}{white}{poppurple}\ifx\relax#1\relax\else\hspace{6pt}{\popxbold\color{white}#1}\fi\par\vspace{7pt}}}
% Exercice résolu : énoncé en haut, \tcblower puis la solution sur fond crème
\newcounter{popexo}\newcounter{popexr}
\newtcolorbox{exoresolu}[1][]{popbase,colframe=poporange,colbacklower=popcream,
  segmentation style={popink,line width=1.2pt,dashed},
  before upper={\stepcounter{popexr}\popboxhead{Exercice résolu \thepopexr}{poporange}{#1}},
  before lower={\poppill{Solution}{popink}{popcream}\par\vspace{6pt}}}
% Exercice (non résolu en place) — la correction est dans la partie Corrigés
\newtcolorbox{exercice}[1][]{popbase,colframe=popink,
  before upper={\stepcounter{popexo}\popboxhead{Exercice \thepopexo}{popink}{#1}}}
% Corrigé
\newtcolorbox{corrige}[1][]{popbase,colframe=popgreen,colback=popgreenL,
  before upper={\popboxhead{Corrigé}{popgreen}{#1}}}
% Objectifs / prérequis / bilan
\newtcolorbox{objectifs}{popbase,colframe=popink,colback=poporangeL,
  before upper={\popboxhead{Objectifs de la leçon}{popink}{}}}
\newtcolorbox{prerequis}{popbase,colframe=popink,colback=popblueL,
  before upper={\popboxhead{Prérequis}{popblue}{}}}
\newtcolorbox{bilan}{popbase,colframe=popink,colback=white,boxrule=2.8pt,
  before upper={\popboxhead{Fiche bilan}{poporange}{L'essentiel à connaître pour l'examen}}}
% Figure encadrée avec légende
\newtcolorbox{popfigure}[1][]{enhanced,breakable=false,colback=white,colframe=popline,boxrule=1.4pt,arc=8pt,
  left=10pt,right=10pt,top=10pt,bottom=8pt,before skip=10pt,after skip=12pt,halign=center,
  lower separated=false,halign lower=center,
  fontlower=\popsemi\fontsize{9}{11.5}\selectfont\color{popmuted},
  #1}
\newcommand{\legende}[1]{\par\vspace{6pt}{\popsemi\fontsize{9}{11.5}\selectfont\color{popmuted}#1}}

% ============================================================
% Signature OnBuch+
% ============================================================
\newcommand{\poplogoinline}[1]{{\poplogo\fontsize{#1}{#1}\selectfont\color{popink}OnBuch\color{poporange}+}}
\newcommand{\signatureonbuch}{%
  \begin{tcolorbox}[enhanced,colback=popink,colframe=popink,boxrule=2.2pt,arc=10pt,
      drop shadow=poporange,left=14pt,right=14pt,top=10pt,bottom=10pt,before skip=0pt,after skip=0pt]
    \tikz[baseline=-0.7ex]\node[circle,fill=popgreen,draw=popcream,line width=1.3pt,minimum size=22pt,inner sep=0]{\popcheck{white}};%
    \hspace{9pt}\begin{minipage}[c]{0.62\linewidth}
      {\popxbold\fontsize{12}{13}\selectfont\color{popcream}Signé\ \textbullet\ {\poplogo OnBuch\color{poporange}+}}\par\vspace{2pt}
      {\raggedright\popsemi\fontsize{8}{10.5}\selectfont\color{popline}\MakeUppercase{Équipe pédagogique OnBuch+}\\\MakeUppercase{Conforme au programme officiel MINESEC}\par}
    \end{minipage}\hfill
    {\poplogo\fontsize{22}{22}\selectfont\color{popcream}OnBuch\color{poporange}+}
  \end{tcolorbox}%
}

% ============================================================
% Page de garde
% #1 titre · #2 matière · #3 niveau · #4 module · #5 n° leçon · #6 durée ·
% #7 description / objectifs (texte ou itemize)
% ============================================================
\newcommand{\pagegarde}[7]{%
  \thispagestyle{empty}
  \newgeometry{a4paper,margin=1.7cm,top=1.6cm,bottom=1.4cm}%
  \begin{tikzpicture}[remember picture,overlay]
    \fill[poporange!18] ([xshift=-3.2cm,yshift=-3.6cm]current page.north east) circle (4.6cm);
    \fill[poppurple!14] ([xshift=2.4cm,yshift=7.6cm]current page.south west) circle (3.8cm);
  \end{tikzpicture}%
  {\poplogo\fontsize{30}{30}\selectfont\color{popink}OnBuch\color{poporange}+}\hfill
  \raisebox{0.6ex}{\poppill{Cours signé \textbullet\ Premium}{popink}{popcream}}\par
  \vspace{1pt}\popsquiggle{poppurple}\par
  \vspace{8pt}
  \begin{tcolorbox}[enhanced,colback=poporange,colframe=popink,boxrule=2.8pt,arc=13pt,
      drop shadow=popink,left=18pt,right=18pt,top=12pt,bottom=13pt,after skip=10pt]
    \poppill{#2 \textbullet\ #3}{popink}{popcream}\hspace{5pt}\poppill{#4}{white}{popink}\par\vspace{8pt}
    {\popdisplay\fontsize{14}{14}\selectfont\color{popink}LEÇON #5}\par\vspace{4pt}
    {\raggedright\hyphenpenalty=10000\exhyphenpenalty=10000\popdisplay\fontsize{25}{27}\selectfont\color{popcream}\MakeUppercase{#1}\par}
    \vspace{8pt}
    {\popxbold\fontsize{9.5}{9.5}\selectfont\color{popink}\MakeUppercase{Durée conseillée : #6}}
  \end{tcolorbox}
  \begin{tcolorbox}[enhanced,colback=white,colframe=popink,boxrule=2.2pt,arc=10pt,
      drop shadow=popink,left=16pt,right=16pt,top=10pt,bottom=10pt,after skip=8pt]
    \poppill{Dans cette leçon}{poppurple}{white}\par\vspace{5pt}
    \popsemi\fontsize{9.3}{12.6}\selectfont\color{popink}#7
  \end{tcolorbox}
  \noindent\begin{minipage}[t]{0.487\linewidth}
    \begin{tcolorbox}[enhanced,colback=popgreen,colframe=popink,boxrule=2.2pt,arc=10pt,
        drop shadow=popink,left=11pt,right=11pt,top=10pt,bottom=10pt,height=3.1cm,valign=center]
      \tcbox[colback=white,colframe=popink,boxrule=1.3pt,arc=5pt,boxsep=0pt,nobeforeafter,
        left=3pt,right=3pt,top=3pt,bottom=3pt,valign=center]{\qrcode[height=1.9cm]{https://play.google.com/store/apps/details?id=com.luvvix.onbuch&hl=fr}}%
      \hspace{8pt}\begin{minipage}[c]{0.5\linewidth}
        {\popdisplay\fontsize{11}{12.5}\selectfont\color{white}\MakeUppercase{Télécharge OnBuch}}\par\vspace{4pt}
        {\raggedright\popsemi\fontsize{8.4}{11}\selectfont\color{white}Gratuit \textbullet\ Google Play\\Tuteur IA, annales, concours, résultats\par}
      \end{minipage}
    \end{tcolorbox}
  \end{minipage}\hfill
  \begin{minipage}[t]{0.487\linewidth}
    \begin{tcolorbox}[enhanced,colback=poppurple,colframe=popink,boxrule=2.2pt,arc=10pt,
        drop shadow=popink,left=11pt,right=11pt,top=10pt,bottom=10pt,height=3.1cm,valign=center]
      \tikz[baseline=-0.7ex]\node[circle,fill=white,draw=popink,line width=1.3pt,minimum size=1.6cm,inner sep=0]{\popdisplay\fontsize{24}{24}\selectfont\color{poppurple}+};%
      \hspace{8pt}\begin{minipage}[c]{0.6\linewidth}
        {\popdisplay\fontsize{11}{12.5}\selectfont\color{white}\MakeUppercase{Passe à OnBuch+}}\par\vspace{4pt}
        {\raggedright\popsemi\fontsize{8.4}{11}\selectfont\color{white}Tous les cours signés, Tuteur IA illimité, fiches PDF — onglet OnBuch+ de l'app\par}
      \end{minipage}
    \end{tcolorbox}
  \end{minipage}
  \vspace{8pt}
  \popcontenu
  \vfill
  \signatureonbuch
  \restoregeometry
  \newpage
}

% Rangée « Ce cours contient » (page de garde)
\newcommand{\poptile}[3]{% #1 couleur #2 gros texte #3 légende
  \begin{tcolorbox}[enhanced,colback=#1,colframe=popink,boxrule=2pt,arc=9pt,shadow={2.5pt}{-2.5pt}{0pt}{popink},
      left=6pt,right=6pt,top=9pt,bottom=9pt,halign=center,valign=center,height=1.75cm,width=\linewidth,nobeforeafter]
    {\popdisplay\fontsize{12.5}{13}\selectfont\color{white}\MakeUppercase{#2}}\par\vspace{3pt}
    {\centering\popsemi\fontsize{7.8}{9.5}\selectfont\color{white}#3\par}
  \end{tcolorbox}}
\newcommand{\popcontenu}{%
  \par\noindent{\popxboldls\fontsize{7.5}{7.5}\selectfont\color{popmuted}CE COURS SIGNÉ CONTIENT}\par\vspace{7pt}
  \noindent\begin{minipage}[t]{0.235\linewidth}\poptile{popblue}{Cours}{complet, illustré, pas à pas}\end{minipage}\hfill
  \begin{minipage}[t]{0.235\linewidth}\poptile{poporange}{Méthodes}{et exercices résolus}\end{minipage}\hfill
  \begin{minipage}[t]{0.235\linewidth}\poptile{poppink}{Exercices}{type examen + corrigés}\end{minipage}\hfill
  \begin{minipage}[t]{0.235\linewidth}\poptile{popgreen}{Fiche bilan}{l'essentiel à retenir}\end{minipage}\par}

% Sommaire
\newtcolorbox{sommaire}{enhanced,colback=white,colframe=popink,boxrule=2.2pt,arc=10pt,
  drop shadow=popink,left=18pt,right=18pt,top=13pt,bottom=13pt,before skip=6pt,after skip=20pt,
  before upper={\poppill{Sommaire}{poppurple}{white}\par\vspace{9pt}\popsemi\fontsize{10.6}{14.5}\selectfont\color{popink}}}

% Signature de fin de cours — poussée tout en bas de la dernière page
\newcommand{\findecours}{%
  \par\vfill\nopagebreak
  \begin{center}\popsquiggle{poporange}\end{center}\vspace{4pt}
  \signatureonbuch
}
\AtBeginDocument{\def\llbracket{\mathopen{[\mkern-3.5mu[}}\def\rrbracket{\mathclose{]\mkern-3.5mu]}}} % intervalles d'entiers (glyphe ⟦ absent de la police)
% Glyphes absents de Fira Math : redéfinis à partir de glyphes disponibles
\AtBeginDocument{%
  \def\setminus{\mathbin{\backslash}}\let\smallsetminus\setminus
  \def\vdots{\mathord{\vbox{\baselineskip3.2pt\lineskiplimit0pt\kern4pt\hbox{.}\hbox{.}\hbox{.}}}}%
  \def\ddots{\mathinner{\mkern1mu\raise6.5pt\hbox{.}\mkern2mu\raise3.6pt\hbox{.}\mkern2mu\raise0.7pt\hbox{.}\mkern1mu}}%
  \def\triangle{\mathord{\scalebox{0.9}{$\symup{\Delta}$}}}%
  \def\nabla{\mathord{\raisebox{\depth}{\scalebox{1}[-1]{$\symup{\Delta}$}}}}%
  \def\checkmark{\popcheck{popdark}}%
  \def\star{\mathbin{\ast}}\def\bigstar{\mathord{\ast}}%
  \def\diamond{\mathbin{\scalebox{0.6}{$\Diamond$}}}%
}

%% FIGFIX-BEGIN v5 — correctifs globaux des figures (docs/onbuch-generation/tools/figfix)
\usepackage{adjustbox}\usepackage{etoolbox}\tcbuselibrary{raster}
% toute figure TikZ/pgfplots plus large que la ligne est réduite à la largeur disponible
\BeforeBeginEnvironment{tikzpicture}{\begin{adjustbox}{max width=\linewidth}}
\AfterEndEnvironment{tikzpicture}{\end{adjustbox}}
% légendes pgfplots lisibles par-dessus les courbes
\pgfplotsset{every axis/.append style={legend style={fill=white,fill opacity=0.92,text opacity=1,draw=black!15,inner sep=3pt}}}
% étiquettes d'arêtes (syntaxe edge["..."]) avec fond blanc
\tikzset{every edge quotes/.append style={fill=white,inner sep=1.2pt}}
%% FIGFIX-END
\begin{document}

\pagegarde{Calcul intégral}{Mathématiques}{Tle D}{Relations et opérations fondamentales dans l'ensemble des nombres réels}{08}{14 h}{Cette leçon introduit l'intégrale d'une fonction continue sur un segment comme aire algébrique sous la courbe, en établit les propriétés fondamentales (linéarité, positivité, relation de Chasles) et développe les techniques de calcul (primitives, intégration par parties, changement de variable affine). Elle se conclut par les applications géométriques majeures au Baccalauréat : calculs d'aires, de volumes de révolution et approximation par la méthode des rectangles.
\vspace{4pt}
\begin{multicols}{2}\small
\begin{itemize}
\item À la fin de cette leçon, je sais définir et noter l'intégrale d'une fonction continue sur un segment [a ; b] et l'interpréter comme une aire algébrique.
\item À la fin de cette leçon, je sais appliquer la linéarité, la positivité et la relation de Chasles pour transformer ou comparer des intégrales.
\item À la fin de cette leçon, je sais étudier le sens de variation de la fonction x ↦ ∫ₐˣ f(t) dt sans utiliser la notion de primitive.
\item À la fin de cette leçon, je sais utiliser l'inégalité de la moyenne pour encadrer une intégrale et calculer la valeur moyenne d'une fonction.
\item À la fin de cette leçon, je sais calculer une intégrale par lecture directe d'une primitive, par intégration par parties ou par changement de variable affine.
\item À la fin de cette leçon, je sais calculer l'intégrale d'une fonction valeur absolue en découpant l'intervalle.
\end{itemize}\end{multicols}}

\begin{sommaire}
\begin{multicols}{2}
\begin{enumerate}
\item Intégrale d'une fonction continue : définition et premières propriétés
\item Propriétés algébriques : linéarité, positivité, relation de Chasles
\item Fonction définie par une intégrale et valeur moyenne
\item Techniques de calcul : primitives, intégration par parties, changement de variable affine
\item Applications géométriques : aires et volumes de révolution
\item Approximation d'une intégrale : méthode des rectangles
\item Activité d'intégration
\item Méthodes et exercices résolus
\item Exercices
\item Corrigés des exercices
\item Fiche bilan
\end{enumerate}
\end{multicols}
\end{sommaire}

\begin{objectifs}
\begin{itemize}
\item À la fin de cette leçon, je sais définir et noter l'intégrale d'une fonction continue sur un segment [a ; b] et l'interpréter comme une aire algébrique.
\item À la fin de cette leçon, je sais appliquer la linéarité, la positivité et la relation de Chasles pour transformer ou comparer des intégrales.
\item À la fin de cette leçon, je sais étudier le sens de variation de la fonction x ↦ ∫ₐˣ f(t) dt sans utiliser la notion de primitive.
\item À la fin de cette leçon, je sais utiliser l'inégalité de la moyenne pour encadrer une intégrale et calculer la valeur moyenne d'une fonction.
\item À la fin de cette leçon, je sais calculer une intégrale par lecture directe d'une primitive, par intégration par parties ou par changement de variable affine.
\item À la fin de cette leçon, je sais calculer l'intégrale d'une fonction valeur absolue en découpant l'intervalle.
\item À la fin de cette leçon, je sais calculer l'aire entre une courbe et l'axe des abscisses ou entre deux courbes, en unités d'aire puis en unités usuelles.
\item À la fin de cette leçon, je sais calculer le volume d'un solide de révolution engendré par la rotation d'une courbe autour de l'axe des abscisses.
\item À la fin de cette leçon, je sais encadrer une intégrale par la méthode des rectangles et contrôler la précision de l'encadrement.
\end{itemize}
\end{objectifs}

\begin{prerequis}
\begin{itemize}
\item Dérivation : dérivées des fonctions usuelles, dérivée d'une composée, tableau de variation (Première et début de Terminale)
\item Fonctions usuelles : polynômes, racine carrée, exponentielle, logarithme népérien, fonctions trigonométriques
\item Lecture graphique : signe d'une fonction, position relative de deux courbes, repère orthogonal et unités d'aire
\item Notion d'aire des figures usuelles (rectangle, triangle, trapèze, disque) et volumes usuels (cylindre, cône, sphère)
\item Sommes finies et manipulation d'inégalités sur les fonctions bornées
\end{itemize}
\end{prerequis}

\newpage

\coursec{Intégrale d'une fonction continue : définition et premières propriétés}

À Bafoussam, une coopérative agricole remplit chaque jour un grand réservoir d'eau destiné à l'irrigation des champs de maïs et de pommes de terre. Le débit de la pompe n'est pas constant : faible au lever du jour, il augmente en milieu de matinée, puis diminue en fin d'après-midi. Le responsable connaît la loi qui donne ce débit (en litres par heure) en fonction de l'heure $t$. \textbf{Question : quel volume total d'eau s'est accumulé dans le réservoir entre 6 h et 18 h ?}

Si le débit était constant, la réponse serait immédiate : volume = débit $\times$ durée, c'est-à-dire l'aire d'un rectangle. Mais ici le débit \emph{varie} : la quantité d'eau ajoutée pendant une petite durée $\Delta t$ autour de l'instant $t$ vaut environ $d(t)\times \Delta t$, et le volume total est la « somme » de toutes ces petites contributions. Additionner une infinité de quantités infiniment petites : voilà exactement l'idée du \cle{calcul intégral}.

Le même besoin apparaît partout : un ingénieur des Ponts et Chaussées qui calcule la surface d'un terrain délimité par une route rectiligne et une rivière sinueuse, un industriel qui détermine le volume d'un silo à maïs obtenu par rotation d'une courbe, un physicien qui déduit la distance parcourue d'un véhicule à partir de sa vitesse instantanée.

\textbf{Problématique de la leçon :} comment donner un sens précis à la notion d'« aire sous une courbe », comment la calculer, et quelles propriétés fondamentales possède-t-elle ?

\courssub{Approche intuitive : l'aire sous la courbe}

\textbf{Le problème de l'aire}\par

Depuis le collège, tu sais calculer l'aire d'un rectangle, d'un triangle, d'un disque. Toutes ces figures ont des bords « simples » : segments de droite ou cercle. Mais quelle est l'aire du domaine situé \emph{sous} la courbe d'une fonction, par exemple sous la parabole d'équation $y = x^2$ entre $x = 0$ et $x = 1$ ? Les formules du collège ne suffisent plus : le bord supérieur est \emph{courbe}.

L'idée, déjà présente chez Archimède (III\textsuperscript{e} siècle av. J.-C.), est d'une simplicité géniale : \textbf{approcher le domaine courbe par des rectangles}, dont on sait parfaitement calculer les aires, puis faire tendre la largeur des rectangles vers zéro.

\begin{experience}[Découpage en rectangles : l'idée d'Archimède]
Considérons la fonction $f(x) = x^2$ sur $[0\,;\,1]$ et cherchons l'aire $\mathcal{A}$ du domaine sous sa courbe.

\textbf{Étape 1 :} partageons $[0\,;\,1]$ en $4$ intervalles de même largeur $\dfrac{1}{4}$. Sur chaque intervalle, construisons le rectangle dont la hauteur est la valeur de $f$ à l'extrémité droite. La somme des aires de ces rectangles vaut :
\[ S_4 = \frac{1}{4}\left[\left(\frac{1}{4}\right)^2 + \left(\frac{2}{4}\right)^2 + \left(\frac{3}{4}\right)^2 + \left(\frac{4}{4}\right)^2\right] = \frac{1}{4}\times\frac{1+4+9+16}{16} = \frac{30}{64} = \num{0,46875}. \]

\textbf{Étape 2 :} avec $8$ rectangles de largeur $\dfrac{1}{8}$, on obtient $S_8 \approx \num{0,3984}$. Avec $16$ rectangles, $S_{16} \approx \num{0,3652}$.

\textbf{Étape 3 :} les valeurs $S_4, S_8, S_{16}, \ldots$ décroissent et se rapprochent d'un nombre limite. On démontrera plus tard que cette limite vaut exactement $\dfrac{1}{3}$ : c'est l'aire cherchée. Plus le découpage est fin, plus les rectangles « épousent » la courbe, et plus la somme de leurs aires approche l'aire exacte.
\end{experience}

\begin{popfigure}
\begin{center}
\begin{tikzpicture}
\begin{axis}[width=0.32\linewidth, axis lines=middle, xlabel={$x$}, ylabel={$y$}, xmin=-0.1, xmax=1.25, ymin=-0.1, ymax=1.3, xtick={1}, ytick=\empty, title={$4$ rectangles}, title style={font=\small}, every axis/.append style={font=\small}]
\addplot[popcurve, domain=0:1.15, samples=100]{x^2};
\addplot[fill=popblueL, draw=popblue] coordinates {(0,0) (0,0.0625) (0.25,0.0625) (0.25,0)} \closedcycle;
\addplot[fill=popblueL, draw=popblue] coordinates {(0.25,0) (0.25,0.25) (0.5,0.25) (0.5,0)} \closedcycle;
\addplot[fill=popblueL, draw=popblue] coordinates {(0.5,0) (0.5,0.5625) (0.75,0.5625) (0.75,0)} \closedcycle;
\addplot[fill=popblueL, draw=popblue] coordinates {(0.75,0) (0.75,1) (1,1) (1,0)} \closedcycle;
\addplot[popcurve, domain=0:1.15, samples=100]{x^2};
\end{axis}
\end{tikzpicture}\hfill
\begin{tikzpicture}
\begin{axis}[width=0.32\linewidth, axis lines=middle, xlabel={$x$}, ylabel={$y$}, xmin=-0.1, xmax=1.25, ymin=-0.1, ymax=1.3, xtick={1}, ytick=\empty, title={$8$ rectangles}, title style={font=\small}, every axis/.append style={font=\small}]
\foreach \i in {0,...,7}{
  \addplot[fill=poporangeL, draw=poporange] coordinates {(\i*0.125,0) (\i*0.125,{(\i*0.125+0.125)^2}) ({\i*0.125+0.125},{(\i*0.125+0.125)^2}) ({\i*0.125+0.125},0)} \closedcycle;
}
\addplot[popcurve, domain=0:1.15, samples=100]{x^2};
\end{axis}
\end{tikzpicture}\hfill
\begin{tikzpicture}
\begin{axis}[width=0.32\linewidth, axis lines=middle, xlabel={$x$}, ylabel={$y$}, xmin=-0.1, xmax=1.25, ymin=-0.1, ymax=1.3, xtick={1}, ytick=\empty, title={$16$ rectangles}, title style={font=\small}, every axis/.append style={font=\small}]
\foreach \i in {0,...,15}{
  \addplot[fill=popgreenL, draw=popgreen] coordinates {(\i*0.0625,0) (\i*0.0625,{(\i*0.0625+0.0625)^2}) ({\i*0.0625+0.0625},{(\i*0.0625+0.0625)^2}) ({\i*0.0625+0.0625},0)} \closedcycle;
}
\addplot[popcurve, domain=0:1.15, samples=100]{x^2};
\end{axis}
\end{tikzpicture}
\end{center}
\legende{Approximation de l'aire sous la parabole $y = x^2$ sur $[0\,;\,1]$ par $4$, $8$ puis $16$ rectangles : la somme des aires converge vers $\frac{1}{3}$.}
\end{popfigure}

\begin{attention}[Ce que dit — et ne dit pas — cette approche]
Le découpage en rectangles est une \emph{approche intuitive} qui justifie la définition. Il ne constitue pas la méthode de calcul pratique : personne ne calcule des sommes de $10\,000$ rectangles à la main ! Nous verrons dans la suite de la leçon des outils bien plus efficaces (primitives, intégration par parties). Mais c'est cette image de « somme infinie de petits rectangles » qu'il faut garder en tête pour comprendre le sens de toutes les formules.
\end{attention}

\courssub{Définition de l'intégrale d'une fonction continue positive}

\begin{definition}[Intégrale d'une fonction continue positive]
Soit $f$ une fonction \cle{continue} et \cle{positive} sur un intervalle $[a\,;\,b]$ (avec $a < b$), de courbe représentative $\mathcal{C}_f$ dans un repère orthogonal $(O\,;\,\vec{\imath},\vec{\jmath})$.

On appelle \cle{intégrale de $f$ entre $a$ et $b$}, et on note $\displaystyle\int_a^b f(x)\,\mathrm{d}x$, l'\textbf{aire}, exprimée en \cle{unités d'aire}, du domaine
\[ \mathcal{D} = \left\{ M(x\,;\,y) \;:\; a \le x \le b \ \text{et}\ 0 \le y \le f(x) \right\}, \]
c'est-à-dire le domaine du plan délimité par la courbe $\mathcal{C}_f$, l'axe des abscisses et les droites verticales d'équations $x = a$ et $x = b$.
\end{definition}

\begin{popfigure}
\begin{center}
\begin{tikzpicture}
\begin{axis}[width=0.72\linewidth, axis lines=middle, xlabel={$x$}, ylabel={$y$}, xmin=-0.3, xmax=5.4, ymin=-0.3, ymax=3.4, xtick={1,4}, xticklabels={$a$,$b$}, ytick=\empty, every axis/.append style={font=\small}]
\addplot[fill=popblueL, draw=none, domain=1:4, samples=100]{0.35*(x-1)*(4-x)+1.2} \closedcycle;
\addplot[popcurve, domain=0.4:4.6, samples=150]{0.35*(x-1)*(4-x)+1.2} node[pos=0.85, above left, popdark]{$\mathcal{C}_f$};
\draw[popline, dashed] (axis cs:1,0) -- (axis cs:1,1.2);
\draw[popline, dashed] (axis cs:4,0) -- (axis cs:4,1.2);
\node[popblue] at (axis cs:2.5,1.6) {$\displaystyle\int_a^b f(x)\,\mathrm{d}x$};
\end{axis}
\end{tikzpicture}
\end{center}
\legende{L'intégrale $\int_a^b f(x)\,\mathrm{d}x$ est l'aire (en u.a.) du domaine hachuré situé sous la courbe d'une fonction positive.}
\end{popfigure}

\courssub{Lecture de la notation $\displaystyle\int_a^b f(x)\,\mathrm{d}x$}

Chaque élément de cette notation a un sens précis, hérité du découpage en rectangles :

\begin{itemize}
\item Le symbole $\displaystyle\int$ est un « S » allongé, initiale du mot latin \emph{summa} (somme) : il rappelle que l'intégrale est une \textbf{somme} infinie de petites contributions.
\item Les nombres $a$ et $b$ sont les \cle{bornes} de l'intégrale : $a$ est la borne inférieure, $b$ la borne supérieure. Elles indiquent \emph{entre quelles valeurs} on somme.
\item $f(x)$ est la hauteur du rectangle élémentaire au point d'abscisse $x$.
\item $\mathrm{d}x$ représente la \textbf{largeur infiniment petite} du rectangle élémentaire. Le produit $f(x)\,\mathrm{d}x$ est donc l'aire d'un rectangle « infiniment fin ».
\item La lettre $x$ est une \cle{variable muette} : elle ne joue aucun rôle dans le résultat. On peut la remplacer par n'importe quelle autre lettre :
\[ \int_a^b f(x)\,\mathrm{d}x = \int_a^b f(t)\,\mathrm{d}t = \int_a^b f(u)\,\mathrm{d}u. \]
\end{itemize}

\begin{aretenir}[L'unité d'aire]
Dans un repère orthogonal, l'\cle{unité d'aire} (u.a.) est l'aire du rectangle construit sur les vecteurs unitaires des axes. Si l'unité graphique sur $(Ox)$ est $a$ cm et celle sur $(Oy)$ est $b$ cm, alors :
\[ 1\ \text{u.a.} = a \times b\ \text{cm}^2. \]
Par exemple, avec \SI{2}{cm} par unité en abscisse et \SI{0,5}{cm} par unité en ordonnée, $1$ u.a. $= 2 \times \num{0,5} = \SI{1}{cm^2}$. Si l'intégrale vaut $6$ u.a., l'aire réelle du domaine est \SI{6}{cm^2}.
\end{aretenir}

\begin{exemplebox}[Premiers calculs par lecture géométrique]
\textbf{1. Intégrale d'une constante.} Calculons $\displaystyle\int_0^3 2\,\mathrm{d}x$. La fonction $x \mapsto 2$ est constante et positive : le domaine sous sa courbe entre $0$ et $3$ est un \textbf{rectangle} de longueur $3$ et de hauteur $2$. Donc :
\[ \int_0^3 2\,\mathrm{d}x = 3 \times 2 = 6\ \text{u.a.} \]

\textbf{2. Intégrale de la fonction identité.} Calculons $\displaystyle\int_0^4 x\,\mathrm{d}x$. La courbe de $x \mapsto x$ est la première bissectrice : le domaine sous la courbe entre $0$ et $4$ est un \textbf{triangle rectangle} de base $4$ et de hauteur $4$. Donc :
\[ \int_0^4 x\,\mathrm{d}x = \frac{4 \times 4}{2} = 8\ \text{u.a.} \]
Aucune formule sophistiquée : une simple lecture géométrique suffit quand le domaine est une figure usuelle.
\end{exemplebox}

\courssub{Extension aux fonctions de signe quelconque : l'aire algébrique}

Que se passe-t-il si la fonction prend des valeurs négatives ? L'aire géométrique, elle, est toujours positive. Mais pour que l'intégrale conserve de bonnes propriétés de calcul (additivité, linéarité), on adopte une convention de \textbf{signe}.

\begin{definition}[Intégrale d'une fonction continue de signe quelconque (admise)]
On admet que toute fonction \cle{continue} sur $[a\,;\,b]$ admet une intégrale entre $a$ et $b$, définie comme l'\cle{aire algébrique} du domaine compris entre la courbe et l'axe des abscisses :
\begin{itemize}
\item les portions de domaine situées \textbf{au-dessus} de l'axe $(Ox)$ sont comptées \textbf{positivement} ;
\item les portions situées \textbf{en dessous} de l'axe $(Ox)$ sont comptées \textbf{négativement}.
\end{itemize}
En particulier, si $f$ est négative sur $[a\,;\,b]$, alors $\displaystyle\int_a^b f(x)\,\mathrm{d}x = -\,\mathcal{A}$, où $\mathcal{A}$ est l'aire géométrique (positive) du domaine entre la courbe et l'axe.
\end{definition}

\begin{popfigure}
\begin{center}
\begin{tikzpicture}
\begin{axis}[width=0.75\linewidth, axis lines=middle, xlabel={$x$}, ylabel={$y$}, xmin=-0.4, xmax=6.6, ymin=-2.4, ymax=2.9, xtick={1,3.14,5.5}, xticklabels={$a$,$c$,$b$}, ytick=\empty, every axis/.append style={font=\small}]
\addplot[fill=popgreenL, draw=none, domain=1:3.14, samples=100]{2*sin(deg(x-1))} \closedcycle;
\addplot[fill=poppinkL, draw=none, domain=3.14:5.5, samples=100]{2*sin(deg(x-1))} \closedcycle;
\addplot[popcurve, domain=0.5:6, samples=200]{2*sin(deg(x-1))} node[pos=0.06, above, popdark]{$\mathcal{C}_f$};
\draw[popline, dashed] (axis cs:1,0) -- (axis cs:1,0);
\draw[popline, dashed] (axis cs:5.5,0) -- (axis cs:5.5,-1.79);
\node[popgreen!60!black] at (axis cs:2,1) {$+\,\mathcal{A}_1$};
\node[popink] at (axis cs:4.4,-1.1) {$-\,\mathcal{A}_2$};
\end{axis}
\end{tikzpicture}
\end{center}
\legende{Aire algébrique : la zone verte (au-dessus de l'axe) compte positivement, la zone rose (en dessous) négativement. Ici $\int_a^b f = \mathcal{A}_1 - \mathcal{A}_2$.}
\end{popfigure}

\begin{attention}[Intégrale $\neq$ aire géométrique]
C'est l'erreur la plus fréquente au Baccalauréat !
\begin{itemize}
\item L'\textbf{intégrale} $\displaystyle\int_a^b f(x)\,\mathrm{d}x$ est un nombre \emph{de signe quelconque} : elle peut être négative, voire nulle alors que la courbe n'est pas l'axe des abscisses (les aires positives et négatives se compensent).
\item L'\textbf{aire géométrique} du domaine entre la courbe et l'axe est \emph{toujours positive}. Pour la calculer quand $f$ change de signe, on découpe aux points où $f$ s'annule et on prend les valeurs absolues : $\mathcal{A} = \mathcal{A}_1 + \mathcal{A}_2$, et non $\mathcal{A}_1 - \mathcal{A}_2$.
\end{itemize}
Retiens : \emph{l'intégrale d'une fonction négative est négative}. Une aire, jamais.
\end{attention}

\courssub{Cas particuliers et conventions}

\begin{propriete}[Premiers cas particuliers]
Soit $f$ une fonction continue sur un intervalle $I$, et $a$, $b$ des éléments de $I$.
\begin{enumerate}
\item \textbf{Intégrale d'une constante :} pour tout réel $k$,
\[ \int_a^b k\,\mathrm{d}x = k\,(b - a). \]
C'est l'aire (algébrique) d'un rectangle de longueur $b - a$ et de hauteur $k$.
\item \textbf{Bornes égales :} $\displaystyle\int_a^a f(x)\,\mathrm{d}x = 0$ : le domaine est réduit à un segment vertical, d'aire nulle.
\item \textbf{Convention d'inversion des bornes :} on pose, par définition,
\[ \int_b^a f(x)\,\mathrm{d}x = -\int_a^b f(x)\,\mathrm{d}x. \]
Cette convention, purement formelle, rendra toutes les formules (notamment la relation de Chasles) valables quelles que soient les positions de $a$ et $b$.
\end{enumerate}
\end{propriete}

\begin{exemplebox}[Application immédiate]
$\displaystyle\int_1^5 3\,\mathrm{d}x = 3\times(5-1) = 12$ ; \quad $\displaystyle\int_2^2 \mathrm{e}^{x^2}\,\mathrm{d}x = 0$ (même sans connaître de primitive !) ; \quad $\displaystyle\int_7^4 x\,\mathrm{d}x = -\int_4^7 x\,\mathrm{d}x = -\frac{7^2 - 4^2}{2} = -\frac{33}{2}$.
\end{exemplebox}

\courssub{Interprétations concrètes : cinématique et débit}

Revenons à notre coopérative de Bafoussam. Le débit $d(t)$ (en L/h) est une fonction continue du temps $t$ (en heures). Pendant une petite durée $\Delta t$, le volume ajouté est environ $d(t)\times\Delta t$ : c'est l'aire d'un petit rectangle de base $\Delta t$ et de hauteur $d(t)$. En sommant toutes ces contributions entre 6 h et 18 h, on obtient :

\begin{aretenir}[Volume et distance comme intégrales]
\begin{itemize}
\item Si $d(t)$ est le \textbf{débit} d'un fluide à l'instant $t$, le volume total écoulé entre les instants $t_1$ et $t_2$ est $\displaystyle V = \int_{t_1}^{t_2} d(t)\,\mathrm{d}t$.
\item Si $v(t)$ est la \textbf{vitesse} d'un mobile à l'instant $t$, la distance parcourue entre $t_1$ et $t_2$ est $\displaystyle D = \int_{t_1}^{t_2} v(t)\,\mathrm{d}t$.
\end{itemize}
Dans les deux cas, l'intégrale « somme » des quantités infiniment petites : c'est exactement le problème de l'accroche. L'unité du résultat est le produit des unités : (L/h) $\times$ h $=$ L ; (km/h) $\times$ h $=$ km.
\end{aretenir}

\begin{exoresolu}[Le réservoir de la coopérative]
Le débit de remplissage du réservoir est modélisé, entre 6 h et 18 h, par $d(t) = 100$ L/h (débit supposé constant pour cette première approche). Quel volume d'eau s'est accumulé entre 6 h et 18 h ? Même question si, sur une période plus courte, le débit croît linéairement selon $d(t) = 20t$ L/h entre $t = 0$ et $t = 5$ h.
\tcblower
\textbf{Solution.} Le volume est l'intégrale du débit.

\textbf{Cas 1 : débit constant.}
\[ V = \int_6^{18} 100\,\mathrm{d}t = 100\times(18 - 6) = 100 \times 12 = 1\,200\ \text{L}. \]
Géométriquement, c'est l'aire d'un rectangle de longueur $12$ et de hauteur $100$.

\textbf{Cas 2 : débit linéaire.} La courbe de $d(t) = 20t$ est une droite passant par l'origine ; le domaine sous la courbe entre $0$ et $5$ est un triangle de base $5$ et de hauteur $d(5) = 100$ :
\[ V = \int_0^5 20t\,\mathrm{d}t = \frac{5 \times 100}{2} = 250\ \text{L}. \]
La coopérative a donc stocké $250$ litres pendant ces cinq heures.
\end{exoresolu}

\begin{savaistu}[Le « S » de Leibniz et les « fluxions » de Newton]
Le symbole $\displaystyle\int$ a été introduit par le philosophe et mathématicien allemand \textbf{Gottfried Wilhelm Leibniz} en 1675 : c'est un « S » allongé, initiale du latin \emph{summa} (« somme »), car l'intégrale est une somme infinie de quantités infiniment petites $f(x)\,\mathrm{d}x$. La notation $\mathrm{d}x$ est aussi due à Leibniz. De son côté, en Angleterre, \textbf{Isaac Newton} développait les mêmes idées avec un vocabulaire différent : il appelait « fluentes » les quantités variables et « fluxions » leurs vitesses de variation. La querelle de priorité entre les deux génies a divisé les mathématiciens européens pendant des décennies ! C'est finalement la notation de Leibniz, plus pratique, qui s'est imposée dans le monde entier — et que tu utilises aujourd'hui au lycée de Yaoundé ou de Douala.
\end{savaistu}

\begin{methode}[Bien démarrer avec les intégrales]
Face à une intégrale $\displaystyle\int_a^b f(x)\,\mathrm{d}x$ :
\begin{enumerate}
\item Vérifier que $f$ est \textbf{continue} sur $[a\,;\,b]$ (condition d'existence).
\item Se demander si le domaine est une \textbf{figure usuelle} (rectangle, triangle, trapèze, demi-disque) : si oui, calculer l'aire directement.
\item Sinon, \textbf{esquisser la courbe} et repérer le signe de $f$ pour interpréter le résultat (aire algébrique).
\item Ne jamais oublier que le résultat s'exprime en \textbf{unités d'aire} et qu'il peut être \textbf{négatif}.
\end{enumerate}
\end{methode}

Dans la section suivante, nous établirons les propriétés algébriques fondamentales de l'intégrale — linéarité, positivité, relation de Chasles — qui en font un véritable outil de calcul.

\coursec{Propriétés algébriques : linéarité, positivité, relation de Chasles}

Dans la section précédente, nous avons défini l'intégrale $\displaystyle\int_a^b f(t)\,\mathrm{d}t$ d'une fonction continue sur un segment $[a\,;b]$ : pour une fonction positive, c'est l'aire, exprimée en unités d'aire, du domaine situé sous la courbe ; dans le cas général, c'est une \emph{aire algébrique}. Nous avons aussi convenu que $\displaystyle\int_a^a f = 0$ et que $\displaystyle\int_b^a f = -\int_a^b f$.

Cette interprétation géométrique va maintenant nous servir de guide. Tout ce que l'on peut dire intuitivement sur les aires — on peut les additionner, les comparer, les découper, les encadrer — va devenir une propriété algébrique précise de l'intégrale. Ces propriétés sont les outils de travail quotidiens du calcul intégral : c'est grâce à elles que l'on pourra, dans les sections suivantes, calculer des intégrales compliquées en les décomposant en morceaux simples, ou encadrer une intégrale sans même la calculer. Prends le temps de bien les assimiler : elles tombent \emph{à tous les coups} au Baccalauréat.

Dans toute cette section, $f$ et $g$ désignent des fonctions \cle{continues} sur un intervalle $I$ de $\mathbb{R}$, et $a$, $b$, $c$ sont des éléments de $I$.

\courssub{Linéarité de l'intégrale}

\textbf{Intuition.} Imagine deux champs rectangulaires côte à côte sur la route de Bamenda : l'aire totale est la somme des deux aires. De même, si tu doubles la hauteur d'un mur, tu doubles la surface à peindre. L'intégrale, généralisation de la notion d'aire, doit donc respecter deux règles naturelles : l'intégrale d'une somme est la somme des intégrales, et l'intégrale d'un multiple est le multiple de l'intégrale.

\begin{propriete}[Linéarité de l'intégrale (admise)]
Pour toutes fonctions $f$ et $g$ continues sur $[a\,;b]$ et pour tous réels $\alpha$ et $\beta$ :
\[ \int_a^b \bigl(\alpha f(t) + \beta g(t)\bigr)\,\mathrm{d}t = \alpha\int_a^b f(t)\,\mathrm{d}t + \beta\int_a^b g(t)\,\mathrm{d}t. \]
En particulier : $\displaystyle\int_a^b (f+g) = \int_a^b f + \int_a^b g$ et $\displaystyle\int_a^b (\lambda f) = \lambda \int_a^b f$.
\end{propriete}

\textbf{Justification graphique pour des fonctions positives.} Supposons $f$ et $g$ positives sur $[a\,;b]$.

\begin{itemize}
\item Pour la \emph{somme} : la courbe de $f+g$ s'obtient en empilant, en chaque abscisse $t$, la hauteur $f(t)$ puis la hauteur $g(t)$. Le domaine sous la courbe de $f+g$ se découpe alors en deux parties : le domaine sous la courbe de $f$ (d'aire $\int_a^b f$) et le domaine compris entre les courbes de $f$ et de $f+g$, qui a exactement la même aire que le domaine sous la courbe de $g$ (on l'a simplement translaté verticalement). D'où $\int_a^b (f+g) = \int_a^b f + \int_a^b g$.
\item Pour le \emph{produit par un scalaire} $\lambda > 0$ : toutes les hauteurs sont multipliées par $\lambda$, donc l'aire est multipliée par $\lambda$ (comme pour le mur à peindre). Le cas $\lambda < 0$ s'obtient en remarquant que multiplier par $-1$ échange les aires algébriques au-dessus et au-dessous de l'axe.
\end{itemize}

\begin{exemplebox}[Linéarité en action]
Sachant que $\displaystyle\int_0^2 t^2\,\mathrm{d}t = \frac{8}{3}$ et $\displaystyle\int_0^2 t\,\mathrm{d}t = 2$, calculons $\displaystyle I = \int_0^2 \left(3t^2 - 5t\right)\mathrm{d}t$ sans connaître de primitive :
\[ I = 3\int_0^2 t^2\,\mathrm{d}t - 5\int_0^2 t\,\mathrm{d}t = 3\times\frac{8}{3} - 5\times 2 = 8 - 10 = -2. \]
Le résultat est négatif : l'aire algébrique sous l'axe l'emporte sur celle au-dessus.
\end{exemplebox}

\begin{attention}[La linéarité ne s'étend pas au produit]
En général, $\displaystyle\int_a^b (f \times g) \neq \left(\int_a^b f\right)\times\left(\int_a^b g\right)$. Par exemple, $\int_0^1 t\,\mathrm{d}t = \frac{1}{2}$, donc $\left(\int_0^1 t\,\mathrm{d}t\right)^2 = \frac{1}{4}$, alors que $\int_0^1 t^2\,\mathrm{d}t = \frac{1}{3}$. L'intégrale d'un produit demande des techniques spécifiques (intégration par parties, section 4).
\end{attention}

\courssub{Positivité et comparaison}

\textbf{Intuition.} Si une courbe reste au-dessus de l'axe des abscisses entre $a$ et $b$ (avec $a < b$), le domaine sous la courbe est une vraie aire, donc un nombre positif ou nul. C'est tout le contenu de la propriété suivante — et pourtant, on va en tirer des conséquences puissantes.

\begin{propriete}[Positivité de l'intégrale]
Soit $f$ une fonction continue sur $[a\,;b]$ avec $a < b$. Si $f(t) \geq 0$ pour tout $t \in [a\,;b]$, alors
\[ \int_a^b f(t)\,\mathrm{d}t \geq 0. \]
\end{propriete}

C'est une conséquence immédiate de la définition : l'intégrale d'une fonction positive sur $[a\,;b]$ (avec $a<b$) est l'aire d'un domaine, donc un réel positif.

\begin{propriete}[Comparaison des intégrales (théorème)]
Soient $f$ et $g$ continues sur $[a\,;b]$ avec $a < b$. Si $f(t) \leq g(t)$ pour tout $t \in [a\,;b]$, alors
\[ \int_a^b f(t)\,\mathrm{d}t \leq \int_a^b g(t)\,\mathrm{d}t. \]
\emph{La démonstration est exigible : elle illustre parfaitement la méthode « ramener une comparaison à une positivité ».}
\end{propriete}

\begin{experience}[Démonstration guidée du théorème de comparaison]
On suppose $f(t) \leq g(t)$ pour tout $t \in [a\,;b]$, avec $a < b$.
\begin{enumerate}
\item \textbf{Fabriquer une fonction positive.} Posons $h = g - f$. Par hypothèse, pour tout $t \in [a\,;b]$ : $h(t) = g(t) - f(t) \geq 0$. De plus, $h$ est continue sur $[a\,;b]$ comme différence de fonctions continues.
\item \textbf{Appliquer la positivité.} Comme $h \geq 0$ sur $[a\,;b]$ et $a < b$, la propriété de positivité donne : $\displaystyle\int_a^b h(t)\,\mathrm{d}t \geq 0$.
\item \textbf{Utiliser la linéarité.} Or $\displaystyle\int_a^b h = \int_a^b (g - f) = \int_a^b g - \int_a^b f$ par linéarité.
\item \textbf{Conclure.} On obtient $\displaystyle\int_a^b g - \int_a^b f \geq 0$, c'est-à-dire $\displaystyle\int_a^b f \leq \int_a^b g$. \hfill$\blacksquare$
\end{enumerate}
\end{experience}

\begin{popfigure}
\begin{center}
\begin{tikzpicture}
\begin{axis}[
  width=0.85\linewidth, height=6cm,
  axis lines=middle,
  xlabel={$x$}, ylabel={$y$},
  xmin=-0.3, xmax=3.6, ymin=-0.4, ymax=3.2,
  xtick={0.5,3}, xticklabels={$a$,$b$},
  ytick=\empty,
  clip=false
]
\addplot[popcurve, domain=0.2:3.3, samples=150, name path=G]{0.35*x^2 - 0.8*x + 2.2};
\addplot[poporange, very thick, domain=0.2:3.3, samples=150, name path=F]{0.15*x^2 - 0.3*x + 0.9};
\addplot[popblue!25] fill between[of=F and G, soft clip={domain=0.5:3}];
\path[name path=AXE] (axis cs:0,0) -- (axis cs:3.6,0);
\addplot[poporange!25] fill between[of=F and AXE, soft clip={domain=0.5:3}];
\draw[dashed, popmuted] (axis cs:0.5,0) -- (axis cs:0.5,1.8875);
\draw[dashed, popmuted] (axis cs:3,0) -- (axis cs:3,2.95);
\node[popblue] at (axis cs:2.9,2.6) {$\mathcal{C}_g$};
\node[poporange] at (axis cs:1.1,0.75) {$\mathcal{C}_f$};
\node[popdark, font=\small] at (axis cs:1.75,1.05) {$\displaystyle\int_a^b f$};
\node[popdark, font=\small] at (axis cs:2.3,1.85) {$\displaystyle\int_a^b (g-f)$};
\end{axis}
\end{tikzpicture}
\end{center}
\legende{Si $f \leq g$ sur $[a\,;b]$, l'aire sous $\mathcal{C}_f$ est inférieure à l'aire sous $\mathcal{C}_g$ : la différence est l'aire du domaine situé entre les deux courbes.}
\end{popfigure}

\begin{attention}[Ordre des bornes : le piège classique du Bac]
La positivité et la comparaison exigent impérativement $a < b$. Si l'on inverse les bornes, l'inégalité \emph{se renverse}, car $\displaystyle\int_b^a f = -\int_a^b f$. Par exemple, $t^2 \geq 0$ donne $\int_0^2 t^2\,\mathrm{d}t = \frac{8}{3} \geq 0$, mais $\int_2^0 t^2\,\mathrm{d}t = -\frac{8}{3} \leq 0$. Avant d'intégrer une inégalité, vérifie toujours que la borne inférieure est bien la plus petite.
\end{attention}

\courssub{Intégration des inégalités : encadrement d'une intégrale}

\textbf{Intuition.} Si sur tout l'intervalle $[a\,;b]$ la courbe de $f$ reste coincée entre deux droites horizontales $y = m$ et $y = M$, alors l'aire sous la courbe est comprise entre les aires de deux rectangles de même largeur $b - a$ et de hauteurs $m$ et $M$.

\begin{propriete}[Intégration des inégalités]
Soit $f$ continue sur $[a\,;b]$ avec $a < b$. Si pour tout $t \in [a\,;b]$ on a $m \leq f(t) \leq M$ (où $m$ et $M$ sont des constantes), alors
\[ m(b-a) \leq \int_a^b f(t)\,\mathrm{d}t \leq M(b-a). \]
\end{propriete}

\textbf{Démonstration.} L'inégalité $m \leq f(t)$ s'écrit : la fonction constante $t \mapsto m$ est inférieure à $f$. Par comparaison (théorème précédent), $\displaystyle\int_a^b m\,\mathrm{d}t \leq \int_a^b f(t)\,\mathrm{d}t$. Or $\displaystyle\int_a^b m\,\mathrm{d}t = m(b-a)$ (aire d'un rectangle). On procède de même avec $f(t) \leq M$. \hfill$\blacksquare$

\begin{popfigure}
\begin{center}
\begin{tikzpicture}
\begin{axis}[
  width=0.85\linewidth, height=6cm,
  axis lines=middle,
  xlabel={$x$}, ylabel={$y$},
  xmin=-0.3, xmax=3.8, ymin=-0.3, ymax=3.4,
  xtick={0.6,3.2}, xticklabels={$a$,$b$},
  ytick={1.0,2.8}, yticklabels={$m$,$M$},
  clip=false
]
\addplot[draw=none, forget plot, name path=hconst1, domain=0.6:3.2] {2.8};
\addplot[draw=none, forget plot, name path=hconst2, domain=0.6:3.2] {0};
\addplot[popgold!35, draw=none] fill between[of=hconst1 and hconst2, soft clip={domain=0.6:3.2}];
\addplot[draw=none, forget plot, name path=hconst3, domain=0.6:3.2] {1.0};
\addplot[draw=none, forget plot, name path=hconst4, domain=0.6:3.2] {0};
\addplot[popgreen!25, draw=none] fill between[of=hconst3 and hconst4, soft clip={domain=0.6:3.2}];
\addplot[popcurve, domain=0.3:3.5, samples=150]{1.9 + 0.7*sin(deg(2.2*x - 1.5))};
\draw[popgold, thick, dashed] (axis cs:0.6,2.8) -- (axis cs:3.2,2.8);
\draw[popgreen, thick, dashed] (axis cs:0.6,1.0) -- (axis cs:3.2,1.0);
\draw[dashed, popmuted] (axis cs:0.6,0) -- (axis cs:0.6,2.8);
\draw[dashed, popmuted] (axis cs:3.2,0) -- (axis cs:3.2,2.8);
\node[popdark, font=\small] at (axis cs:1.9,1.9) {$\displaystyle\int_a^b f$};
\node[popgreen!60!black, font=\small] at (axis cs:1.9,0.5) {$m(b-a)$};
\node[popgold!70!black, font=\small] at (axis cs:1.9,3.0) {rectangle de hauteur $M$};
\end{axis}
\end{tikzpicture}
\end{center}
\legende{L'aire sous la courbe est comprise entre l'aire du rectangle de hauteur $m$ (vert) et celle du rectangle de hauteur $M$ (jaune).}
\end{popfigure}

\begin{exemplebox}[Encadrer sans calculer]
Montrons que $\displaystyle \frac{4}{3} \leq \int_1^3 \frac{x^2+1}{x^2+2}\,\mathrm{d}x \leq \frac{20}{11}$.

Posons, pour tout $x \in [1\,;3]$, $f(x) = \dfrac{x^2+1}{x^2+2}$. La démarche se fait en trois étapes.

\textbf{Étape 1 : étudier les variations de $f$ sur $[1\,;3]$.} Écrivons
\[ f(x) = \frac{x^2+2-1}{x^2+2} = 1 - \frac{1}{x^2+2}. \]
Sur $[1\,;3]$, la fonction $x \mapsto x^2 + 2$ est strictement croissante et à valeurs strictement positives, donc $x \mapsto \dfrac{1}{x^2+2}$ est strictement décroissante, et par suite $f$ est strictement \emph{croissante} sur $[1\,;3]$.

\textbf{Étape 2 : en déduire un encadrement de $f(x)$.} Pour tout $x \in [1\,;3]$ :
\[ f(1) \leq f(x) \leq f(3) \quad\text{c'est-à-dire}\quad \frac{2}{3} \leq f(x) \leq \frac{10}{11}, \]
car $f(1) = \dfrac{1+1}{1+2} = \dfrac{2}{3}$ et $f(3) = \dfrac{9+1}{9+2} = \dfrac{10}{11}$.

\textbf{Étape 3 : intégrer l'encadrement.} La fonction $f$ est continue sur $[1\,;3]$ et $1 < 3$ : on peut appliquer l'intégration des inégalités avec $m = \dfrac{2}{3}$, $M = \dfrac{10}{11}$ et $b - a = 2$ :
\[ \frac{2}{3}\times(3-1) \leq \int_1^3 f(x)\,\mathrm{d}x \leq \frac{10}{11}\times(3-1), \quad\text{soit}\quad \frac{4}{3} \leq \int_1^3 \frac{x^2+1}{x^2+2}\,\mathrm{d}x \leq \frac{20}{11}. \]
Numériquement : $\dfrac{4}{3} \approx 1{,}33$ et $\dfrac{20}{11} \approx 1{,}82$ : l'intégrale vaut donc environ $1{,}6$, sans que nous ayons eu besoin de la calculer exactement.
\end{exemplebox}

\begin{methode}[Encadrer une intégrale sans la calculer]
Pour encadrer $\displaystyle\int_a^b f(x)\,\mathrm{d}x$ (avec $a < b$) :
\begin{enumerate}
\item Encadrer $f(x)$ sur $[a\,;b]$ par deux constantes $m$ et $M$, en utilisant les variations de $f$, la composée avec une fonction de référence (inverse, carré, racine\ldots) ou des inégalités connues.
\item Vérifier que $f$ est continue sur $[a\,;b]$ et que $a < b$.
\item Multiplier l'encadrement par la largeur $b - a$ : $\; m(b-a) \leq \displaystyle\int_a^b f \leq M(b-a)$.
\end{enumerate}
\end{methode}

\begin{exoresolu}[Encadrement d'une intégrale]
\textbf{Énoncé.} Soit $I = \displaystyle\int_0^1 \frac{\mathrm{d}t}{1+t^2}$. Montrer que $\dfrac{1}{2} \leq I \leq 1$.
\tcblower
\textbf{Solution.} Pour tout $t \in [0\,;1]$, on a $0 \leq t^2 \leq 1$, donc $1 \leq 1 + t^2 \leq 2$. La fonction inverse étant décroissante sur $]0\,;+\infty[$, on obtient
\[ \frac{1}{2} \leq \frac{1}{1+t^2} \leq 1. \]
La fonction $t \mapsto \dfrac{1}{1+t^2}$ est continue sur $[0\,;1]$ et $0 < 1$ : on peut intégrer cet encadrement :
\[ \frac{1}{2}(1-0) \leq \int_0^1 \frac{\mathrm{d}t}{1+t^2} \leq 1\times(1-0), \quad\text{soit}\quad \frac{1}{2} \leq I \leq 1. \]
(On verra plus tard que $I = \dfrac{\pi}{4} \approx 0{,}785$ : l'encadrement est bien vérifié.)
\end{exoresolu}

\courssub{Relation de Chasles}

\textbf{Intuition.} Un commerçant de Douala mesure la surface de sa boutique en deux lots : la salle de vente et la réserve. L'aire totale est la somme des deux aires. De même, l'aire sous une courbe entre $a$ et $c$ est la somme des aires entre $a$ et $b$ puis entre $b$ et $c$.

\begin{propriete}[Relation de Chasles (admise dans le cas général)]
Soit $f$ continue sur un intervalle $I$. Pour tous $a$, $b$, $c$ dans $I$ :
\[ \int_a^b f(t)\,\mathrm{d}t + \int_b^c f(t)\,\mathrm{d}t = \int_a^c f(t)\,\mathrm{d}t. \]
\end{propriete}

\textbf{Justification graphique} dans le cas $a < b < c$ et $f$ positive : le domaine sous la courbe entre $a$ et $c$ est la \emph{réunion disjointe} (à une frontière près, d'aire nulle) du domaine entre $a$ et $b$ et du domaine entre $b$ et $c$. Les aires s'ajoutent donc.

\begin{popfigure}
\begin{center}
\begin{tikzpicture}
\begin{axis}[
  width=0.85\linewidth, height=6cm,
  axis lines=middle,
  xlabel={$x$}, ylabel={$y$},
  xmin=-0.3, xmax=5.4, ymin=-0.3, ymax=3.4,
  xtick={0.7,2.6,4.6}, xticklabels={$a$,$b$,$c$},
  ytick=\empty,
  clip=false
]
\addplot[popcurve, domain=0.3:5.0, samples=200, name path=F]{1.6 + 0.6*sin(deg(1.6*x)) + 0.25*x};
\addplot[draw=none, forget plot, name path=hconst5, domain=0.7:2.6] {0};
\addplot[popblue!30, draw=none] fill between[of=F and hconst5, soft clip={domain=0.7:2.6}];
\addplot[draw=none, forget plot, name path=hconst6, domain=2.6:4.6] {0};
\addplot[poporange!30, draw=none] fill between[of=F and hconst6, soft clip={domain=2.6:4.6}];
\draw[dashed, popmuted] (axis cs:0.7,0) -- (axis cs:0.7,2.35);
\draw[dashed, popmuted] (axis cs:2.6,0) -- (axis cs:2.6,2.1);
\draw[dashed, popmuted] (axis cs:4.6,0) -- (axis cs:4.6,2.5);
\node[popblue, font=\small] at (axis cs:1.6,0.8) {$\mathcal{A}_1 = \displaystyle\int_a^b f$};
\node[poporange!80!black, font=\small] at (axis cs:3.7,0.9) {$\mathcal{A}_2 = \displaystyle\int_b^c f$};
\node[popdark] at (axis cs:4.9,2.7) {$\mathcal{C}_f$};
\end{axis}
\end{tikzpicture}
\end{center}
\legende{Relation de Chasles : l'aire totale entre $a$ et $c$ est la somme des aires des deux domaines adjacents : $\mathcal{A}_1 + \mathcal{A}_2 = \displaystyle\int_a^c f$.}
\end{popfigure}

\begin{attention}[Chasles fonctionne dans tous les ordres]
Contrairement à la positivité, la relation de Chasles est valable \emph{quel que soit l'ordre} de $a$, $b$, $c$, grâce à la convention $\int_b^a f = -\int_a^b f$. Par exemple, si $a < c < b$ : $\int_a^b f = \int_a^c f + \int_c^b f$ reste vraie. C'est exactement l'analogue intégral de la relation de Chasles sur les vecteurs : $\overrightarrow{AB} + \overrightarrow{BC} = \overrightarrow{AC}$.
\end{attention}

\courssub{Applications : découpage selon le signe et fonctions par morceaux}

La relation de Chasles est l'outil qui rend calculables deux situations très fréquentes au Baccalauréat.

\textbf{Première situation : découper selon le signe de $f$.} Pour calculer l'\emph{aire totale} (positive) entre une courbe et l'axe des abscisses lorsque $f$ change de signe, on découpe l'intégrale aux points où $f$ s'annule, et l'on prend l'opposé des morceaux où $f$ est négative.

\textbf{Deuxième situation : fonctions définies par morceaux.} Si l'expression de $f$ change en un point de raccordement, on découpe l'intégrale à ce point.

\begin{methode}[Intégrale d'une fonction définie par morceaux]
Pour calculer $\displaystyle\int_a^b f$ lorsque $f$ est définie par des expressions différentes sur des sous-intervalles :
\begin{enumerate}
\item Repérer les \cle{points de raccordement} $x_1 < x_2 < \cdots$ où l'expression de $f$ change, situés dans $[a\,;b]$.
\item Écrire, par la relation de Chasles : $\displaystyle\int_a^b f = \int_a^{x_1} f + \int_{x_1}^{x_2} f + \cdots + \int_{x_k}^{b} f$.
\item Sur chaque morceau, remplacer $f$ par l'expression qui y correspond et calculer séparément.
\item Additionner les résultats.
\end{enumerate}
\end{methode}

\begin{exoresolu}[Fonction définie par morceaux]
\textbf{Énoncé.} Une entreprise de livraison de Yaoundé facture sa consommation d'énergie selon le tarif ENEO suivant (modélisé) : la puissance absorbée en kW pendant une journée est donnée par
\[ f(t) = \begin{cases} 2 & \text{si } 0 \leq t \leq 4 \\ 6 - t & \text{si } 4 < t \leq 6 \end{cases} \]
où $t$ est le temps en heures. Calculer $\displaystyle\int_0^6 f(t)\,\mathrm{d}t$.
\tcblower
\textbf{Solution.} Le point de raccordement est $t = 4$. Par la relation de Chasles :
\begin{align*}
\int_0^6 f(t)\,\mathrm{d}t &= \int_0^4 f(t)\,\mathrm{d}t + \int_4^6 f(t)\,\mathrm{d}t \\
&= \int_0^4 2\,\mathrm{d}t + \int_4^6 (6-t)\,\mathrm{d}t.
\end{align*}
La première intégrale est l'aire d'un rectangle : $\int_0^4 2\,\mathrm{d}t = 2 \times 4 = 8$.
La seconde est l'aire d'un triangle de base $2$ et de hauteur $2$ (car $6 - t$ décroît de $2$ à $0$ sur $[4\,;6]$) : $\int_4^6 (6-t)\,\mathrm{d}t = \dfrac{2 \times 2}{2} = 2$.
Donc $\displaystyle\int_0^6 f(t)\,\mathrm{d}t = 8 + 2 = 10$ : l'énergie consommée sur les six heures est de $10$ kWh.
\end{exoresolu}

\begin{exoresolu}[Découpage selon le signe]
\textbf{Énoncé.} On donne $\displaystyle\int_0^1 (2t-1)\,\mathrm{d}t = 0$ et $\displaystyle\int_0^{1/2} (2t-1)\,\mathrm{d}t = -\frac{1}{4}$. En déduire l'aire $\mathcal{A}$ du domaine compris entre la droite d'équation $y = 2x - 1$, l'axe des abscisses et les droites $x = 0$ et $x = 1$.
\tcblower
\textbf{Solution.} La fonction $f(t) = 2t - 1$ s'annule en $t = \dfrac{1}{2}$ ; elle est négative sur $\left[0\,;\dfrac{1}{2}\right]$ et positive sur $\left[\dfrac{1}{2}\,;1\right]$. L'intégrale $\int_0^1 f = 0$ ne donne \emph{pas} l'aire : les deux morceaux se compensent algébriquement. Par Chasles :
\[ \int_0^1 f = \int_0^{1/2} f + \int_{1/2}^1 f \implies 0 = -\frac{1}{4} + \int_{1/2}^1 f \implies \int_{1/2}^1 f = \frac{1}{4}. \]
L'aire totale est la somme des aires (positives) des deux morceaux :
\[ \mathcal{A} = -\int_0^{1/2} f + \int_{1/2}^1 f = \frac{1}{4} + \frac{1}{4} = \frac{1}{2} \text{ unité d'aire}. \]
\end{exoresolu}

\begin{attention}[Intégrale nulle ne signifie pas aire nulle]
Si $f$ change de signe sur $[a\,;b]$, l'intégrale $\int_a^b f$ peut être nulle alors que le domaine entre la courbe et l'axe a une aire strictement positive. Pour obtenir l'aire géométrique, il faut découper par Chasles aux zéros de $f$ et prendre la valeur absolue de chaque morceau — c'est le rôle de l'intégrale de $|f|$, que nous étudierons dans la section suivante.
\end{attention}

\begin{aretenir}[L'essentiel de la section]
\begin{itemize}
\item \textbf{Linéarité} : $\displaystyle\int_a^b (\alpha f + \beta g) = \alpha \int_a^b f + \beta \int_a^b g$.
\item \textbf{Positivité} ($a < b$) : $f \geq 0 \implies \displaystyle\int_a^b f \geq 0$.
\item \textbf{Comparaison} ($a < b$) : $f \leq g \implies \displaystyle\int_a^b f \leq \int_a^b g$ (se démontre par positivité de $g - f$).
\item \textbf{Intégration des inégalités} ($a < b$) : $m \leq f \leq M \implies m(b-a) \leq \displaystyle\int_a^b f \leq M(b-a)$.
\item \textbf{Chasles} : $\displaystyle\int_a^b f + \int_b^c f = \int_a^c f$, valable dans tous les ordres.
\item Vigilance permanente : positivité et comparaison exigent $a < b$ ; inverser les bornes change le signe de l'intégrale.
\end{itemize}
\end{aretenir}

\begin{savaistu}[Chasles, un nom de géomètre]
La « relation de Chasles » porte le nom de Michel Chasles (1793--1880), mathématicien français célèbre pour ses travaux de géométrie. La relation vectorielle $\overrightarrow{AB} + \overrightarrow{BC} = \overrightarrow{AC}$ que tu as rencontrée en Seconde et la relation intégrale $\int_a^b + \int_b^c = \int_a^c$ expriment la même idée profonde : une grandeur \emph{additive} le long d'un chemin ne dépend que du point de départ et du point d'arrivée. C'est cette même idée qui fait fonctionner les compteurs d'énergie d'ENEO ou les applications de suivi de distance des motos-taxis : on additionne les petites contributions successives pour obtenir le total.
\end{savaistu}

\begin{exercice}
Soit $f$ la fonction définie sur $[0\,;5]$ par $f(x) = |x - 3|$.
\begin{enumerate}
\item Tracer la courbe de $f$ et déterminer le signe de $x - 3$ selon les valeurs de $x$.
\item En découpant avec la relation de Chasles au point d'abscisse $3$, calculer $\displaystyle\int_0^5 f(x)\,\mathrm{d}x$ à l'aide d'aires de triangles.
\item Montrer, sans nouveau calcul, que $\displaystyle 0 \leq \int_1^4 f(x)\,\mathrm{d}x \leq 6$.
\end{enumerate}
\end{exercice}

\begin{corrige}[Exercice]
\begin{enumerate}
\item $f(x) = x - 3$ si $x \geq 3$ et $f(x) = 3 - x$ si $x \leq 3$. La courbe est formée de deux demi-droites se raccordant au point $(3\,;0)$, en forme de « V ».
\item Par Chasles : $\displaystyle\int_0^5 f = \int_0^3 (3-x)\,\mathrm{d}x + \int_3^5 (x-3)\,\mathrm{d}x$. Le premier morceau est l'aire d'un triangle de base $3$ et de hauteur $3$ : $\dfrac{3 \times 3}{2} = \dfrac{9}{2}$. Le second est l'aire d'un triangle de base $2$ et de hauteur $2$ : $\dfrac{2 \times 2}{2} = 2$. Donc $\displaystyle\int_0^5 f = \frac{9}{2} + 2 = \frac{13}{2}$.
\item Sur $[1\,;4]$, $f$ est positive, donc l'intégrale est positive. De plus $f(x) = |x-3| \leq 2$ sur $[1\,;4]$ (le maximum est atteint en $x = 1$, où $f(1) = 2$). Par intégration des inégalités : $\displaystyle\int_1^4 f \leq 2 \times (4-1) = 6$. D'où $0 \leq \displaystyle\int_1^4 f \leq 6$.
\end{enumerate}
\end{corrige}

\coursec{Fonction définie par une intégrale et valeur moyenne}

Dans la section précédente, nous avons établi les propriétés algébriques de l'intégrale : linéarité, positivité et relation de Chasles. Nous allons maintenant exploiter ces outils pour répondre à une question nouvelle et passionnante : que se passe-t-il lorsque la borne supérieure de l'intégrale devient \emph{variable} ? L'intégrale $\int_a^x f(t)\,\mathrm{d}t$ dépend alors de $x$ : elle définit une \cle{fonction} de $x$, dont nous allons étudier le sens de variation \emph{sans utiliser la notion de primitive}, comme l'exige le programme. Nous définirons ensuite la \cle{valeur moyenne} d'une fonction sur un intervalle, notion très concrète (vitesse moyenne d'un taxi-brousse entre Douala et Yaoundé, débit moyen du fleuve Wouri), et nous terminerons par le calcul de l'intégrale d'une valeur absolue.

\courssub{Fonction définie par une intégrale}

\textbf{Une idée simple pour commencer.}\par
Imagine un compteur d'eau de la CAMWATER : à chaque instant $t$, le débit instantané est $f(t)$ (en litres par heure). La quantité totale d'eau consommée entre l'instant initial $a$ et l'instant $x$ est $\int_a^x f(t)\,\mathrm{d}t$. Quand $x$ augmente, cette quantité augmente aussi (puisque le débit est positif) : nous venons de définir une fonction croissante de $x$. C'est exactement l'idée de cette section.

\begin{definition}[Fonction définie par une intégrale]
Soit $f$ une fonction \cle{continue} sur un intervalle $I$ et $a$ un point de $I$. On définit sur $I$ la fonction $F$ par :
\[ F(x) = \int_a^x f(t)\,\mathrm{d}t. \]
En particulier, $F(a) = \int_a^a f(t)\,\mathrm{d}t = 0$.
\end{definition}

\begin{attention}[La variable d'intégration est muette]
Dans l'écriture $\int_a^x f(t)\,\mathrm{d}t$, la variable $t$ est une \cle{variable muette} : on aurait pu écrire $\int_a^x f(u)\,\mathrm{d}u$ ou $\int_a^x f(s)\,\mathrm{d}s$, le résultat serait identique. En revanche, il est \textbf{interdit} d'écrire $\int_a^x f(x)\,\mathrm{d}x$ : la lettre $x$ désignerait alors à la fois la borne (qui varie) et la variable d'intégration (qui parcourt $[a\,;x]$), ce qui est une confusion grave. Retiens : \emph{la borne s'appelle $x$, la variable d'intégration s'appelle $t$}.
\end{attention}

\textbf{Exemple.} Soit $f(t) = 3$ (fonction constante) et $a = 0$. Alors $F(x) = \int_0^x 3\,\mathrm{d}t = 3x$ (aire d'un rectangle de longueur $x$ et de hauteur $3$). On vérifie bien $F(0) = 0$, et $F$ est croissante car $f \geq 0$. Ce constat est général : c'est l'objet du théorème fondamental qui suit.

\courssub{Sens de variation de $F$ : le signe de $f$ commande}

Voici le théorème central de cette section. Sa démonstration est \textbf{exigible} : elle n'utilise que la relation de Chasles et la positivité de l'intégrale, vues dans la section précédente.

\begin{propriete}[Théorème : signe de $f$ et variations de $F$ — démonstration exigible]
Soit $f$ une fonction continue sur un intervalle $I$, $a \in I$, et $F(x) = \int_a^x f(t)\,\mathrm{d}t$.
\begin{itemize}
\item Si $f$ est \textbf{positive} sur $I$, alors $F$ est \textbf{croissante} sur $I$.
\item Si $f$ est \textbf{négative} sur $I$, alors $F$ est \textbf{décroissante} sur $I$.
\end{itemize}
\end{propriete}

\begin{experience}[Démonstration guidée du théorème]
On suppose $f$ positive sur $I$. Soient $x_1$ et $x_2$ deux points de $I$ tels que $x_1 < x_2$. On veut comparer $F(x_1)$ et $F(x_2)$.

\textbf{Étape 1 : exprimer la différence.} Par la relation de Chasles :
\[ F(x_2) - F(x_1) = \int_a^{x_2} f(t)\,\mathrm{d}t - \int_a^{x_1} f(t)\,\mathrm{d}t = \int_{x_1}^{x_2} f(t)\,\mathrm{d}t. \]

\textbf{Étape 2 : appliquer la positivité.} Sur l'intervalle $[x_1\,;x_2]$ (avec $x_1 < x_2$), la fonction $f$ est positive. Par positivité de l'intégrale :
\[ \int_{x_1}^{x_2} f(t)\,\mathrm{d}t \geq 0. \]

\textbf{Étape 3 : conclure.} On obtient $F(x_2) - F(x_1) \geq 0$, c'est-à-dire $F(x_1) \leq F(x_2)$. La fonction $F$ est donc croissante sur $I$. $\square$

Le cas $f$ négative se traite de la même façon : l'intégrale $\int_{x_1}^{x_2} f(t)\,\mathrm{d}t$ est alors négative, donc $F(x_2) \leq F(x_1)$ et $F$ est décroissante.
\end{experience}

\begin{aretenir}[Règle pratique]
Le \cle{sens de variation} de $F(x) = \int_a^x f(t)\,\mathrm{d}t$ est donné par le \cle{signe} de $f$ :
\begin{center}
\begin{tabularx}{\linewidth}{|l|X|X|}
\hline
\rowcolor{popblueL} Signe de $f$ sur $I$ & Variation de $F$ sur $I$ & Image mentale \\
\hline
$f \geq 0$ & $F$ croissante & l'aire s'accumule \\
\hline
$f \leq 0$ & $F$ décroissante & l'aire « négative » s'accumule \\
\hline
$f$ change de signe en $c$ & $F$ change de variation en $c$ & extremum de $F$ en $c$ \\
\hline
\end{tabularx}
\end{center}
\end{aretenir}

\textbf{Lecture graphique : construire le tableau de variations de $F$.}\par
Lorsque l'énoncé donne la courbe de $f$ (et non son expression), on lit le signe de $f$ (position de la courbe par rapport à l'axe des abscisses) et on en déduit les variations de $F$, en n'oubliant pas que $F(a) = 0$.

\begin{exoresolu}[Tableau de variations de $F$ à partir de la courbe de $f$]
\textbf{Énoncé.} Soit $f$ la fonction définie sur $[0\,;4]$ par $f(x) = x - 2$, et $F(x) = \int_0^x f(t)\,\mathrm{d}t$. Étudier le sens de variation de $F$ sur $[0\,;4]$ et dresser son tableau de variations (on admettra les valeurs $F(2) = -2$ et $F(4) = 0$, qui seront justifiées après l'étude des primitives).
\tcblower
\textbf{Solution.} La fonction $f$ est continue sur $[0\,;4]$ (fonction affine), donc $F$ est bien définie sur $[0\,;4]$.

\textbf{Signe de $f$ :} $f(x) = x - 2 \geq 0 \iff x \geq 2$. Donc $f \leq 0$ sur $[0\,;2]$ et $f \geq 0$ sur $[2\,;4]$.

\textbf{Variations de $F$ :} d'après le théorème, $F$ est décroissante sur $[0\,;2]$ et croissante sur $[2\,;4]$. De plus $F(0) = 0$. D'où le tableau :
\begin{center}
\begin{tabular}{|c|ccccc|}
\hline
$x$ & $0$ & \hspace{1.2cm} & $2$ & \hspace{1.2cm} & $4$ \\
\hline
Signe de $f$ & \multicolumn{2}{c|}{$-$} & $0$ & \multicolumn{2}{c|}{$+$} \\
\hline
Variations de $F$ & $0$ & $\searrow$ & $-2$ & $\nearrow$ & $0$ \\
\hline
\end{tabular}
\end{center}
$F$ admet donc un minimum en $x = 2$, égal à $-2$. Remarque que $F$ peut prendre des valeurs négatives : une intégrale n'est pas toujours une aire positive !
\end{exoresolu}

\begin{popfigure}
\begin{tikzpicture}
\begin{axis}[width=0.85\linewidth, height=5.5cm, axis lines=middle, xlabel={$x$}, ylabel={$f(x)$}, xmin=-0.3, xmax=4.4, ymin=-2.6, ymax=2.6, xtick={0,2,4}, ytick={-2,2}, domain=0:4]
\addplot[popcurve, samples=100] {x-2};
\addplot[poporange, dashed, domain=0:4, samples=2] {0};
\node[popblue] at (axis cs:3.4,1.9) {$\mathcal{C}_f$};
\node[popmuted, font=\small] at (axis cs:1,-0.5) {$f<0$};
\node[popmuted, font=\small] at (axis cs:3.2,0.5) {$f>0$};
\end{axis}
\end{tikzpicture}

\vspace{0.3cm}
\begin{center}
\begin{tabular}{|c|ccccc|}
\hline
$x$ & $0$ & \hspace{1cm} & $2$ & \hspace{1cm} & $4$ \\
\hline
Signe de $f$ & \multicolumn{2}{c|}{$-$} & $0$ & \multicolumn{2}{c|}{$+$} \\
\hline
$F$ & $0$ & $\searrow$ & $-2$ & $\nearrow$ & $0$ \\
\hline
\end{tabular}
\end{center}
\legende{Correspondance entre le signe de $f$ (courbe) et les variations de $F(x)=\int_0^x f(t)\,\mathrm{d}t$ (tableau).}
\end{popfigure}

\begin{exemplebox}[Le cas du sinus]
Soit $f(x) = \sin x$ sur $[0\,;\pi]$ et $F(x) = \int_0^x \sin t\,\mathrm{d}t$. Comme $\sin t \geq 0$ pour tout $t \in [0\,;\pi]$, la fonction $F$ est \textbf{croissante} sur $[0\,;\pi]$, avec $F(0) = 0$. Nous calculerons plus loin (grâce aux primitives) que $F(x) = 1 - \cos x$, ce qui confirmera ce résultat : $F$ passe de $0$ à $2$ en croissant.
\end{exemplebox}

\courssub{Inégalité de la moyenne et valeur moyenne}

\textbf{Intuition.}\par
Si la température à Bafoussam reste toute la journée entre $18\,^\circ\mathrm{C}$ et $26\,^\circ\mathrm{C}$, alors la température moyenne de la journée est forcément comprise entre $18$ et $26$. Cette évidence se traduit mathématiquement par l'\cle{inégalité de la moyenne}.

\begin{propriete}[Inégalité de la moyenne]
Soit $f$ une fonction continue sur $[a\,;b]$ ($a < b$) telle que, pour tout $x \in [a\,;b]$ :
\[ m \leq f(x) \leq M. \]
Alors :
\[ m(b-a) \leq \int_a^b f(x)\,\mathrm{d}x \leq M(b-a), \qquad \text{c'est-à-dire} \qquad m \leq \frac{1}{b-a}\int_a^b f(x)\,\mathrm{d}x \leq M. \]
\end{propriete}

\begin{experience}[Démonstration de l'inégalité de la moyenne]
\textbf{Étape 1 : intégrer l'encadrement.} Pour tout $x \in [a\,;b]$, on a $m \leq f(x) \leq M$. Par positivité de l'intégrale (l'intégration conserve les inégalités larges sur $[a\,;b]$ avec $a < b$) :
\[ \int_a^b m\,\mathrm{d}x \leq \int_a^b f(x)\,\mathrm{d}x \leq \int_a^b M\,\mathrm{d}x. \]
\textbf{Étape 2 : calculer les intégrales des constantes.} L'intégrale d'une constante $k$ sur $[a\,;b]$ vaut $k(b-a)$ (aire d'un rectangle) :
\[ m(b-a) \leq \int_a^b f(x)\,\mathrm{d}x \leq M(b-a). \]
\textbf{Étape 3 : diviser par $b - a > 0$.} On obtient l'inégalité annoncée. $\square$
\end{experience}

\begin{definition}[Valeur moyenne d'une fonction]
Soit $f$ une fonction continue sur $[a\,;b]$ ($a < b$). On appelle \cle{valeur moyenne} de $f$ sur $[a\,;b]$ le réel :
\[ \mu = \frac{1}{b-a}\int_a^b f(x)\,\mathrm{d}x. \]
D'après l'inégalité de la moyenne, si $m \leq f \leq M$ sur $[a\,;b]$, alors $m \leq \mu \leq M$.
\end{definition}

\textbf{Interprétation graphique.}\par
L'égalité $\int_a^b f(x)\,\mathrm{d}x = \mu \times (b-a)$ signifie que, pour une fonction positive, l'aire sous la courbe est égale à l'aire du \cle{rectangle} de base $[a\,;b]$ et de hauteur $\mu$ : la valeur moyenne est la « hauteur constante équivalente ».

\begin{popfigure}
\begin{tikzpicture}
\begin{axis}[width=0.85\linewidth, height=6cm, axis lines=middle, xlabel={$x$}, ylabel={$y$}, xmin=-0.3, xmax=4.4, ymin=-0.3, ymax=3.2, xtick={1,4}, xticklabels={$a$,$b$}, ytick=\empty, domain=1:4]
\addplot[popblueL, draw=none, domain=1:4, samples=100] {1.5+0.8*sin(deg(x-1)*1.2)+0.2*(x-1)} \closedcycle;
\addplot[popcurve, samples=150] {1.5+0.8*sin(deg(x-1)*1.2)+0.2*(x-1)};
\addplot[poporange, thick, dashed, domain=1:4, samples=2] {2.1};
\draw[poporange, thick, dashed] (axis cs:1,0) -- (axis cs:1,2.1);
\draw[poporange, thick, dashed] (axis cs:4,0) -- (axis cs:4,2.1);
\node[poporange] at (axis cs:0.55,2.1) {$\mu$};
\node[popblue] at (axis cs:3.6,3.0) {$\mathcal{C}_f$};
\node[popmuted, font=\small] at (axis cs:2.5,1.0) {aire $=$ aire du rectangle};
\end{axis}
\end{tikzpicture}
\legende{La valeur moyenne $\mu$ : le rectangle de hauteur $\mu$ a la même aire que le domaine sous la courbe.}
\end{popfigure}

\textbf{Interprétation physique.}\par
Si $v(t)$ est la vitesse instantanée d'un véhicule entre les instants $a$ et $b$, alors $\frac{1}{b-a}\int_a^b v(t)\,\mathrm{d}t$ est sa \cle{vitesse moyenne} : distance parcourue divisée par la durée. De même, le débit moyen d'une rivière sur une journée est la valeur moyenne du débit instantané.

\begin{exemplebox}[{Valeur moyenne du sinus sur $[0\,;\pi]$}]
On admet provisoirement (le calcul sera refait rigoureusement après l'étude des primitives) que $\int_0^\pi \sin x\,\mathrm{d}x = 2$. La valeur moyenne de $\sin$ sur $[0\,;\pi]$ est donc :
\[ \mu = \frac{1}{\pi - 0}\int_0^\pi \sin x\,\mathrm{d}x = \frac{2}{\pi} \approx 0{,}64. \]
C'est cohérent : $\sin$ varie entre $0$ et $1$ sur $[0\,;\pi]$, et l'inégalité de la moyenne garantit bien $0 \leq \mu \leq 1$.
\end{exemplebox}

\begin{exoresolu}[Encadrer une intégrale puis estimer une valeur moyenne]
\textbf{Énoncé.} Soit $f(x) = \dfrac{1}{1+x^2}$ sur $[0\,;1]$.
\begin{enumerate}
\item Montrer que pour tout $x \in [0\,;1]$, $\dfrac{1}{2} \leq f(x) \leq 1$.
\item En déduire un encadrement de $I = \int_0^1 f(x)\,\mathrm{d}x$ puis de la valeur moyenne $\mu$ de $f$ sur $[0\,;1]$.
\end{enumerate}
\tcblower
\textbf{Solution.}
\begin{enumerate}
\item Pour $x \in [0\,;1]$, on a $0 \leq x^2 \leq 1$, donc $1 \leq 1 + x^2 \leq 2$. La fonction inverse étant décroissante sur $]0\,;+\infty[$ : $\dfrac{1}{2} \leq \dfrac{1}{1+x^2} \leq 1$.
\item On intègre cet encadrement sur $[0\,;1]$ (inégalité de la moyenne avec $m = \frac12$, $M = 1$, $b - a = 1$) :
\[ \frac{1}{2}(1-0) \leq I \leq 1(1-0) \quad \text{soit} \quad \frac{1}{2} \leq I \leq 1. \]
Comme $\mu = \dfrac{I}{1-0} = I$, on a aussi $\dfrac{1}{2} \leq \mu \leq 1$. (En fait, nous verrons que $I = \frac{\pi}{4} \approx 0{,}785$, bien dans l'encadrement.)
\end{enumerate}
\end{exoresolu}

\courssub{Intégrale d'une fonction valeur absolue}

\textbf{Problème.}\par
On veut calculer $\int_a^b |f(x)|\,\mathrm{d}x$. Géométriquement, c'est l'\cle{aire géométrique totale} entre la courbe de $f$ et l'axe des abscisses, comptée \emph{positivement} partout, même là où $f$ est négative. L'idée est simple : là où $f \geq 0$, on a $|f| = f$ ; là où $f \leq 0$, on a $|f| = -f$. On découpe donc l'intégrale avec la relation de Chasles.

\begin{methode}[Calculer $\int_a^b |f(x)|\,\mathrm{d}x$]
\begin{enumerate}
\item Résoudre l'équation $f(x) = 0$ sur $[a\,;b]$ et dresser le \textbf{tableau de signes} de $f$.
\item \textbf{Découper} l'intervalle $[a\,;b]$ aux points où $f$ change de signe, grâce à la relation de Chasles.
\item Sur chaque morceau, remplacer $|f|$ par $f$ (si $f \geq 0$) ou par $-f$ (si $f \leq 0$).
\item Calculer chaque intégrale et additionner les résultats.
\end{enumerate}
\end{methode}

\begin{exoresolu}[Calcul de $\int_{-1}^{2} |x|\,\mathrm{d}x$]
\textbf{Énoncé.} Calculer $J = \int_{-1}^{2} |x|\,\mathrm{d}x$ et interpréter géométriquement le résultat.
\tcblower
\textbf{Solution.} \textbf{Étape 1 :} $x = 0 \iff x = 0$ ; $x \leq 0$ sur $[-1\,;0]$ et $x \geq 0$ sur $[0\,;2]$.

\textbf{Étapes 2 et 3 :} par la relation de Chasles, en remplaçant $|x|$ par $-x$ sur $[-1\,;0]$ et par $x$ sur $[0\,;2]$ :
\[ J = \int_{-1}^{0} (-x)\,\mathrm{d}x + \int_{0}^{2} x\,\mathrm{d}x. \]
\textbf{Étape 4 :} chaque intégrale est l'aire d'un triangle :
\begin{align*}
\int_{-1}^{0} (-x)\,\mathrm{d}x &= \text{aire du triangle de base } 1 \text{ et de hauteur } 1 = \frac{1 \times 1}{2} = \frac{1}{2}, \\
\int_{0}^{2} x\,\mathrm{d}x &= \text{aire du triangle de base } 2 \text{ et de hauteur } 2 = \frac{2 \times 2}{2} = 2.
\end{align*}
Donc $J = \dfrac{1}{2} + 2 = \dfrac{5}{2}$. C'est l'aire totale entre la courbe de $f(x) = x$ et l'axe des abscisses sur $[-1\,;2]$.
\end{exoresolu}

\begin{attention}[Piège classique : $\int_a^b |f| \neq \left|\int_a^b f\right|$]
Beaucoup d'élèves écrivent $\int_a^b |f(x)|\,\mathrm{d}x = \left|\int_a^b f(x)\,\mathrm{d}x\right|$. C'est \textbf{faux en général} ! Contre-exemple avec $f(x) = x$ sur $[-1\,;2]$ :
\[ \int_{-1}^{2} x\,\mathrm{d}x = -\frac{1}{2} + 2 = \frac{3}{2} \quad \text{donc} \quad \left|\int_{-1}^{2} x\,\mathrm{d}x\right| = \frac{3}{2}, \]
alors que nous avons calculé $\int_{-1}^{2} |x|\,\mathrm{d}x = \dfrac{5}{2}$. Et $\dfrac{5}{2} \neq \dfrac{3}{2}$ !

Explication : dans $\int_a^b f$, les aires négatives \emph{compensent} les aires positives ; dans $\int_a^b |f|$, tout compte positivement. On a seulement l'inégalité $\left|\int_a^b f\right| \leq \int_a^b |f|$.
\end{attention}

\begin{popfigure}
\begin{tikzpicture}
\begin{axis}[width=0.85\linewidth, height=6cm, axis lines=middle, xlabel={$x$}, ylabel={$y$}, xmin=-1.6, xmax=2.6, ymin=-1.6, ymax=2.6, xtick={-1,0,2}, ytick={-1,1,2}, domain=-1.4:2.4]
\addplot[popcurve, samples=100, domain=-1.4:0] {x};
\addplot[popcurve, samples=100, domain=0:2.4] {x};
\addplot[poporange, thick, dashed, samples=100, domain=-1.4:0] {-x};
\addplot[popgreenL, draw=none, domain=-1:0, samples=50] {-x} \closedcycle;
\addplot[popblueL, draw=none, domain=0:2, samples=50] {x} \closedcycle;
\node[popblue] at (axis cs:1.9,1.5) {$y=|f(x)|$};
\node[poporange] at (axis cs:-1.15,0.6) {replié};
\node[popmuted, font=\small] at (axis cs:-0.5,-1.2) {$y=f(x)$};
\end{axis}
\end{tikzpicture}
\legende{La courbe de $|f|$ s'obtient en « repliant » au-dessus de l'axe les parties où $f$ est négative (pointillés orange). L'aire totale coloriée vaut $\int_{-1}^{2}|x|\,\mathrm{d}x = \frac{5}{2}$.}
\end{popfigure}

\begin{savaistu}[Pourquoi cette fonction $F$ est-elle si importante ?]
La fonction $F(x) = \int_a^x f(t)\,\mathrm{d}t$ est la clé de voûte de tout le calcul intégral : le \emph{théorème fondamental de l'analyse} (que nous rencontrerons dans la prochaine section) affirme que $F$ est dérivable et que $F' = f$. Autrement dit, intégrer et dériver sont des opérations \emph{réciproques} l'une de l'autre — découverte faite indépendamment par Newton et Leibniz au XVII\textsuperscript{e} siècle. Le résultat de cette section en est l'avant-goût : le signe de $f$ donne les variations de $F$, exactement comme le signe d'une dérivée donne les variations d'une fonction !
\end{savaistu}

\begin{exercice}[Pour t'entraîner]
Soit $f$ la fonction définie sur $[-2\,;3]$ par $f(x) = x^2 - 1$, et $F(x) = \int_{-2}^x f(t)\,\mathrm{d}t$.
\begin{enumerate}
\item Étudier le signe de $f$ sur $[-2\,;3]$ et en déduire le tableau de variations de $F$.
\item Montrer que $\dfrac{1}{3} \leq \int_2^3 \dfrac{1}{x}\,\mathrm{d}x \leq \dfrac{1}{2}$ en utilisant l'inégalité de la moyenne.
\item Calculer $\int_{-2}^{3} |x^2 - 1|\,\mathrm{d}x$ en admettant qu'une primitive de $x^2 - 1$ est $\dfrac{x^3}{3} - x$.
\end{enumerate}
\end{exercice}

\begin{corrige}[Exercice]
\begin{enumerate}
\item $f(x) = x^2 - 1 = (x-1)(x+1)$ s'annule en $-1$ et $1$. Signe : $f \geq 0$ sur $[-2\,;-1]$, $f \leq 0$ sur $[-1\,;1]$, $f \geq 0$ sur $[1\,;3]$. Donc $F$ est croissante sur $[-2\,;-1]$, décroissante sur $[-1\,;1]$, croissante sur $[1\,;3]$, avec $F(-2) = 0$. $F$ admet un maximum local en $-1$ et un minimum local en $1$.
\item Sur $[2\,;3]$, la fonction inverse est décroissante : $\dfrac{1}{3} \leq \dfrac{1}{x} \leq \dfrac{1}{2}$. Par l'inégalité de la moyenne avec $b - a = 1$ : $\dfrac{1}{3} \leq \int_2^3 \dfrac{\mathrm{d}x}{x} \leq \dfrac{1}{2}$.
\item D'après le tableau de signes de la question 1, on découpe par Chasles :
\[ \int_{-2}^{3} |x^2-1|\,\mathrm{d}x = \int_{-2}^{-1} (x^2-1)\,\mathrm{d}x - \int_{-1}^{1} (x^2-1)\,\mathrm{d}x + \int_{1}^{3} (x^2-1)\,\mathrm{d}x. \]
Avec $G(x) = \dfrac{x^3}{3} - x$ : $\int_{-2}^{-1} = G(-1)-G(-2) = \dfrac{2}{3} - \left(-\dfrac{2}{3}\right) = \dfrac{4}{3}$ ; $\int_{-1}^{1} = G(1)-G(-1) = -\dfrac{2}{3} - \dfrac{2}{3} = -\dfrac{4}{3}$ ; $\int_{1}^{3} = G(3)-G(1) = 6 - \left(-\dfrac{2}{3}\right) = \dfrac{20}{3}$. Total : $\dfrac{4}{3} + \dfrac{4}{3} + \dfrac{20}{3} = \dfrac{28}{3}$.
\end{enumerate}
\end{corrige}

\coursec{Techniques de calcul : primitives, intégration par parties, changement de variable affine}

Dans la section précédente, nous avons découvert la fonction $F : x \mapsto \displaystyle\int_a^x f(t)\,\mathrm{d}t$ et admis un fait capital : cette fonction est dérivable et $F' = f$. Autrement dit, \emph{intégrer} et \emph{dériver} sont deux opérations réciproques l'une de l'autre. C'est exactement ce pont entre dérivation et intégration qui va nous fournir, dans cette section, des outils de calcul efficaces. Jusqu'ici, nous savions estimer une intégrale par des encadrements ; nous allons maintenant apprendre à la \cle{calculer exactement}, grâce à trois techniques : la lecture directe d'une primitive, l'intégration par parties et le changement de variable affine. Ce sont ces trois gestes que le correcteur du Baccalauréat attend de toi, avec une rédaction soignée.

\courssub{Le lien fondamental : intégrale et primitives}

\courssub{Primitives : rappels et tableau de référence}

\begin{definition}[Primitive d'une fonction continue]
Soit $f$ une fonction continue sur un intervalle $I$. On appelle \cle{primitive} de $f$ sur $I$ toute fonction $F$ dérivable sur $I$ telle que $F' = f$.
\end{definition}

Deux primitives d'une même fonction sur un intervalle diffèrent d'une constante : si $F$ et $G$ sont deux primitives de $f$, alors $G(x) = F(x) + k$ pour une constante réelle $k$. Chercher une primitive, c'est donc \og remonter la dérivation \fg{} : on lit le tableau des dérivées \emph{à l'envers}.

\begin{aretenir}[Tableau des primitives usuelles]
\begin{center}
\begin{tabularx}{\linewidth}{|l|X|l|}
\hline
\rowcolor{popblueL} \textbf{Fonction $f$} & \textbf{Une primitive $F$} & \textbf{Intervalle} \\
\hline
$x^n$ ($n \in \mathbb{Z}$, $n \neq -1$) & $\dfrac{x^{n+1}}{n+1}$ & $\mathbb{R}$ si $n \geq 0$ ; $\mathbb{R}^*$ si $n \leq -2$ \\
\hline
$\dfrac{1}{x}$ & $\ln|x|$ & $]0\,;+\infty[$ ou $]-\infty\,;0[$ \\
\hline
$\dfrac{1}{\sqrt{x}}$ & $2\sqrt{x}$ & $]0\,;+\infty[$ \\
\hline
$\mathrm{e}^x$ & $\mathrm{e}^x$ & $\mathbb{R}$ \\
\hline
$\cos x$ & $\sin x$ & $\mathbb{R}$ \\
\hline
$\sin x$ & $-\cos x$ & $\mathbb{R}$ \\
\hline
$\ln x$ & $x\ln x - x$ & $]0\,;+\infty[$ \\
\hline
\end{tabularx}
\end{center}
\end{aretenir}

\begin{propriete}[Formes composées usuelles]
Soit $u$ une fonction dérivable. On reconnaîtra les formes suivantes (à connaître par c\oe ur) :
\begin{center}
\begin{tabularx}{\linewidth}{|X|X|}
\hline
\rowcolor{popblueL} \textbf{Forme $f$} & \textbf{Une primitive $F$} \\
\hline
$u' u^n$ ($n \neq -1$) & $\dfrac{u^{n+1}}{n+1}$ \\
\hline
$\dfrac{u'}{u}$ (avec $u > 0$) & $\ln u$ \\
\hline
$u' \mathrm{e}^{u}$ & $\mathrm{e}^{u}$ \\
\hline
$\dfrac{u'}{u^2}$ (avec $u \neq 0$) & $-\dfrac{1}{u}$ \\
\hline
$u' \sqrt{u}$ (avec $u > 0$) & $\dfrac{2}{3}u\sqrt{u}$ \\
\hline
\end{tabularx}
\end{center}
\end{propriete}

\begin{attention}[Le piège de l'exponentielle composée]
Une primitive de $\mathrm{e}^{2x}$ n'est \textbf{pas} $\mathrm{e}^{2x}$ ! En dérivant $\mathrm{e}^{2x}$, on obtient $2\mathrm{e}^{2x}$ : il y a un facteur $2$ en trop. La bonne primitive est $\dfrac{1}{2}\mathrm{e}^{2x}$. De manière générale, une primitive de $\mathrm{e}^{ax+b}$ (avec $a \neq 0$) est $\dfrac{1}{a}\mathrm{e}^{ax+b}$. Vérifie toujours ta primitive en la dérivant mentalement : c'est un réflexe qui ne coûte rien et qui sauve des points.
\end{attention}

\courssub{Le théorème fondamental du calcul intégral}

Voici le résultat central de toute cette leçon. Nous l'admettons, mais il est la conséquence directe de ce que nous avons vu dans la section 3 : la fonction $x \mapsto \displaystyle\int_a^x f(t)\,\mathrm{d}t$ est une primitive de $f$.

\begin{propriete}[Théorème fondamental (admis)]
Soit $f$ une fonction continue sur $[a\,;b]$ et $F$ une primitive \emph{quelconque} de $f$ sur $[a\,;b]$. Alors :
\[ \int_a^b f(x)\,\mathrm{d}x = F(b) - F(a). \]
On note cette différence $\big[F(x)\big]_a^b = F(b) - F(a)$.
\end{propriete}

Pourquoi le choix de la primitive n'a-t-il aucune importance ? Si $G$ est une autre primitive, alors $G = F + k$, donc $G(b) - G(a) = F(b) + k - F(a) - k = F(b) - F(a)$ : la constante $k$ s'élimine toujours. C'est pourquoi on écrit \og une \fg{} primitive et non \og la \fg{} primitive.

\begin{exemplebox}[Premier calcul par lecture directe]
Calculons $I = \displaystyle\int_1^3 (2x + 1)\,\mathrm{d}x$. Une primitive de $x \mapsto 2x+1$ est $F(x) = x^2 + x$. Donc :
\[ I = \big[x^2 + x\big]_1^3 = (9 + 3) - (1 + 1) = 12 - 2 = 10. \]
Géométriquement, c'est l'aire du trapèze sous la droite d'équation $y = 2x+1$ entre $x=1$ et $x=3$ : on retrouve $\dfrac{(3 + 7)\times 2}{2} = 10$. Le calcul intégral est cohérent avec la géométrie !
\end{exemplebox}

\begin{methode}[Calculer une intégrale par lecture directe d'une primitive]
\begin{enumerate}
\item Identifier la \textbf{forme} de $f$ : puissance, $u'/u$, $u'\mathrm{e}^u$, etc.
\item Écrire une primitive candidate $F$, en \textbf{ajustant le coefficient constant} si nécessaire (par exemple le $\tfrac{1}{a}$ devant $\mathrm{e}^{ax+b}$).
\item \textbf{Vérifier} en dérivant $F$ : on doit retomber exactement sur $f$.
\item Appliquer $\displaystyle\int_a^b f = F(b) - F(a)$ et simplifier.
\end{enumerate}
\end{methode}

\begin{exoresolu}[Lecture directe avec ajustement de coefficient]
Énoncé. Calculer $J = \displaystyle\int_0^1 \dfrac{2x}{x^2+1}\,\mathrm{d}x$.
\tcblower
Solution détaillée. On reconnaît la forme $\dfrac{u'}{u}$ avec $u(x) = x^2 + 1$, car $u'(x) = 2x$ : le numérateur est exactement la dérivée du dénominateur, et $u(x) > 0$ sur $[0\,;1]$. Une primitive est donc $F(x) = \ln(x^2+1)$. Vérification : $F'(x) = \dfrac{2x}{x^2+1}$. C'est correct. Alors :
\[ J = \big[\ln(x^2+1)\big]_0^1 = \ln 2 - \ln 1 = \ln 2. \]
Conclusion : $J = \ln 2 \approx 0{,}69$.
\end{exoresolu}

\courssub{Intégration par parties}

Certaines fonctions, comme $x\mathrm{e}^x$ ou $\ln x$, n'apparaissent pas dans notre tableau et ne ressemblent à aucune forme composée. Il faut alors un outil qui transforme l'intégrale d'un produit en une intégrale plus simple : c'est l'\cle{intégration par parties} (IPP).

\begin{propriete}[Formule d'intégration par parties (admise)]
Soient $u$ et $v$ deux fonctions dérivables sur $[a\,;b]$, de dérivées $u'$ et $v'$ continues sur $[a\,;b]$. Alors :
\[ \int_a^b u(x)\,v'(x)\,\mathrm{d}x = \big[u(x)v(x)\big]_a^b - \int_a^b u'(x)\,v(x)\,\mathrm{d}x. \]
\end{propriete}

D'où vient cette formule ? Elle n'est que la version intégrée de la dérivée d'un produit : $(uv)' = u'v + uv'$. En intégrant cette égalité entre $a$ et $b$, on obtient $\big[uv\big]_a^b = \displaystyle\int_a^b u'v + \int_a^b uv'$, puis on isole $\displaystyle\int_a^b uv'$. La démonstration guidée est proposée en exercice en fin de section : elle est formatrice et tout à fait à ta portée.

\begin{methode}[Réussir une intégration par parties]
Au Baccalauréat, l'énoncé \textbf{fournit toujours} le choix de $u$ et $v'$ (conformément au programme). Ta démarche :
\begin{enumerate}
\item Poser clairement $u(x) = \ldots$ et $v'(x) = \ldots$ comme indiqué.
\item Calculer $u'(x)$ (dériver $u$) et $v(x)$ (primitiver $v'$).
\item Écrire la formule : $\displaystyle\int_a^b uv' = \big[uv\big]_a^b - \int_a^b u'v$.
\item Vérifier que la nouvelle intégrale $\displaystyle\int_a^b u'v$ est \textbf{plus simple} que celle de départ ; sinon, le choix de $u$ et $v'$ est probablement inadapté.
\item Conclure par le calcul exact, puis numérique si demandé.
\end{enumerate}
Les cas types au Bac : $x\mathrm{e}^x$ et $x\cos x$ (on dérive le $x$ pour le faire disparaître) ; $\ln x$ (on pose $u = \ln x$ et $v' = 1$).
\end{methode}

\begin{exemplebox}[Le calcul de référence : $\displaystyle\int_0^1 x\mathrm{e}^x\,\mathrm{d}x = 1$]
Posons, comme le suggérerait l'énoncé, $u(x) = x$ et $v'(x) = \mathrm{e}^x$. Alors $u'(x) = 1$ et $v(x) = \mathrm{e}^x$. Les dérivées sont continues, la formule s'applique :
\begin{align*}
\int_0^1 x\mathrm{e}^x\,\mathrm{d}x &= \big[x\mathrm{e}^x\big]_0^1 - \int_0^1 1 \times \mathrm{e}^x\,\mathrm{d}x \\
&= \big(1\cdot\mathrm{e}^1 - 0\big) - \big[\mathrm{e}^x\big]_0^1 \\
&= \mathrm{e} - (\mathrm{e} - 1) = 1.
\end{align*}
Remarque l'efficacité de la méthode : dériver $x$ l'a fait disparaître, et l'intégrale restante $\displaystyle\int_0^1 \mathrm{e}^x\,\mathrm{d}x$ était immédiate.
\end{exemplebox}

\begin{exoresolu}[Intégrer le logarithme népérien]
Énoncé. En posant $u(x) = \ln x$ et $v'(x) = 1$, calculer $K = \displaystyle\int_1^{\mathrm{e}} \ln x\,\mathrm{d}x$.
\tcblower
Solution détaillée. On a $u(x) = \ln x$, donc $u'(x) = \dfrac{1}{x}$ ; et $v'(x) = 1$, donc $v(x) = x$. Les fonctions $u'$ et $v'$ sont continues sur $[1\,;\mathrm{e}]$. La formule d'intégration par parties donne :
\begin{align*}
K &= \big[x\ln x\big]_1^{\mathrm{e}} - \int_1^{\mathrm{e}} \dfrac{1}{x}\times x\,\mathrm{d}x \\
&= \big(\mathrm{e}\ln\mathrm{e} - 1\times\ln 1\big) - \int_1^{\mathrm{e}} 1\,\mathrm{d}x \\
&= (\mathrm{e} - 0) - \big[x\big]_1^{\mathrm{e}} = \mathrm{e} - (\mathrm{e} - 1) = 1.
\end{align*}
Conclusion : $K = 1$. Au passage, ce calcul justifie la primitive usuelle de $\ln$ : $x\ln x - x$.
\end{exoresolu}

\begin{attention}[Pièges de rédaction en IPP]
\begin{itemize}
\item N'oublie pas le \textbf{signe moins} devant la deuxième intégrale : c'est l'erreur la plus fréquente.
\item Le crochet $\big[u(x)v(x)\big]_a^b$ se calcule \emph{avant} de traiter l'intégrale restante ; ne mélange pas les deux étapes.
\item Justifie brièvement les hypothèses ($u'$ et $v'$ continues) : une phrase suffit et montre ta rigueur.
\item Si après l'IPP la nouvelle intégrale est \emph{plus compliquée} (par exemple en dérivant $\mathrm{e}^x$ au lieu de $x$ dans $\int x\mathrm{e}^x$), c'est le signal qu'il faut échanger les rôles de $u$ et $v'$.
\end{itemize}
\end{attention}

\courssub{Changement de variable affine}

Dernière technique : lorsque la fonction à intégrer est de la forme $f(\alpha x + \beta)$ — une fonction composée avec une fonction \emph{affine} — on simplifie le problème en posant $t = \alpha x + \beta$. C'est le \cle{changement de variable affine}.

\begin{propriete}[Changement de variable affine]
Soit $f$ une fonction continue et $\alpha$, $\beta$ deux réels avec $\alpha \neq 0$. Alors :
\[ \int_a^b f(\alpha x + \beta)\,\mathrm{d}x = \frac{1}{\alpha}\int_{\alpha a + \beta}^{\alpha b + \beta} f(t)\,\mathrm{d}t. \]
En pratique : on pose $t = \alpha x + \beta$, donc $\mathrm{d}t = \alpha\,\mathrm{d}x$, c'est-à-dire $\mathrm{d}x = \dfrac{\mathrm{d}t}{\alpha}$, et on transforme les bornes : quand $x = a$, $t = \alpha a + \beta$ ; quand $x = b$, $t = \alpha b + \beta$.
\end{propriete}

\begin{popfigure}
\begin{center}
\begin{tikzpicture}[scale=1, >=Stealth]
% Ligne des x
\draw[->, thick, popdark] (0,0) -- (7,0) node[right] {$x$};
\fill[popblue] (1,0) circle (2pt) node[below=2pt] {$a$};
\fill[popblue] (5.5,0) circle (2pt) node[below=2pt] {$b$};
\draw[popblue, very thick] (1,0.12) -- (5.5,0.12);
\node[popblue, above] at (3.25,0.15) {on intègre en $x$ : $\mathrm{d}x$};
% Flèches de transformation
\draw[->, very thick, poporange, bend left=25] (1,0.55) to node[above, sloped, font=\small] {$t = \alpha x + \beta$} (1,2.45);
\draw[->, very thick, poporange, bend right=25] (5.5,0.55) to (5.5,2.45);
% Ligne des t
\draw[->, thick, popdark] (0,3) -- (7,3) node[right] {$t$};
\fill[popgreen] (1,3) circle (2pt) node[above=2pt] {$\alpha a + \beta$};
\fill[popgreen] (5.5,3) circle (2pt) node[above=2pt] {$\alpha b + \beta$};
\draw[popgreen, very thick] (1,2.88) -- (5.5,2.88);
\node[popgreen, below] at (3.25,2.85) {on intègre en $t$ : $\mathrm{d}x = \dfrac{\mathrm{d}t}{\alpha}$};
\end{tikzpicture}
\end{center}
\legende{Schéma du changement de variable affine : la variable $x$ parcourt $[a\,;b]$, la nouvelle variable $t = \alpha x + \beta$ parcourt $[\alpha a + \beta\,;\,\alpha b + \beta]$, et $\mathrm{d}x$ devient $\mathrm{d}t/\alpha$.}
\end{popfigure}

\begin{methode}[Appliquer un changement de variable affine]
\begin{enumerate}
\item Poser $t = \alpha x + \beta$ et écrire $\mathrm{d}t = \alpha\,\mathrm{d}x$, donc $\mathrm{d}x = \dfrac{\mathrm{d}t}{\alpha}$.
\item \textbf{Transformer les deux bornes} : nouvelle borne basse $= \alpha a + \beta$, nouvelle borne haute $= \alpha b + \beta$.
\item Réécrire l'intégrale entièrement en $t$, sans oublier le \textbf{facteur} $\dfrac{1}{\alpha}$.
\item Intégrer en $t$ (primitive usuelle) et conclure. Il est inutile de revenir à $x$ puisque les bornes ont été changées.
\end{enumerate}
\end{methode}

\begin{exemplebox}[Trois applications directes]
\textbf{a)} $\displaystyle\int_0^1 \mathrm{e}^{3x+1}\,\mathrm{d}x$. Posons $t = 3x+1$, $\mathrm{d}x = \dfrac{\mathrm{d}t}{3}$ ; bornes : $x=0 \Rightarrow t=1$ ; $x=1 \Rightarrow t=4$.
\[ \int_0^1 \mathrm{e}^{3x+1}\,\mathrm{d}x = \frac{1}{3}\int_1^4 \mathrm{e}^t\,\mathrm{d}t = \frac{1}{3}\big[\mathrm{e}^t\big]_1^4 = \frac{\mathrm{e}^4 - \mathrm{e}}{3}. \]
\textbf{b)} $\displaystyle\int_0^{\pi/2} \sin(2x)\,\mathrm{d}x$. Posons $t = 2x$, $\mathrm{d}x = \dfrac{\mathrm{d}t}{2}$ ; bornes : $0 \to 0$ et $\dfrac{\pi}{2} \to \pi$.
\[ \int_0^{\pi/2} \sin(2x)\,\mathrm{d}x = \frac{1}{2}\int_0^{\pi} \sin t\,\mathrm{d}t = \frac{1}{2}\big[-\cos t\big]_0^{\pi} = \frac{1}{2}\big(-\cos\pi + \cos 0\big) = \frac{1}{2}(1 + 1) = 1. \]
\textbf{c)} $\displaystyle\int_1^2 \dfrac{\mathrm{d}x}{2x+1}$. Posons $t = 2x+1$, $\mathrm{d}x = \dfrac{\mathrm{d}t}{2}$ ; bornes : $1 \to 3$ et $2 \to 5$.
\[ \int_1^2 \frac{\mathrm{d}x}{2x+1} = \frac{1}{2}\int_3^5 \frac{\mathrm{d}t}{t} = \frac{1}{2}\big[\ln t\big]_3^5 = \frac{1}{2}\ln\frac{5}{3}. \]
\end{exemplebox}

\begin{exoresolu}[Un calcul complet avec une racine carrée]
Énoncé. À l'aide du changement de variable $t = 2x+1$, calculer $L = \displaystyle\int_0^4 \sqrt{2x+1}\,\mathrm{d}x$.
\tcblower
Solution détaillée. Posons $t = 2x + 1$. Alors $\mathrm{d}t = 2\,\mathrm{d}x$, donc $\mathrm{d}x = \dfrac{\mathrm{d}t}{2}$.
Bornes : quand $x = 0$, $t = 1$ ; quand $x = 4$, $t = 9$. On remplace :
\[ L = \int_1^9 \sqrt{t}\times\frac{\mathrm{d}t}{2} = \frac{1}{2}\int_1^9 t^{1/2}\,\mathrm{d}t. \]
Une primitive de $t^{1/2}$ est $\dfrac{t^{3/2}}{3/2} = \dfrac{2}{3}t^{3/2} = \dfrac{2}{3}t\sqrt{t}$. Donc :
\begin{align*}
L &= \frac{1}{2}\left[\frac{2}{3}t\sqrt{t}\right]_1^9 = \frac{1}{3}\big[t\sqrt{t}\big]_1^9 \\
&= \frac{1}{3}\big(9\sqrt{9} - 1\sqrt{1}\big) = \frac{1}{3}(27 - 1) = \frac{26}{3}.
\end{align*}
Conclusion : $L = \dfrac{26}{3} \approx 8{,}67$. Retiens l'enchaînement : changement de variable $\to$ nouvelles bornes $\to$ facteur $\tfrac12$ $\to$ primitive de $\sqrt{t}$.
\end{exoresolu}

\begin{attention}[Les trois erreurs classiques du changement de variable]
\begin{itemize}
\item \textbf{Oublier de changer les bornes} : si tu intègres en $t$ avec les bornes en $x$, le résultat est faux. Les bornes appartiennent à la variable d'intégration.
\item \textbf{Oublier le facteur $\dfrac{1}{\alpha}$} : écrire $\displaystyle\int_0^1 \mathrm{e}^{3x+1}\mathrm{d}x = \int_1^4 \mathrm{e}^t\,\mathrm{d}t$ est faux ; il manque le $\tfrac13$.
\item \textbf{Mélanger les variables} : dans une même intégrale, on ne doit voir que $x$ (avant) ou que $t$ (après), jamais les deux.
\end{itemize}
\end{attention}

\courssub{Quelle technique choisir ? Organigramme de décision}

Face à une intégrale au Baccalauréat, pose-toi les questions dans cet ordre :

\begin{popfigure}
\begin{center}
\begin{tikzpicture}[
  node distance=0.9cm and 1.6cm,
  decision/.style={diamond, draw=poporange, fill=poporangeL, text width=3.4cm, align=center, inner sep=1pt, aspect=2.2, font=\small},
  action/.style={rectangle, rounded corners, draw=popblue, fill=popblueL, text width=3.6cm, align=center, font=\small},
  lab/.style={font=\small\itshape, popdark}
]
\node[decision] (q1) {$f$ est-elle dans le tableau ou de forme $u'u^n$, $\frac{u'}{u}$, $u'\mathrm{e}^u$ ?};
\node[action, right=of q1] (a1) {Primitive directe : $\int_a^b f = F(b)-F(a)$};
\node[decision, below=of q1] (q2) {$f$ est-elle un produit ($x\mathrm{e}^x$, $x\cos x$) ou $\ln x$ ?};
\node[action, right=of q2] (a2) {Intégration par parties ($u$, $v'$ fournis par l'énoncé)};
\node[decision, below=of q2] (q3) {$f$ est-elle de la forme $g(\alpha x + \beta)$ ?};
\node[action, right=of q3] (a3) {Changement de variable $t = \alpha x + \beta$};
\node[action, below=of q3] (a4) {Décomposer $f$ (linéarité) puis traiter chaque morceau};
\draw[->, thick, popdark] (q1) -- node[lab, above] {oui} (a1);
\draw[->, thick, popdark] (q1) -- node[lab, right] {non} (q2);
\draw[->, thick, popdark] (q2) -- node[lab, above] {oui} (a2);
\draw[->, thick, popdark] (q2) -- node[lab, right] {non} (q3);
\draw[->, thick, popdark] (q3) -- node[lab, above] {oui} (a3);
\draw[->, thick, popdark] (q3) -- node[lab, right] {non} (a4);
\end{tikzpicture}
\end{center}
\legende{Organigramme de décision : trois questions suffisent pour choisir la bonne technique de calcul.}
\end{popfigure}

\begin{savaistu}[Complément utile au Bac]
Même lorsque l'énoncé ne dit rien, reconnaître les formes $\dfrac{u'}{u}$ et $u'\mathrm{e}^u$ accélère considérablement les calculs. Par exemple, $\displaystyle\int_0^1 \frac{\mathrm{e}^x}{\mathrm{e}^x + 1}\,\mathrm{d}x = \big[\ln(\mathrm{e}^x+1)\big]_0^1 = \ln\frac{\mathrm{e}+1}{2}$, car le numérateur est la dérivée du dénominateur. Ce réflexe, combiné à la linéarité, permet de traiter la grande majorité des intégrales proposées en Terminale.
\end{savaistu}

\begin{exercice}[Démonstration guidée : la formule d'intégration par parties]
Soient $u$ et $v$ deux fonctions dérivables sur $[a\,;b]$, à dérivées continues.
\begin{enumerate}
\item Rappeler la formule de dérivation du produit $(uv)'$.
\item En intégrant cette égalité entre $a$ et $b$, montrer que $\big[u(x)v(x)\big]_a^b = \displaystyle\int_a^b u'(x)v(x)\,\mathrm{d}x + \int_a^b u(x)v'(x)\,\mathrm{d}x$.
\item En déduire la formule d'intégration par parties.
\end{enumerate}
\end{exercice}

\begin{corrige}[Exercice 1]
1. $(uv)' = u'v + uv'$, c'est-à-dire $\big(u(x)v(x)\big)' = u'(x)v(x) + u(x)v'(x)$.
2. La fonction $uv$ est une primitive de $(uv)'$. Par le théorème fondamental :
\[ \int_a^b (uv)'(x)\,\mathrm{d}x = \big[u(x)v(x)\big]_a^b. \]
Par linéarité de l'intégrale, le membre de gauche vaut $\displaystyle\int_a^b u'v + \int_a^b uv'$, d'où l'égalité demandée.
3. En isolant $\displaystyle\int_a^b u(x)v'(x)\,\mathrm{d}x$, on obtient :
\[ \int_a^b u(x)v'(x)\,\mathrm{d}x = \big[u(x)v(x)\big]_a^b - \int_a^b u'(x)v(x)\,\mathrm{d}x, \]
qui est exactement la formule d'intégration par parties.
\end{corrige}

\begin{exercice}[Entraînement au choix de technique]
Calculer les intégrales suivantes en précisant la technique utilisée :
\[ A = \int_1^2 \frac{3}{x}\,\mathrm{d}x \;;\qquad B = \int_0^{\pi} x\cos x\,\mathrm{d}x \;\;(\text{poser } u = x,\ v' = \cos x) \;;\qquad C = \int_0^1 \frac{\mathrm{d}x}{3x+2}. \]
\end{exercice}

\begin{corrige}[Exercice 2]
\textbf{Calcul de $A$} (primitive directe) : $A = 3\big[\ln x\big]_1^2 = 3(\ln 2 - \ln 1) = 3\ln 2$.

\textbf{Calcul de $B$} (intégration par parties) : posons $u(x) = x$ et $v'(x) = \cos x$, donc $u'(x) = 1$ et $v(x) = \sin x$. Les dérivées sont continues sur $[0\,;\pi]$. La formule donne :
\[ B = \big[x\sin x\big]_0^{\pi} - \int_0^{\pi} \sin x\,\mathrm{d}x. \]
Or $\big[x\sin x\big]_0^{\pi} = \pi\sin\pi - 0 = 0$ et $\displaystyle\int_0^{\pi}\sin x\,\mathrm{d}x = \big[-\cos x\big]_0^{\pi} = -\cos\pi + \cos 0 = 1 + 1 = 2$. Donc $B = 0 - 2 = -2$.

\textbf{Calcul de $C$} (changement de variable $t = 3x+2$) : on a $\mathrm{d}x = \dfrac{\mathrm{d}t}{3}$ ; bornes : $x=0 \Rightarrow t=2$ et $x=1 \Rightarrow t=5$.
\[ C = \frac{1}{3}\int_2^5 \frac{\mathrm{d}t}{t} = \frac{1}{3}\big[\ln t\big]_2^5 = \frac{1}{3}\ln\frac{5}{2}. \]
\end{corrige}

\begin{aretenir}[L'essentiel de la section]
\begin{itemize}
\item \textbf{Théorème fondamental} : si $F$ est une primitive de $f$ continue sur $[a\,;b]$, alors $\displaystyle\int_a^b f(x)\,\mathrm{d}x = \big[F(x)\big]_a^b = F(b) - F(a)$.
\item \textbf{Intégration par parties} : $\displaystyle\int_a^b uv' = \big[uv\big]_a^b - \int_a^b u'v$, avec $u$ et $v'$ fournis par l'énoncé.
\item \textbf{Changement de variable affine} : $\displaystyle\int_a^b f(\alpha x + \beta)\,\mathrm{d}x = \frac{1}{\alpha}\int_{\alpha a+\beta}^{\alpha b+\beta} f(t)\,\mathrm{d}t$ ; on change les bornes et on n'oublie pas le facteur $\dfrac{1}{\alpha}$.
\item Dans tous les cas : \textbf{vérifie ta primitive en la dérivant}.
\end{itemize}
\end{aretenir}

Tu disposes maintenant de la boîte à outils complète du calcul intégral. Dans la prochaine section, nous mettrons ces techniques au service de la géométrie : calculs d'aires entre courbes et de volumes de solides de révolution — des applications spectaculaires où chaque franc d'effort investi ici sera remboursé avec intérêts.

\coursec{Applications géométriques : aires et volumes de révolution}

Dans la section précédente, tu as appris à \emph{calculer} des intégrales : lecture directe d'une primitive, intégration par parties, changement de variable affine. Il est temps de récolter les fruits de ce travail. Car l'intégrale n'est pas qu'un exercice de calcul : c'est d'abord un \cle{outil de mesure}. Mesurer l'aire d'un terrain aux contours courbes, le volume d'un silo à maïs, la capacité d'un réservoir d'eau : voilà ce que l'intégrale permet de faire avec une précision parfaite. Cette section te donne les deux grandes applications géométriques exigibles au Baccalauréat : les \cle{aires} et les \cle{volumes de révolution}.

\courssub{Aire entre une courbe et l'axe des abscisses}

\textbf{L'unité d'aire : le cadre du repère.}\par
Plaçons-nous dans un repère orthogonal $(O\,;\vec{\imath},\vec{\jmath})$. L'\cle{unité d'aire} (notée u.a.) est l'aire du rectangle construit sur les vecteurs unitaires : $1\ \text{u.a.} = \|\vec{\imath}\| \times \|\vec{\jmath}\|$. Si le repère est orthonormé avec des graduations en centimètres, alors $1\ \text{u.a.} = 1\ \text{cm}^2$. Mais si l'axe des abscisses est gradué avec $2$ cm par unité et l'axe des ordonnées avec $3$ cm par unité, alors $1\ \text{u.a.} = 2 \times 3 = 6\ \text{cm}^2$ : il faut le recalculer à chaque énoncé !

\begin{definition}[Aire sous une courbe positive]
Soit $f$ une fonction \cle{continue} et \cle{positive} sur un intervalle $[a\,;b]$ ($a < b$), de courbe représentative $\mathcal{C}_f$ dans un repère orthogonal. L'\cle{aire} du domaine délimité par $\mathcal{C}_f$, l'axe des abscisses et les droites d'équations $x=a$ et $x=b$ est :
\[ \mathcal{A} = \int_a^b f(x)\,\mathrm{d}x \quad \text{(en unités d'aire).} \]
\end{definition}

Cette définition est cohérente avec tout ce que nous avons construit : pour une fonction positive, l'intégrale \emph{est} l'aire. Mais que faire si la courbe passe \emph{sous} l'axe des abscisses ?

\begin{propriete}[Aire et signe de la fonction]
Soit $f$ continue sur $[a\,;b]$.
\begin{itemize}
\item Si $f \geq 0$ sur $[a\,;b]$, alors $\mathcal{A} = \displaystyle\int_a^b f(x)\,\mathrm{d}x$.
\item Si $f \leq 0$ sur $[a\,;b]$, alors $\mathcal{A} = -\displaystyle\int_a^b f(x)\,\mathrm{d}x$.
\item Si $f$ change de signe, on \cle{découpe} l'intervalle aux points où $f$ s'annule (relation de Chasles) et on additionne les aires des morceaux, ce qui revient à calculer :
\[ \mathcal{A} = \int_a^b \bigl|f(x)\bigr|\,\mathrm{d}x. \]
\end{itemize}
Dans tous les cas, une aire est un nombre \textbf{positif}.
\end{propriete}

\begin{attention}[Aire $\neq$ intégrale]
L'intégrale $\int_a^b f$ est un nombre \emph{algébrique} : elle peut être négative ou nulle alors que le domaine a une aire non nulle. Par exemple, $\int_{-1}^{1} x\,\mathrm{d}x = 0$, mais l'aire entre la droite $y=x$ et l'axe sur $[-1\,;1]$ vaut $1$ u.a. (deux triangles de surface $\tfrac12$). Au Bac, si l'on te demande une \emph{aire}, vérifie d'abord le \cle{signe} de $f$ et découpe si nécessaire.
\end{attention}

\begin{exemplebox}[Aire avec changement de signe]
Soit $f(x) = x^2 - 1$ sur $[0\,;2]$. La fonction s'annule en $x=1$ : elle est négative sur $[0\,;1]$ et positive sur $[1\,;2]$. Une primitive est $F(x) = \dfrac{x^3}{3} - x$. L'aire totale est :
\begin{align*}
\mathcal{A} &= -\int_0^1 (x^2-1)\,\mathrm{d}x + \int_1^2 (x^2-1)\,\mathrm{d}x \\
&= -\bigl[F(1)-F(0)\bigr] + \bigl[F(2)-F(1)\bigr] \\
&= -\left(\tfrac{1}{3}-1\right) + \left(\tfrac{8}{3}-2-\tfrac{1}{3}+1\right) = \tfrac{2}{3} + \tfrac{4}{3} = 2\ \text{u.a.}
\end{align*}
Si le repère est orthonormé avec $1$ cm par unité, $\mathcal{A} = 2\ \text{cm}^2$.
\end{exemplebox}

\begin{propriete}[Conversion en unités concrètes]
Si une unité sur l'axe $(Ox)$ correspond à une longueur $p$ (par exemple $p = 2$ cm) et une unité sur $(Oy)$ à une longueur $q$, alors :
\[ \mathcal{A}\ (\text{en cm}^2) = \mathcal{A}\ (\text{en u.a.}) \times p \times q. \]
\end{propriete}

\courssub{Aire entre deux courbes}

\textbf{Intuition.}\par
Imaginons deux routes vues du ciel : l'une (la courbe de $f$) passe toujours au-dessus de l'autre (la courbe de $g$) entre les kilomètres $a$ et $b$. La surface de la « bande » comprise entre les deux routes, c'est l'aire sous $f$ \emph{moins} l'aire sous $g$. D'où la propriété suivante.

\begin{propriete}[Aire entre deux courbes]
Soient $f$ et $g$ deux fonctions continues sur $[a\,;b]$ telles que $f(x) \geq g(x)$ pour tout $x \in [a\,;b]$. L'aire du domaine délimité par les courbes $\mathcal{C}_f$, $\mathcal{C}_g$ et les droites $x=a$, $x=b$ est :
\[ \mathcal{A} = \int_a^b \bigl(f(x)-g(x)\bigr)\,\mathrm{d}x \quad \text{(u.a.)} \]
\end{propriete}

\begin{experience}[Démonstration (exigible)]
\textbf{Cas où $f \geq g \geq 0$.} L'aire sous $\mathcal{C}_f$ vaut $\int_a^b f$ et celle sous $\mathcal{C}_g$ vaut $\int_a^b g$. L'aire entre les courbes est la différence : $\int_a^b f - \int_a^b g = \int_a^b (f-g)$ par linéarité.

\textbf{Cas général.} Les fonctions $f$ et $g$ sont continues sur $[a\,;b]$, donc $g$ admet un minimum $m$ sur $[a\,;b]$. Posons $k = |m| + 1$ et translatons tout verticalement : $\tilde{f} = f + k$ et $\tilde{g} = g + k$ sont alors strictement positives, avec $\tilde{f} \geq \tilde{g}$. La translation ne change pas l'aire du domaine entre les courbes, donc :
\[ \mathcal{A} = \int_a^b \bigl(\tilde{f}-\tilde{g}\bigr) = \int_a^b \bigl(f-g\bigr). \]
La formule est donc valable quelles que soient les positions des courbes par rapport à l'axe.
\end{experience}

\begin{popfigure}
\begin{center}
\begin{tikzpicture}[scale=2.2]
% Functions: f(x) = -(x-1)^2 + 2, g(x) = 1. Intersections at x=0 and x=2. f >= g on [0,2].
\fill[popblueL] plot[domain=0:2,samples=60] (\x,{ -(\x-1)*(\x-1) + 2 })
 -- plot[domain=2:0,samples=60] (\x,{1}) -- cycle;
\draw[->,popmuted] (-0.4,0) -- (2.8,0) node[below left] {$x$};
\draw[->,popmuted] (0,-0.4) -- (0,2.6) node[below left] {$y$};
\draw[popblue,very thick] plot[domain=-0.3:2.3,samples=60] (\x,{ -(\x-1)*(\x-1) + 2 }) node[right] {$\mathcal{C}_f$};
\draw[poporange,very thick] plot[domain=-0.3:2.3,samples=60] (\x,{1}) node[right] {$\mathcal{C}_g$};
\draw[dashed,popmuted] (0,0) node[below] {$a$} -- (0,2);
\draw[dashed,popmuted] (2,0) node[below] {$b$} -- (2,1);
\filldraw[popdark] (0,1) circle (0.03) (0,2) circle (0.03) (2,1) circle (0.03) (2,2) circle (0.03);
\node[popdark] at (1,1.5) {$\mathcal{A}$};
\end{tikzpicture}
\end{center}
\legende{Les courbes de $f$ et $g$ se coupent en $x=a=0$ et $x=b=2$ ; sur $[0\,;2]$, $f \geq g$, et l'aire hachurée vaut $\mathcal{A} = \displaystyle\int_a^b \bigl(f(x)-g(x)\bigr)\,\mathrm{d}x$.}
\end{popfigure}

\begin{methode}[Calcul d'aire entre deux courbes — démarche type Bac]
\begin{enumerate}
\item \textbf{Chercher les intersections} : résoudre $f(x) = g(x)$. Les solutions fournissent (souvent) les bornes $a$ et $b$.
\item \textbf{Étudier la position relative} : étudier le signe de $f(x) - g(x)$ sur $[a\,;b]$ (tableau de signes) pour savoir quelle courbe est au-dessus. Si le signe change, découper avec Chasles.
\item \textbf{Intégrer l'écart} : calculer $\int_a^b (f-g)$ avec une primitive de $f-g$.
\item \textbf{Convertir} : multiplier par $p \times q$ (unités des axes) si l'énoncé demande des cm$^2$.
\end{enumerate}
\end{methode}

\begin{exoresolu}[Aire entre une parabole et une droite]
\textbf{Énoncé.} Dans un repère orthonormé (unité : $1$ cm), on donne $f(x) = x^2$ et $g(x) = x + 2$. Calculer l'aire du domaine compris entre $\mathcal{C}_f$ et $\mathcal{C}_g$.
\tcblower
\textbf{Solution.} \textbf{1) Intersections :} $x^2 = x+2 \iff x^2 - x - 2 = 0 \iff (x-2)(x+1)=0$, donc $x=-1$ ou $x=2$.

\textbf{2) Position relative :} sur $[-1\,;2]$, le trinôme $x^2-x-2$ est négatif (signe contraire à celui de $x^2$ entre les racines), donc $f(x) \leq g(x)$ : la droite est au-dessus.

\textbf{3) Intégration :}
\begin{align*}
\mathcal{A} &= \int_{-1}^{2} \bigl(g(x)-f(x)\bigr)\,\mathrm{d}x = \int_{-1}^{2} \left(-x^2+x+2\right)\mathrm{d}x \\
&= \left[-\dfrac{x^3}{3}+\dfrac{x^2}{2}+2x\right]_{-1}^{2} = \left(-\tfrac{8}{3}+2+4\right) - \left(\tfrac{1}{3}+\tfrac{1}{2}-2\right) = \dfrac{9}{2}\ \text{u.a.}
\end{align*}
\textbf{4) Conversion :} le repère est orthonormé d'unité $1$ cm, donc $\mathcal{A} = 4,5\ \text{cm}^2$.
\end{exoresolu}

\courssub{Volume d'un solide de révolution}

\textbf{Intuition : la tranche de saucisson.}\par
Faisons tourner la courbe d'une fonction positive $f$ autour de l'axe $(Ox)$ : elle engendre un solide (pense à un vase sur un tour de potier). Pour mesurer son volume, découpons-le en tranches très fines, perpendiculairement à $(Ox)$. Chaque tranche, située à l'abscisse $x$ et d'épaisseur $\mathrm{d}x$, est quasiment un \cle{cylindre} de rayon $f(x)$ : son volume est $\pi [f(x)]^2\,\mathrm{d}x$. En « sommant » toutes ces tranches de $a$ à $b$, on obtient le théorème fondamental.

\begin{propriete}[Volume de révolution autour de $(Ox)$ — admis]
Soit $f$ une fonction continue et positive sur $[a\,;b]$. Le solide engendré par la rotation autour de l'axe $(Ox)$ du domaine situé sous la courbe de $f$ a pour volume :
\[ V = \pi \int_a^b \bigl[f(x)\bigr]^2\,\mathrm{d}x \quad \text{(en unités de volume).} \]
Si les deux axes ont la même unité de longueur (par exemple $1$ cm), alors $1\ \text{u.v.} = (1\ \text{cm})^3 = 1\ \text{cm}^3$.
\end{propriete}

\begin{popfigure}
\begin{center}
\begin{tikzpicture}[scale=1.6]
% axe
\draw[->,popmuted] (-0.3,0) -- (5,0) node[below left] {$x$};
\draw[->,popmuted] (0,-1.6) -- (0,1.9) node[below left] {$y$};
% courbe
\draw[popblue,very thick] plot[domain=0:4,samples=80] (\x,{0.6+0.35*(\x)-0.06*(\x)^2}) node[above] {$y=f(x)$};
\draw[popblue!60] plot[domain=0:4,samples=80] (\x,{-0.6-0.35*(\x)+0.06*(\x)^2});
% ellipses extrémités
\draw[popblue] (0,0) ellipse (0.12 and 0.6);
\draw[popblue] (4,0) ellipse (0.15 and 1.04);
% tranche cylindrique
\def\xa{2.2}\def\xb{2.7}
\pgfmathsetmacro\ra{0.6+0.35*\xa-0.06*\xa*\xa}
\pgfmathsetmacro\rb{0.6+0.35*\xb-0.06*\xb*\xb}
\fill[poporangeL,opacity=0.8] (\xa,-\ra) -- (\xb,-\rb) -- (\xb,\rb) -- (\xa,\ra) -- cycle;
\draw[poporange,thick] (\xa,0) ellipse (0.13 and \ra);
\draw[poporange,thick] (\xb,0) ellipse (0.14 and \rb);
\draw[poporange,thick] (\xa,\ra) -- (\xb,\rb);
\draw[poporange,thick] (\xa,-\ra) -- (\xb,-\rb);
% annotations
\draw[<->,popdark] (\xa,-0.35) -- (\xa,-\ra+0.06) ;
\node[popdark,left] at (\xa,-0.9) {\footnotesize $f(x)$};
\draw[<->,popdark] (\xa,1.35) -- (\xb,1.35) node[midway,above] {\footnotesize $\mathrm{d}x$};
\draw[dashed,popmuted] (\xa,0) -- (\xa,\ra);
\node at (3.4,-1.35) {\footnotesize tranche : cylindre de volume $\pi[f(x)]^2\mathrm{d}x$};
\end{tikzpicture}
\end{center}
\legende{La rotation de la courbe de $f$ autour de $(Ox)$ engendre un solide. Une tranche d'épaisseur $\mathrm{d}x$ est un cylindre de rayon $f(x)$, donc de volume $\pi[f(x)]^2\,\mathrm{d}x$ ; le volume total est la somme intégrale de ces tranches.}
\end{popfigure}

\begin{exemplebox}[Volume du cône : on retrouve la formule du collège]
Faisons tourner le segment d'équation $y = \dfrac{r}{h}x$ sur $[0\,;h]$ autour de $(Ox)$ : on obtient un cône de révolution de hauteur $h$ et de rayon de base $r$. Son volume :
\begin{align*}
V &= \pi \int_0^h \left(\dfrac{r}{h}x\right)^2 \mathrm{d}x = \pi\,\dfrac{r^2}{h^2}\int_0^h x^2\,\mathrm{d}x = \pi\,\dfrac{r^2}{h^2}\left[\dfrac{x^3}{3}\right]_0^h = \pi\,\dfrac{r^2}{h^2}\cdot\dfrac{h^3}{3} = \dfrac{1}{3}\pi r^2 h.
\end{align*}
C'est exactement la formule apprise au collège — cette fois \emph{démontrée}.
\end{exemplebox}

\begin{popfigure}
\begin{center}
\begin{tikzpicture}[scale=1.1]
\draw[->,popmuted] (-0.3,0) -- (5,0) node[below left] {$x$};
\draw[->,popmuted] (0,-2.6) -- (0,2.9) node[below left] {$y$};
% segment générateur
\draw[popblue,very thick] (0,0) -- (4,2) node[right] {$y=\dfrac{x}{2}$};
\draw[popblue!60] (0,0) -- (4,-2);
% cône
\draw[popdark,thick] (4,0) ellipse (0.35 and 2);
\draw[popdark,thick] (0,0) -- (4,2);
\draw[popdark,thick] (0,0) -- (4,-2);
% cotes
\draw[<->,poporange,thick] (4.6,0) -- (4.6,2) node[midway,right] {$r=2$};
\draw[<->,poporange,thick] (0,-2.4) -- (4,-2.4) node[midway,below] {$h=4$};
\filldraw[popdark] (0,0) circle (0.05) node[above left] {$O$};
\draw[dashed,popmuted] (4,0) -- (4,2);
\end{tikzpicture}
\end{center}
\legende{Rotation du segment $y=\dfrac{x}{2}$ sur $[0\,;4]$ autour de $(Ox)$ : le solide engendré est un cône de hauteur $h=4$ et de rayon $r=2$, de volume $V=\pi\int_0^4 \frac{x^2}{4}\,\mathrm{d}x = \frac{16\pi}{3}$ u.v.}
\end{popfigure}

\begin{exoresolu}[Volume d'une sphère]
\textbf{Énoncé.} Le demi-cercle d'équation $y = \sqrt{R^2 - x^2}$, $x \in [-R\,;R]$, tourne autour de $(Ox)$. Montrer que le volume de la sphère obtenue est $V = \dfrac{4}{3}\pi R^3$.
\tcblower
\textbf{Solution.} On a $[f(x)]^2 = R^2 - x^2$, donc :
\begin{align*}
V &= \pi \int_{-R}^{R} \left(R^2 - x^2\right)\mathrm{d}x = \pi\left[R^2 x - \dfrac{x^3}{3}\right]_{-R}^{R} \\
&= \pi\left[\left(R^3 - \dfrac{R^3}{3}\right) - \left(-R^3 + \dfrac{R^3}{3}\right)\right] = \pi \times \dfrac{4R^3}{3} = \dfrac{4}{3}\pi R^3.
\end{align*}
Application numérique : pour un ballon de rayon $R = 11$ cm, $V = \dfrac{4}{3}\pi \times 1331 \approx \SI{5575}{cm^3}$, soit environ $5,6$ litres.
\end{exoresolu}

\begin{attention}[Les trois erreurs qui coûtent des points au Bac]
\begin{itemize}
\item \textbf{Oublier le $\pi$ ou le carré} : la formule est $V = \pi\int_a^b [f(x)]^2\,\mathrm{d}x$, pas $\int_a^b f$. Avant de calculer, écris d'abord $[f(x)]^2$ explicitement.
\item \textbf{Oublier la conversion} : une aire en u.a. n'est en cm$^2$ qu'après multiplication par $p \times q$ ; un volume en u.v. n'est en cm$^3$ que si l'unité des axes est le centimètre ($1\ \text{u.v.} = p^3$ si les axes ont la même unité $p$).
\item \textbf{Intégrer $f-g$ alors que les courbes se croisent} sur l'intervalle : il faut découper par Chasles et intégrer $|f-g|$ morceau par morceau, sinon les aires se compensent et le résultat est faux.
\end{itemize}
\end{attention}

\begin{savaistu}[Archimède, 2000 ans avant Newton]
La méthode des tranches cylindriques n'a rien de moderne : c'est \cle{Archimède} de Syracuse (III\textsuperscript{e} siècle av. J.-C.) qui l'inventa sous le nom de « méthode d'exhaustion ». Il calcula ainsi le volume de la sphère et découvrit que celui-ci vaut les deux tiers du volume du cylindre qui la contient — résultat si beau à ses yeux qu'il demanda qu'on grave une sphère inscrite dans un cylindre sur sa tombe. Il faudra attendre Newton et Leibniz, près de 2000 ans plus tard, pour transformer cette intuition en calcul intégral systématique. Aujourd'hui, les ingénieurs de la SONARA ou de la CAMTEL utilisent exactement ces formules pour dimensionner cuves, réservoirs et silos.
\end{savaistu}

\begin{exercice}[Silo à maïs]
Un silo est modélisé par la rotation autour de $(Ox)$ de la courbe de $f(x) = \sqrt{x+1}$ sur $[0\,;3]$ (unité : $1$ m sur chaque axe).
\begin{enumerate}
\item Calculer le volume du silo en m$^3$.
\item En déduire sa contenance en litres ($1\ \text{m}^3 = 1000$ L).
\end{enumerate}
\end{exercice}

\begin{corrige}[Exercice : silo à maïs]
\textbf{1)} $V = \pi\displaystyle\int_0^3 (x+1)\,\mathrm{d}x = \pi\left[\dfrac{x^2}{2}+x\right]_0^3 = \pi\left(\dfrac{9}{2}+3\right) = \dfrac{15\pi}{2}\ \text{m}^3 \approx 23,56\ \text{m}^3$.

\textbf{2)} Contenance : $\dfrac{15\pi}{2}\times 1000 \approx \num{23562}$ litres.
\end{corrige}

Tu disposes désormais de toute la panoplie : aires sous une courbe, aires entre deux courbes, volumes de révolution. La dernière section te montrera comment \emph{approcher} une intégrale lorsqu'aucune primitive n'est disponible : la méthode des rectangles.

\coursec{Approximation d'une intégrale : méthode des rectangles}

Dans la section précédente, nous avons calculé des aires et des volumes de révolution grâce aux primitives : tout reposait sur la formule $\displaystyle\int_a^b f(x)\,dx = F(b)-F(a)$. Mais que faire lorsqu'on ne connaît \emph{aucune} primitive de $f$ ? Ce cas n'est pas exotique : il est la règle dans les applications réelles. Cette dernière section répond à la question par une idée d'une simplicité remarquable, vieille de plus de deux mille ans : découper l'aire en \cle{rectangles}.

\courssub{Pourquoi approcher une intégrale ?}

\begin{attention}[Des fonctions sans primitive « exprimable »]
Certaines fonctions pourtant très simples et continues n'admettent pas de primitive exprimable à l'aide des fonctions usuelles (polynômes, exponentielle, logarithme, fonctions trigonométriques). C'est un résultat mathématique démontré (théorie de Liouville), pas un manque d'astuce de notre part. Par exemple :
\[ f(x)=e^{x^2}, \qquad g(x)=e^{-x^2}, \qquad h(x)=\frac{\sin x}{x}, \qquad k(x)=\sqrt{1+x^3}. \]
Pourtant, ces fonctions étant continues, leurs intégrales \emph{existent}. L'intégrale $\displaystyle\int_0^1 e^{-x^2}\,dx$ intervient par exemple en probabilités (loi normale), et $\displaystyle\int_0^x \frac{\sin t}{t}\,dt$ en traitement du signal. Il faut donc savoir en calculer des \cle{valeurs approchées}.
\end{attention}

L'idée est géométrique : l'intégrale d'une fonction continue positive est une aire. Or une aire « compliquée » peut être approchée par une somme d'aires de rectangles, figures dont l'aire est élémentaire : $\text{base}\times\text{hauteur}$. Plus les rectangles sont fins, meilleure est l'approximation.

\begin{savaistu}[Des compteurs qui intègrent en continu]
Les caisses enregistreuses des stations-service de Douala ou de Yaoundé mesurent le \emph{débit} de carburant à chaque instant et « intègrent » ce débit pour afficher le volume délivré et le prix à payer. De même, les compteurs d'eau de la CAMWATER cumulent le débit pour obtenir le volume consommé. Concrètement, l'appareil mesure le débit à intervalles réguliers très courts et additionne les petits volumes $\text{débit}\times\text{durée}$ : c'est exactement une méthode des rectangles, exécutée en temps réel ! Archimède, lui, calculait déjà l'aire d'une parabole en l'approchant par des sommes d'aires de figures simples, trois siècles avant notre ère.
\end{savaistu}

\courssub{Subdivision régulière et sommes de rectangles}

\begin{definition}[{Subdivision régulière de $[a\,;b]$}]
Soit $f$ une fonction continue sur $[a\,;b]$ et $n$ un entier naturel non nul. On partage $[a\,;b]$ en $n$ intervalles de même longueur
\[ h=\frac{b-a}{n} \quad (\text{appelée \cle{pas} ou \cle{amplitude} de la subdivision}). \]
Les points de la subdivision sont
\[ x_k = a + k\,h, \qquad k\in\{0,1,\dots,n\}, \]
de sorte que $x_0=a$, $x_n=b$ et $x_{k+1}-x_k=h$.
\end{definition}

Sur chaque intervalle $[x_k\,;\,x_{k+1}]$, on remplace la courbe par un segment horizontal, c'est-à-dire qu'on approche l'aire sous la courbe par celle d'un rectangle de base $h$. Deux choix naturels s'offrent à nous pour la hauteur :

\begin{definition}[Sommes des rectangles à gauche et à droite]
Pour $f$ continue et \textbf{positive} sur $[a\,;b]$ :
\begin{itemize}
\item la \cle{somme des rectangles à gauche} : $\displaystyle G_n = h\sum_{k=0}^{n-1} f(x_k) = h\big[f(x_0)+f(x_1)+\dots+f(x_{n-1})\big]$ ;
\item la \cle{somme des rectangles à droite} : $\displaystyle D_n = h\sum_{k=0}^{n-1} f(x_{k+1}) = h\big[f(x_1)+f(x_2)+\dots+f(x_{n})\big]$.
\end{itemize}
Chaque terme $h\times f(x_k)$ est l'aire d'un rectangle de base $h$ et de hauteur $f(x_k)$.
\end{definition}

\begin{attention}[Le piège classique : oublier le facteur $h$]
Une erreur très fréquente au Bac consiste à additionner les valeurs $f(x_k)$ et à donner cette somme comme réponse. La somme $\sum f(x_k)$ est une somme de \emph{hauteurs} ; une aire exige de la multiplier par la base $h=\dfrac{b-a}{n}$. Retiens : \textbf{somme des hauteurs, puis multiplication par $h$}.
\end{attention}

\begin{popfigure}
\begin{center}
\begin{tikzpicture}
\begin{axis}[width=0.92\linewidth, height=6.5cm, axis lines=middle, xlabel={$x$}, ylabel={$y$}, xmin=-0.1, xmax=1.35, ymin=0, ymax=1.6, xtick={0,0.2,0.4,0.6,0.8,1}, xticklabels={$x_0$,$x_1$,$x_2$,$x_3$,$x_4$,$x_5$}, ytick=\empty, legend style={at={(0.98,0.55)},anchor=east,font=\small}]
\addplot[popcurve, domain=0:1.15, samples=100] {0.4+x+0.15*x^2};
\addlegendentry{$\mathcal{C}_f$ (croissante)}
\foreach \i/\x in {0/0, 1/0.2, 2/0.4, 3/0.6, 4/0.8}{
  \addplot[draw=popblue, fill=popblueL, fill opacity=0.75] coordinates {(\x,0) (\x,{0.4+\x+0.15*\x*\x}) (\x+0.2,{0.4+\x+0.15*\x*\x}) (\x+0.2,0)} \closedcycle;
}
\end{axis}
\end{tikzpicture}
\legende{Rectangles à gauche pour une fonction croissante ($n=5$) : la hauteur est prise en l'extrémité gauche de chaque intervalle ; l'aire totale des rectangles \emph{sous-estime} l'aire sous la courbe.}
\end{center}
\end{popfigure}

\begin{popfigure}
\begin{center}
\begin{tikzpicture}
\begin{axis}[width=0.92\linewidth, height=6.5cm, axis lines=middle, xlabel={$x$}, ylabel={$y$}, xmin=-0.1, xmax=1.35, ymin=0, ymax=1.6, xtick={0,0.2,0.4,0.6,0.8,1}, xticklabels={$x_0$,$x_1$,$x_2$,$x_3$,$x_4$,$x_5$}, ytick=\empty, legend style={at={(0.98,0.55)},anchor=east,font=\small}]
\addplot[popcurve, domain=0:1.15, samples=100] {0.4+x+0.15*x^2};
\addlegendentry{$\mathcal{C}_f$ (croissante)}
\foreach \i/\x in {0/0, 1/0.2, 2/0.4, 3/0.6, 4/0.8}{
  \addplot[draw=poporange, fill=poporangeL, fill opacity=0.8] coordinates {(\x,0) (\x,{0.4+\x+0.2+0.15*(\x+0.2)*(\x+0.2)}) (\x+0.2,{0.4+\x+0.2+0.15*(\x+0.2)*(\x+0.2)}) (\x+0.2,0)} \closedcycle;
}
\end{axis}
\end{tikzpicture}
\legende{Rectangles à droite pour la même fonction ($n=5$) : la hauteur est prise en l'extrémité droite ; l'aire totale \emph{surestime} l'aire sous la courbe. L'intégrale est « coincée » entre les deux sommes.}
\end{center}
\end{popfigure}

\courssub{Encadrement de l'intégrale pour une fonction monotone}

\begin{propriete}[Encadrement par les sommes de rectangles]
Soit $f$ une fonction continue, positive et \cle{monotone} sur $[a\,;b]$. Alors :
\begin{itemize}
\item si $f$ est \textbf{croissante} sur $[a\,;b]$ : $\quad G_n \;\le\; \displaystyle\int_a^b f(x)\,dx \;\le\; D_n$ ;
\item si $f$ est \textbf{décroissante} sur $[a\,;b]$ : $\quad D_n \;\le\; \displaystyle\int_a^b f(x)\,dx \;\le\; G_n$.
\end{itemize}
De plus, dans les deux cas, l'écart entre les deux sommes vaut :
\[ |D_n - G_n| = h\,\big|f(b)-f(a)\big| = \frac{b-a}{n}\,\big|f(b)-f(a)\big|. \]
\end{propriete}

\begin{experience}[Justification graphique de l'encadrement]
Supposons $f$ croissante et positive sur $[a\,;b]$, et raisonnons sur un seul intervalle $[x_k\,;\,x_{k+1}]$.
\begin{enumerate}
\item Comme $f$ est croissante, pour tout $x\in[x_k\,;\,x_{k+1}]$ : $\; f(x_k)\le f(x)\le f(x_{k+1})$.
\item Par positivité de l'intégrale, en intégrant cet encadrement sur $[x_k\,;\,x_{k+1}]$ :
\[ h\,f(x_k) \;\le\; \int_{x_k}^{x_{k+1}} f(x)\,dx \;\le\; h\,f(x_{k+1}). \]
Graphiquement : le rectangle de hauteur $f(x_k)$ est \emph{sous} la courbe, celui de hauteur $f(x_{k+1})$ la \emph{dépasse}.
\item On somme ces inégalités pour $k$ de $0$ à $n-1$. Par la relation de Chasles, les intégrales s'ajoutent en $\displaystyle\int_a^b f$ ; les membres de gauche et de droite donnent $G_n$ et $D_n$. D'où $G_n\le \displaystyle\int_a^b f \le D_n$.
\item Pour l'écart, observons que la différence $D_n-G_n$ est une somme télescopique :
\begin{align*}
D_n - G_n &= h\big[f(x_1)-f(x_0)\big] + h\big[f(x_2)-f(x_1)\big] + \dots + h\big[f(x_n)-f(x_{n-1})\big] \\
&= h\big[f(x_n)-f(x_0)\big] = h\big[f(b)-f(a)\big].
\end{align*}
Graphiquement, si l'on fait glisser horizontalement tous les petits rectangles d'« excès » (entre courbe et rectangles à droite), ils s'empilent exactement dans une colonne de largeur $h$ et de hauteur $f(b)-f(a)$.
\end{enumerate}
Le cas décroissant est identique, les inégalités étant renversées. \hfill$\blacksquare$
\end{experience}

\begin{attention}[Fonction décroissante : l'encadrement s'inverse]
Pour une fonction \textbf{décroissante}, c'est la somme à gauche $G_n$ qui \emph{surestime} l'intégrale, et la somme à droite $D_n$ qui la \emph{sous-estime}. Avant d'écrire l'encadrement, vérifie toujours le sens de variation de $f$ sur $[a\,;b]$. Astuce de rédaction : écris l'encadrement sous la forme $\min(G_n,D_n)\le \displaystyle\int_a^b f\le \max(G_n,D_n)$ si tu as un doute, puis précise l'ordre grâce à la monotonie.
\end{attention}

\begin{popfigure}
\begin{center}
\begin{tikzpicture}
\begin{axis}[width=0.45\linewidth, height=5.5cm, axis lines=middle, xlabel={$x$}, ylabel={$y$}, xmin=-0.1, xmax=1.3, ymin=0, ymax=1.7, xtick=\empty, ytick=\empty, title={\small $n=4$ : encadrement grossier}, title style={font=\small}]
\addplot[popcurve, domain=0:1.1, samples=100] {0.4+x+0.15*x^2};
\foreach \x in {0, 0.25, 0.5, 0.75}{
  \addplot[draw=popblue, fill=popblueL, fill opacity=0.75] coordinates {(\x,0) (\x,{0.4+\x+0.15*\x*\x}) (\x+0.25,{0.4+\x+0.15*\x*\x}) (\x+0.25,0)} \closedcycle;
}
\end{axis}
\end{tikzpicture}
\hfill
\begin{tikzpicture}
\begin{axis}[width=0.45\linewidth, height=5.5cm, axis lines=middle, xlabel={$x$}, ylabel={$y$}, xmin=-0.1, xmax=1.3, ymin=0, ymax=1.7, xtick=\empty, ytick=\empty, title={\small $n=16$ : encadrement précis}, title style={font=\small}]
\addplot[popcurve, domain=0:1.1, samples=100] {0.4+x+0.15*x^2};
\foreach \x in {0, 0.0625, 0.125, 0.1875, 0.25, 0.3125, 0.375, 0.4375, 0.5, 0.5625, 0.625, 0.6875, 0.75, 0.8125, 0.875, 0.9375}{
  \addplot[draw=popgreen, fill=popgreenL, fill opacity=0.75] coordinates {(\x,0) (\x,{0.4+\x+0.15*\x*\x}) (\x+0.0625,{0.4+\x+0.15*\x*\x}) (\x+0.0625,0)} \closedcycle;
}
\end{axis}
\end{tikzpicture}
\legende{Quand $n$ augmente, le pas $h=\frac{b-a}{n}$ diminue et l'écart $h\,|f(b)-f(a)|$ tend vers $0$ : les rectangles épousent de mieux en mieux la courbe et l'encadrement se resserre autour de la valeur exacte de l'intégrale.}
\end{center}
\end{popfigure}

\courssub{Méthode pratique et contrôle de la précision}

\begin{methode}[Méthode des rectangles : mise en œuvre]
Pour approcher $\displaystyle\int_a^b f(x)\,dx$ avec $n$ rectangles :
\begin{enumerate}
\item \textbf{Calculer le pas} : $h=\dfrac{b-a}{n}$.
\item \textbf{Dresser le tableau de valeurs} : calculer $f(x_0), f(x_1), \dots, f(x_n)$ où $x_k=a+kh$ (à la calculatrice, en mode table avec un pas égal à $h$).
\item \textbf{Sommer} : pour les rectangles à gauche, additionner $f(x_0)$ à $f(x_{n-1})$ ; pour les rectangles à droite, additionner $f(x_1)$ à $f(x_n)$.
\item \textbf{Multiplier par $h$} : $G_n=h\sum_{k=0}^{n-1}f(x_k)$ et $D_n=h\sum_{k=0}^{n-1}f(x_{k+1})$.
\item \textbf{Conclure} par l'encadrement adapté au sens de variation de $f$, et donner éventuellement la valeur approchée $\dfrac{G_n+D_n}{2}$.
\end{enumerate}
\end{methode}

\begin{methode}[Obtenir une précision $\varepsilon$ imposée]
Pour que l'encadrement ait une amplitude au plus égale à $\varepsilon$ (fonction monotone) :
\begin{enumerate}
\item Écrire la condition $|D_n-G_n|\le\varepsilon$, c'est-à-dire $\dfrac{b-a}{n}\,|f(b)-f(a)|\le\varepsilon$.
\item En déduire $\; n \ge \dfrac{(b-a)\,|f(b)-f(a)|}{\varepsilon}$.
\item Choisir pour $n$ le \textbf{plus petit entier} vérifiant cette inégalité (partie entière supérieure).
\item Calculer alors $G_n$ et $D_n$ : toute valeur de l'encadrement (par exemple le milieu $\frac{G_n+D_n}{2}$) approche l'intégrale à $\varepsilon$ près.
\end{enumerate}
\end{methode}

\begin{exemplebox}[Encadrement de $\displaystyle\int_0^1 e^{-x^2}\,dx$ avec $n=4$]
La fonction $f(x)=e^{-x^2}$ est continue, positive et \textbf{décroissante} sur $[0\,;1]$ (sa dérivée $-2x\,e^{-x^2}$ est négative sur $]0;1]$). Le pas est $h=\dfrac{1-0}{4}=0{,}25$. Tableau de valeurs (arrondis à $10^{-4}$) :
\begin{center}
\begin{tabular}{|l|c|c|c|c|c|}
\hline
\rowcolor{popblueL} $x_k$ & $0$ & $0{,}25$ & $0{,}5$ & $0{,}75$ & $1$ \\
\hline
$f(x_k)=e^{-x_k^2}$ & $1$ & $0{,}9394$ & $0{,}7788$ & $0{,}5698$ & $0{,}3679$ \\
\hline
\end{tabular}
\end{center}
Somme à gauche (qui surestime, car $f$ décroît) :
\[ G_4 = 0{,}25\,(1+0{,}9394+0{,}7788+0{,}5698) = 0{,}25\times 3{,}2880 \approx 0{,}8220. \]
Somme à droite (qui sous-estime) :
\[ D_4 = 0{,}25\,(0{,}9394+0{,}7788+0{,}5698+0{,}3679) = 0{,}25\times 2{,}6559 \approx 0{,}6640. \]
D'où l'encadrement : $\; 0{,}664 \le \displaystyle\int_0^1 e^{-x^2}\,dx \le 0{,}822$.
L'amplitude vaut bien $h\,|f(1)-f(0)| = 0{,}25\times(1-e^{-1}) \approx 0{,}25\times 0{,}6321 \approx 0{,}158$, conformément à la théorie. La valeur approchée par le milieu est $\dfrac{G_4+D_4}{2}\approx 0{,}743$ (la valeur exacte, obtenue par des méthodes plus fines, vaut environ $0{,}747$ : notre approximation avec seulement 4 rectangles est déjà honorable !).
\end{exemplebox}

\begin{exoresolu}[Choisir $n$ pour une précision donnée]
\textbf{Énoncé.} Soit $f(x)=e^{-x^2}$ sur $[0\,;1]$. Déterminer un entier $n$ tel que l'encadrement de $\displaystyle I=\int_0^1 e^{-x^2}\,dx$ par la méthode des rectangles ait une amplitude inférieure ou égale à $\varepsilon = 10^{-2}$. Donner alors une valeur approchée de $I$ à $10^{-2}$ près.
\tcblower
\textbf{Solution détaillée.}
\begin{enumerate}
\item $f$ est décroissante sur $[0\,;1]$, donc l'amplitude de l'encadrement est
\[ G_n - D_n = \frac{b-a}{n}\big|f(b)-f(a)\big| = \frac{1}{n}\big(1-e^{-1}\big). \]
\item On résout : $\dfrac{1-e^{-1}}{n}\le 10^{-2} \iff n \ge 100\,(1-e^{-1}) \approx 63{,}2$.
\item Le plus petit entier convenable est donc $\boxed{n=64}$.
\item Avec $n=64$, on est assuré que $G_{64}-D_{64}\le 10^{-2}$. Le milieu $\dfrac{G_{64}+D_{64}}{2}$ diffère alors de $I$ d'au plus $\dfrac{10^{-2}}{2}=0{,}005$. Un calcul à la calculatrice ou sur tableur donne $D_{64}\approx 0{,}7418$ et $G_{64}\approx 0{,}7517$, d'où l'encadrement rigoureux $0{,}7418 \le I \le 0{,}7517$. Une valeur approchée de $I$ à $10^{-2}$ près (arrondie) est donc $\boxed{0{,}75}$.
\end{enumerate}
\end{exoresolu}

\courssub{Algorithme de calcul (complément utile au Bac)}

La somme de cent ou mille termes ne se fait pas à la main : on la confie à une machine. Voici l'algorithme des rectangles à gauche, en langage naturel puis en Python.

\begin{methode}[Algorithme des rectangles à gauche]
\textbf{Langage naturel :}
\begin{enumerate}
\item Saisir $a$, $b$, $n$ et définir la fonction $f$.
\item $h \leftarrow (b-a)/n$ ; \quad $S \leftarrow 0$.
\item Pour $k$ allant de $0$ à $n-1$ : \quad $S \leftarrow S + f(a+k\,h)$.
\item Afficher $h\times S$.
\end{enumerate}
\textbf{En Python :}
\begin{center}
\begin{tabular}{l}
\texttt{from math import exp}\\
\texttt{def f(x): return exp(-x**2)}\\
\texttt{a, b, n = 0, 1, 1000}\\
\texttt{h = (b - a)/n}\\
\texttt{S = sum(f(a + k*h) for k in range(n))}\\
\texttt{print(h * S)  \# affiche environ 0,747}
\end{tabular}
\end{center}
Pour les rectangles à droite, il suffit de faire varier $k$ de $1$ à $n$ (c'est-à-dire remplacer \texttt{range(n)} par \texttt{range(1, n+1)}).
\end{methode}

\begin{aretenir}[L'essentiel de la méthode des rectangles]
\begin{itemize}
\item Subdivision régulière : $h=\dfrac{b-a}{n}$, \; $x_k=a+kh$.
\item Rectangles à gauche : $G_n=h\displaystyle\sum_{k=0}^{n-1}f(x_k)$ ; à droite : $D_n=h\displaystyle\sum_{k=0}^{n-1}f(x_{k+1})$.
\item Si $f$ est \textbf{croissante} : $G_n\le\displaystyle\int_a^b f\le D_n$ ; si $f$ est \textbf{décroissante} : $D_n\le\displaystyle\int_a^b f\le G_n$.
\item Amplitude de l'encadrement : $|D_n-G_n|=\dfrac{b-a}{n}\,|f(b)-f(a)|$, qui tend vers $0$ quand $n$ grandit.
\item Précision $\varepsilon$ : choisir $n\ge\dfrac{(b-a)\,|f(b)-f(a)|}{\varepsilon}$ et prendre $\dfrac{G_n+D_n}{2}$ comme valeur approchée.
\item Ne jamais oublier de multiplier la somme des hauteurs par $h$.
\end{itemize}
\end{aretenir}

\begin{exercice}[À toi de jouer]
Soit $f(x)=\dfrac{1}{1+x^2}$ sur $[0\,;1]$.
\begin{enumerate}
\item Justifier que $f$ est décroissante sur $[0\,;1]$.
\item Avec $n=4$, calculer $G_4$ et $D_4$ (on donnera les valeurs exactes sous forme fractionnaire avant l'arrondi).
\item En déduire un encadrement de $\displaystyle\int_0^1 \frac{dx}{1+x^2}$ et vérifier que l'amplitude vaut $\dfrac{1}{8}$.
\item Quelle valeur de $n$ garantit une amplitude inférieure à $10^{-3}$ ?
\end{enumerate}
\end{exercice}

\begin{corrige}[Exercice : méthode des rectangles pour $1/(1+x^2)$]
\begin{enumerate}
\item $f'(x)=\dfrac{-2x}{(1+x^2)^2}\le 0$ sur $[0\,;1]$, donc $f$ est décroissante sur $[0\,;1]$.
\item $h=\dfrac{1}{4}=0{,}25$ ; $f(0)=1$, $f(0{,}25)=\dfrac{16}{17}\approx 0{,}9412$, $f(0{,}5)=\dfrac{4}{5}=0{,}8$, $f(0{,}75)=\dfrac{16}{25}=0{,}64$, $f(1)=\dfrac{1}{2}=0{,}5$.
\[ G_4 = 0{,}25\left(1+\tfrac{16}{17}+\tfrac{4}{5}+\tfrac{16}{25}\right) \approx 0{,}8453, \qquad D_4 = 0{,}25\left(\tfrac{16}{17}+\tfrac{4}{5}+\tfrac{16}{25}+\tfrac{1}{2}\right) \approx 0{,}7203. \]
\item Comme $f$ décroît : $D_4\le \displaystyle\int_0^1\frac{dx}{1+x^2}\le G_4$, soit $0{,}7203 \le I \le 0{,}8453$. Amplitude : $G_4-D_4 = h\big(f(0)-f(1)\big) = 0{,}25\times\left(1-\tfrac12\right)=\dfrac18=0{,}125$. \hfill (On sait par ailleurs que $I=\dfrac{\pi}{4}\approx 0{,}785$ : bien dans l'encadrement !)
\item Il faut $\dfrac{1}{n}\times\dfrac12 \le 10^{-3}$, soit $n\ge 500$. On choisit $n=500$.
\end{enumerate}
\end{corrige}

Avec cette méthode des rectangles s'achève notre étude du calcul intégral : tu sais désormais définir une intégrale, en exploiter les propriétés, la calculer exactement par des primitives, l'intégration par parties ou un changement de variable, l'appliquer aux aires et aux volumes, et — lorsque le calcul exact est impossible — l'encadrer avec une précision entièrement contrôlée. C'est cette combinaison de rigueur et d'efficacité qui fait de l'intégration l'un des outils les plus puissants des mathématiques du Baccalauréat... et des sciences de l'ingénieur.

\newpage

\coursec{Activité d'intégration}

\courssub{Objectifs de l'activité}

Cette activité d'intégration te place dans une situation authentique de gestion de l'eau potable dans la région de l'Ouest du Cameroun. En travaillant par groupes de trois ou quatre élèves, tu vas mobiliser \textbf{l'ensemble des savoir-faire} de la leçon « Calcul intégral » :

\begin{itemize}
\item modéliser une grandeur physique (un débit) par une fonction et interpréter une intégrale comme une quantité accumulée (un volume) ;
\item étudier les variations d'une fonction et exploiter sa représentation graphique ;
\item calculer une intégrale par \cle{lecture directe d'une primitive} et par une \cle{intégration par parties} guidée ;
\item calculer la \cle{valeur moyenne} d'une fonction sur un intervalle et l'interpréter concrètement ;
\item vérifier un résultat par la \cle{méthode des rectangles} et discuter la précision de l'approximation ;
\item calculer le \cle{volume d'un solide de révolution} et confronter deux grandeurs pour prendre une décision.
\end{itemize}

\courssub{Situation}

La coopérative agricole « Nyambé Progress », située à Bafoussam (région de l'Ouest), alimente ses plantations de maraîchage grâce à un réservoir d'eau collectif. Le réservoir est rempli chaque jour par une pompe solaire couplée à un forage. L'ingénieur hydraulicien de la coopérative a enregistré le débit d'entrée de l'eau pendant plusieurs journées ensoleillées et a constaté que le débit (exprimé en mètres cubes par heure, m$^3$/h) peut être modélisé, pour une journée type, par la fonction :
\[ d(t) = t\,e^{-t/6} + 1, \qquad t \in [0\,;\,12], \]
où $t$ désigne le temps écoulé, en heures, à partir de 6 h du matin (ainsi $t = 0$ correspond à 6 h et $t = 12$ correspond à 18 h). Le terme constant $1$ correspond au débit minimal de la pompe ; le terme $t\,e^{-t/6}$ traduit le renforcement puis l'affaiblissement de l'ensoleillement au cours de la journée.

Le réservoir lui-même est un \textbf{solide de révolution} : sa paroi interne est engendrée par la rotation autour de l'axe $(Ox)$ de la courbe d'équation $y = 2\sqrt{x}$ pour $x \in [0\,;\,4]$, les unités étant exprimées en mètres. Il est initialement vide à 6 h.

Le président de la coopérative se pose trois questions :
\begin{enumerate}
\item Quel volume total d'eau entre dans le réservoir entre 6 h et 18 h ?
\item Quel est le débit moyen sur cette journée ?
\item Le réservoir est-il assez grand pour contenir toute l'eau collectée ?
\end{enumerate}

\courssub{Consignes}

Travail par groupes, puis restitution en grand groupe. Toutes les valeurs approchées seront données à $10^{-2}$ près.

\begin{enumerate}
\item \textbf{Étude du débit.} Justifier que $d$ est dérivable sur $[0\,;\,12]$ et montrer que $d'(t) = e^{-t/6}\left(1 - \dfrac{t}{6}\right)$. Dresser le tableau de variations de $d$ sur $[0\,;\,12]$. À quelle heure le débit est-il maximal ? Tracer l'allure de la courbe de $d$ dans un repère orthogonal (unités : 1 cm pour 2 h en abscisse, 2 cm pour 1 m$^3$/h en ordonnée).
\item \textbf{Volume total d'eau.} Vérifier que $d$ est continue et positive sur $[0\,;\,12]$. Expliquer pourquoi le volume total d'eau entré entre 6 h et 18 h est $V_{\text{eau}} = \displaystyle\int_0^{12} d(t)\,\mathrm{d}t$, et préciser son unité.
\item \textbf{Intégration par parties guidée.} On pose $I = \displaystyle\int_0^{12} t\,e^{-t/6}\,\mathrm{d}t$. En posant $u(t) = t$ et $v'(t) = e^{-t/6}$, calculer $I$ à l'aide de la formule d'intégration par parties (on choisira $v(t) = -6e^{-t/6}$).
\item En déduire la valeur exacte de $V_{\text{eau}}$, puis une valeur approchée à $10^{-2}$ près.
\item \textbf{Débit moyen.} Calculer le débit moyen $\bar{d}$ sur la journée. Interpréter ce résultat pour le président de la coopérative.
\item \textbf{Vérification par les rectangles.} Partager $[0\,;\,12]$ en $n = 6$ intervalles d'amplitude $h = 2$. Calculer la somme des aires des rectangles à gauche $R_g$ et à droite $R_d$. Comparer avec la valeur exacte : l'ordre de grandeur est-il confirmé ? Discuter.
\item \textbf{Capacité du réservoir.} Calculer le volume $V_{\text{rés}} = \pi \displaystyle\int_0^{4} \big(2\sqrt{x}\big)^2\,\mathrm{d}x$. Toute l'eau collectée tient-elle dans le réservoir ?
\end{enumerate}

\courssub{Ressources à mobiliser}

\begin{aretenir}[Boîte à outils de la leçon]
\begin{itemize}
\item \textbf{Interprétation} : si $f$ est continue et positive sur $[a\,;\,b]$, $\displaystyle\int_a^b f(t)\,\mathrm{d}t$ mesure l'aire sous la courbe ; si $f$ est un débit, c'est la quantité accumulée entre les instants $a$ et $b$.
\item \textbf{Linéarité} : $\displaystyle\int_a^b (f + g) = \int_a^b f + \int_a^b g$.
\item \textbf{Intégration par parties} : si $u$ et $v$ sont dérivables, de dérivées continues sur $[a\,;\,b]$, alors $\displaystyle\int_a^b u\,v' = \big[u\,v\big]_a^b - \int_a^b u'\,v$.
\item \textbf{Valeur moyenne} : $\mu = \dfrac{1}{b-a}\displaystyle\int_a^b f(t)\,\mathrm{d}t$.
\item \textbf{Rectangles} : avec $n$ subdivisions d'amplitude $h = \dfrac{b-a}{n}$, $R_g = h\displaystyle\sum_{k=0}^{n-1} f(a + kh)$ et $R_d = h\displaystyle\sum_{k=1}^{n} f(a + kh)$.
\item \textbf{Volume de révolution} autour de $(Ox)$ : $V = \pi \displaystyle\int_a^b \big[f(x)\big]^2\,\mathrm{d}x$.
\end{itemize}
\end{aretenir}

\begin{attention}[Conditions d'application à ne pas oublier]
\begin{itemize}
\item L'intégration par parties exige que $u'$ et $v'$ soient \textbf{continues} sur $[a\,;\,b]$ : c'est le cas ici car $u(t) = t$ et $v(t) = -6e^{-t/6}$ sont de classe $\mathcal{C}^1$ sur $\mathbb{R}$.
\item L'interprétation de l'intégrale comme une aire n'est valable que si la fonction est \textbf{positive} sur l'intervalle : il faut donc le vérifier avant de conclure (question 2).
\item L'encadrement $R_g \leq \displaystyle\int_a^b f \leq R_d$ n'est garanti que si $f$ est \textbf{croissante} sur $[a\,;\,b]$ (inégalités inversées si $f$ est décroissante). Ici $d$ n'est pas monotone sur $[0\,;\,12]$ : la comparaison demandée à la question 6 est un contrôle d'ordre de grandeur, pas un encadrement rigoureux.
\end{itemize}
\end{attention}

\courssub{Démarche et corrigé commenté}

\begin{exoresolu}[Corrigé collectif de l'activité]
Résoudre l'ensemble des questions de la situation du réservoir de Bafoussam.
\tcblower
\textbf{1. Variations du débit.} La fonction $t \mapsto t\,e^{-t/6}$ est dérivable sur $\mathbb{R}$ comme produit de la fonction affine $t \mapsto t$ et de la composée $t \mapsto e^{-t/6}$, toutes deux dérivables ; donc $d$ est dérivable sur $[0\,;\,12]$. Par la formule de la dérivée d'un produit $(uv)' = u'v + uv'$ :
\[ d'(t) = 1 \times e^{-t/6} + t \times \left(-\tfrac{1}{6}\right)e^{-t/6} = e^{-t/6}\left(1 - \frac{t}{6}\right). \]
Comme $e^{-t/6} > 0$ pour tout $t$, le signe de $d'(t)$ est celui de $1 - \dfrac{t}{6}$ : $d'(t) > 0$ pour $t < 6$, $d'(6) = 0$ et $d'(t) < 0$ pour $t > 6$. Donc $d$ est \textbf{strictement croissante sur $[0\,;\,6]$} et \textbf{strictement décroissante sur $[6\,;\,12]$}. Valeurs remarquables :
\[ d(0) = 1, \qquad d(6) = 6e^{-1} + 1 \approx 3{,}21, \qquad d(12) = 12e^{-2} + 1 \approx 2{,}62. \]
Le débit est maximal à $t = 6$, c'est-à-dire à \textbf{12 h}, moment où l'ensoleillement est le plus fort : le modèle est cohérent avec la réalité. La courbe monte de $1$ à environ $3{,}21$ puis redescend vers $2{,}62$ (voir la figure ci-dessous).

\textbf{2. Volume total.} Pour tout $t \in [0\,;\,12]$, $t\,e^{-t/6} \geq 0$, donc $d(t) \geq 1 > 0$ : la fonction $d$ est continue et strictement positive sur $[0\,;\,12]$. Le débit $d(t)$ est la « vitesse » à laquelle le volume augmente : le volume entré entre les instants $0$ et $12$ est l'accumulation du débit, c'est-à-dire
\[ V_{\text{eau}} = \int_0^{12} d(t)\,\mathrm{d}t, \]
exprimé en m$^3$ (produit d'un débit en m$^3$/h par un temps en h). Graphiquement, puisque $d$ est positive, c'est l'aire du domaine situé sous la courbe tracée à la question 1, au-dessus de l'axe des abscisses, entre les droites $t = 0$ et $t = 12$.

\textbf{3. Intégration par parties.} Posons $u(t) = t$ et $v'(t) = e^{-t/6}$, d'où $u'(t) = 1$ et $v(t) = -6e^{-t/6}$. Les fonctions $u'$ et $v'$ sont continues sur $[0\,;\,12]$, la formule s'applique :
\begin{align*}
I &= \Big[ -6t\,e^{-t/6} \Big]_0^{12} - \int_0^{12} 1 \times (-6)\,e^{-t/6}\,\mathrm{d}t \\
&= \left( -72e^{-2} - 0 \right) + 6 \int_0^{12} e^{-t/6}\,\mathrm{d}t \\
&= -72e^{-2} + 6\,\Big[ -6e^{-t/6} \Big]_0^{12} \\
&= -72e^{-2} + 6\left(-6e^{-2} + 6\right) = -72e^{-2} - 36e^{-2} + 36 = 36 - 108e^{-2}.
\end{align*}

\textbf{4. Volume exact.} Par linéarité de l'intégrale :
\[ V_{\text{eau}} = I + \int_0^{12} 1\,\mathrm{d}t = 36 - 108e^{-2} + 12 = 48 - 108e^{-2} \approx 33{,}38 \ \text{m}^3. \]

\textbf{5. Débit moyen.} Par définition de la valeur moyenne :
\[ \bar{d} = \dfrac{1}{12 - 0}\,V_{\text{eau}} = \dfrac{48 - 108e^{-2}}{12} = 4 - 9e^{-2} \approx 2{,}78 \ \text{m}^3/\text{h}. \]
Concrètement : tout s'est passé \emph{comme si} la pompe avait débité environ $2{,}78$ m$^3$/h de façon constante pendant 12 h.

\textbf{6. Rectangles.} Avec $h = 2$ et les valeurs $d(0) = 1$ ; $d(2) \approx 2{,}43$ ; $d(4) \approx 3{,}05$ ; $d(6) \approx 3{,}21$ ; $d(8) \approx 3{,}11$ ; $d(10) \approx 2{,}89$ ; $d(12) \approx 2{,}62$ :
\[ R_g = 2\,(1 + 2{,}43 + 3{,}05 + 3{,}21 + 3{,}11 + 2{,}89) \approx 31{,}38, \]
\[ R_d = 2\,(2{,}43 + 3{,}05 + 3{,}21 + 3{,}11 + 2{,}89 + 2{,}62) \approx 34{,}62. \]
La valeur exacte $33{,}38$ est bien comprise entre $31{,}38$ et $34{,}62$ : l'ordre de grandeur est confirmé. Remarque : $d$ n'étant pas monotone sur $[0\,;\,12]$ (elle croît puis décroît), ces sommes ne constituent pas un encadrement garanti ; la coïncidence ici est heureuse et sera discutée en classe. Pour obtenir un encadrement rigoureux, on appliquerait la méthode séparément sur $[0\,;\,6]$ (où $d$ est croissante) et sur $[6\,;\,12]$ (où elle est décroissante), grâce à la relation de Chasles.

\textbf{7. Capacité du réservoir.} Comme $\big(2\sqrt{x}\big)^2 = 4x$ :
\[ V_{\text{rés}} = \pi \int_0^{4} 4x\,\mathrm{d}x = \pi \Big[ 2x^2 \Big]_0^{4} = \pi(32 - 0) = 32\pi \approx 100{,}53 \ \text{m}^3. \]
Comme $33{,}38 < 100{,}53$, \textbf{toute l'eau collectée tient largement dans le réservoir} : il n'est rempli qu'à environ un tiers de sa capacité.
\end{exoresolu}

\begin{popfigure}
\begin{center}
\begin{tikzpicture}
\begin{axis}[
  width=0.9\linewidth, height=6cm,
  axis lines=middle,
  xlabel={$t$ (h)}, ylabel={$d(t)$ (m$^3$/h)},
  xmin=-0.5, xmax=13, ymin=0, ymax=3.8,
  xtick={0,2,4,6,8,10,12}, ytick={1,2,3},
  domain=0:12, samples=200, clip=false
]
\addplot[popcurve] {x*exp(-x/6)+1};
\addplot[poporange, dashed] coordinates {(6,0) (6,3.207)};
\node[popdark, anchor=south] at (axis cs:6,3.25) {\small max : $d(6) \approx 3{,}21$};
\node[popmuted, anchor=west] at (axis cs:0.2,0.85) {\small $d(0)=1$};
\node[popmuted, anchor=north west] at (axis cs:11.6,2.5) {\small $d(12)\approx 2{,}62$};
\end{axis}
\end{tikzpicture}
\end{center}
\legende{Courbe du débit $d(t) = t\,e^{-t/6} + 1$ sur $[0\,;\,12]$ : croissance jusqu'à midi ($t = 6$), puis décroissance. L'aire sous la courbe représente le volume d'eau collecté.}
\end{popfigure}

\courssub{Bilan de l'activité}

Cette activité a permis de mettre en évidence plusieurs idées fondamentales de la leçon. D'abord, l'\cle{intégrale} n'est pas qu'un objet abstrait : elle mesure une \textbf{quantité accumulée} (ici un volume d'eau) à partir d'une grandeur instantanée (un débit), ce qui justifie son interprétation comme aire sous la courbe pour une fonction continue et positive. Ensuite, tu as constaté que le calcul effectif d'une intégrale repose sur des \textbf{techniques complémentaires} : lecture directe d'une primitive pour les termes simples, intégration par parties pour le produit $t\,e^{-t/6}$. La \cle{valeur moyenne} a donné un sens concret au quotient $\dfrac{1}{b-a}\displaystyle\int_a^b f$ : un débit constant équivalent. La méthode des \cle{rectangles} a joué son rôle historique et pratique : approcher et \textbf{contrôler} un résultat, tout en rappelant que les encadrements garantis exigent la monotonie de la fonction — à défaut, on découpe l'intervalle grâce à la relation de Chasles. Enfin, le calcul du \cle{volume de révolution} a montré comment l'intégrale permet de passer d'un profil plan ($y = 2\sqrt{x}$) à une capacité réelle, et de trancher une question de décision : le réservoir de la coopérative de Bafoussam est correctement dimensionné. Tu es maintenant prêt à traiter les exercices du Baccalauréat mobilisant ces compétences, seules ou combinées.

\newpage

\coursec{Méthodes et exercices résolus}

\courssub{Calculer une intégrale par lecture directe d'une primitive}

\begin{methode}[Lecture directe d'une primitive]
Pour calculer $I=\displaystyle\int_a^b f(x)\,\mathrm{d}x$ lorsque $f$ est une forme usuelle :
\begin{enumerate}
\item \textbf{Identifier la forme} de $f$ : fonction usuelle ($x^n$, $\mathrm{e}^x$, $\cos$, $\sin$, $\dfrac{1}{x}$...) ou forme composée $u'\times(\text{formule en }u)$.
\item \textbf{Écrire une primitive} $F$ de $f$ sur $[a\,;b]$, en vérifiant mentalement que $F'=f$.
\item \textbf{Appliquer le théorème fondamental} : $I=\big[F(x)\big]_a^b=F(b)-F(a)$.
\item \textbf{Conclure} par une valeur exacte, éventuellement suivie d'une valeur approchée.
\end{enumerate}
Formules composées à connaître par cœur (avec $u$ dérivable, une primitive de chaque fonction de la colonne de gauche est donnée à droite) :
\[ u'\mathrm{e}^{u}\ \xrightarrow{\ \text{primitive}\ }\ \mathrm{e}^{u},\qquad u'u^n\ \xrightarrow{\ \text{primitive}\ }\ \frac{u^{n+1}}{n+1}\ (n\neq -1),\qquad \frac{u'}{u}\ \xrightarrow{\ \text{primitive}\ }\ \ln|u|, \]
\[ u'\cos u\ \xrightarrow{\ \text{primitive}\ }\ \sin u,\qquad u'\sin u\ \xrightarrow{\ \text{primitive}\ }\ -\cos u. \]
\end{methode}

\begin{exoresolu}[Intégrale par lecture directe (type Bac)]
Calculer les intégrales suivantes :
\[ I=\int_0^1 \left(3x^2-4x+1\right)\mathrm{d}x \qquad ; \qquad J=\int_0^{\ln 2} x\,\mathrm{e}^{x^2}\,\mathrm{d}x. \]
\tcblower
\textbf{Calcul de $I$.} La fonction $f:x\mapsto 3x^2-4x+1$ est un polynôme, donc continue sur $[0\,;1]$ : l'intégrale existe. Une primitive de $x\mapsto x^n$ est $x\mapsto \dfrac{x^{n+1}}{n+1}$, donc une primitive de $f$ est
\[ F(x)=x^3-2x^2+x. \]
Par le théorème fondamental du calcul intégral :
\begin{align*}
I &= \big[F(x)\big]_0^1 = F(1)-F(0) \\
&= (1-2+1)-(0) = 0.
\end{align*}
\textbf{Conclusion :} $I=0$. (Ce résultat nul n'est pas choquant : $f$ change de signe sur $[0\,;1]$ — en effet $f(x)=(3x-1)(x-1)$ — les aires se compensent.)

\medskip
\textbf{Calcul de $J$.} On reconnaît la forme $u'\,\mathrm{e}^{u}$ avec $u(x)=x^2$, donc $u'(x)=2x$. Or $x\,\mathrm{e}^{x^2}=\dfrac{1}{2}\times 2x\,\mathrm{e}^{x^2}=\dfrac{1}{2}u'(x)\mathrm{e}^{u(x)}$. Une primitive est donc
\[ G(x)=\frac{1}{2}\mathrm{e}^{x^2}. \]
La fonction $x\mapsto x\mathrm{e}^{x^2}$ est continue sur $\mathbb{R}$, donc sur $[0\,;\ln 2]$, et :
\begin{align*}
J &= \Big[\frac{1}{2}\mathrm{e}^{x^2}\Big]_0^{\ln 2} = \frac{1}{2}\mathrm{e}^{(\ln 2)^2}-\frac{1}{2}\mathrm{e}^{0} \\
&= \frac{1}{2}\left(\mathrm{e}^{(\ln 2)^2}-1\right).
\end{align*}
\textbf{Conclusion :} $J=\dfrac{1}{2}\left(\mathrm{e}^{(\ln 2)^2}-1\right)\approx 0{,}31$.
\end{exoresolu}

\begin{attention}[Ne pas confondre $(\ln 2)^2$ et $\ln 4$]
Dans le calcul de $J$ ci-dessus, la borne supérieure est $\ln 2$ et l'exposant final est $(\ln 2)^2$, carré du nombre $\ln 2$ : on a $(\ln 2)^2\approx 0{,}48$. Il ne faut surtout pas écrire $\mathrm{e}^{(\ln 2)^2}=2$ : la propriété $\mathrm{e}^{\ln a}=a$ ne s'applique que lorsque l'exposant est exactement un logarithme, pas son carré. De même, $(\ln 2)^2\neq \ln(2^2)=\ln 4$.
\end{attention}

\courssub{Calculer une intégrale par intégration par parties}

\begin{methode}[Intégration par parties (IPP)]
L'IPP s'utilise pour intégrer un \textbf{produit} de deux fonctions de natures différentes (polynôme $\times$ exponentielle, polynôme $\times$ cosinus, $\ln \times$ polynôme...). Au Bac, l'énoncé fournit toujours le choix de $u$ et $v'$.
\begin{enumerate}
\item \textbf{Poser} $u(x)$ et $v'(x)$ comme indiqué par l'énoncé ; en déduire $u'(x)$ et une primitive $v(x)$.
\item \textbf{Justifier} que $u$ et $v$ sont de classe $\mathcal{C}^1$ (dérivables à dérivée continue) sur $[a\,;b]$.
\item \textbf{Appliquer la formule} :
\[ \int_a^b u(x)\,v'(x)\,\mathrm{d}x = \big[u(x)v(x)\big]_a^b - \int_a^b u'(x)\,v(x)\,\mathrm{d}x. \]
\item \textbf{Calculer} le crochet et la nouvelle intégrale (plus simple que la première), puis conclure.
\end{enumerate}
\emph{Réflexe :} on dérive la fonction qui se simplifie par dérivation (polynôme, $\ln$) et on primitive l'autre.
\end{methode}

\begin{exoresolu}[Intégration par parties (type Bac)]
On considère l'intégrale $K=\displaystyle\int_0^1 x\,\mathrm{e}^{x}\,\mathrm{d}x$.
\begin{enumerate}
\item En posant $u(x)=x$ et $v'(x)=\mathrm{e}^x$, calculer $K$ par intégration par parties.
\item En déduire la valeur de $L=\displaystyle\int_0^1 (x+2)\mathrm{e}^x\,\mathrm{d}x$.
\end{enumerate}
\tcblower
\textbf{1.} Posons, comme l'indique l'énoncé :
\[ u(x)=x \quad\text{donc}\quad u'(x)=1 \qquad ; \qquad v'(x)=\mathrm{e}^x \quad\text{donc}\quad v(x)=\mathrm{e}^x. \]
Les fonctions $u$ et $v$ sont dérivables sur $[0\,;1]$ et leurs dérivées $u'$ et $v'$ sont continues sur $[0\,;1]$ : la formule d'intégration par parties s'applique :
\begin{align*}
K &= \big[x\,\mathrm{e}^x\big]_0^1 - \int_0^1 1\times \mathrm{e}^x\,\mathrm{d}x \\
&= \left(1\times\mathrm{e}^1 - 0\times\mathrm{e}^0\right) - \big[\mathrm{e}^x\big]_0^1 \\
&= \mathrm{e} - (\mathrm{e}-1) = 1.
\end{align*}
\textbf{Conclusion :} $K=1$.

\medskip
\textbf{2.} Par linéarité de l'intégrale :
\begin{align*}
L &= \int_0^1 x\,\mathrm{e}^x\,\mathrm{d}x + 2\int_0^1 \mathrm{e}^x\,\mathrm{d}x = K + 2\big[\mathrm{e}^x\big]_0^1 \\
&= 1 + 2(\mathrm{e}-1) = 2\mathrm{e}-1.
\end{align*}
\textbf{Conclusion :} $L=2\mathrm{e}-1\approx 4{,}44$.
\end{exoresolu}

\courssub{Calculer une intégrale par changement de variable affine}

\begin{methode}[Changement de variable affine]
Pour calculer $\displaystyle\int_a^b f(x)\,\mathrm{d}x$ en posant $t=\alpha x+\beta$ (avec $\alpha\neq 0$) :
\begin{enumerate}
\item \textbf{Poser} le changement de variable et calculer la différentielle : si $t=\alpha x+\beta$, alors $\mathrm{d}t=\alpha\,\mathrm{d}x$, donc $\mathrm{d}x=\dfrac{\mathrm{d}t}{\alpha}$.
\item \textbf{Transformer les bornes} : quand $x=a$, $t=\alpha a+\beta$ ; quand $x=b$, $t=\alpha b+\beta$. \emph{On ne garde jamais les anciennes bornes avec la nouvelle variable.}
\item \textbf{Exprimer l'intégrale} entièrement en fonction de $t$.
\item \textbf{Calculer} l'intégrale obtenue et conclure.
\end{enumerate}
Cas fréquent au Bac : si $F$ est une primitive de $f$, alors
\[ \int_a^b f(\alpha x+\beta)\,\mathrm{d}x=\frac{1}{\alpha}\big[F(\alpha x+\beta)\big]_a^b, \]
c'est-à-dire \og primitive de la composée $=\dfrac{1}{\alpha}F(\alpha x+\beta)$ \fg.
\end{methode}

\begin{exoresolu}[Changement de variable affine (type Bac)]
Calculer $M=\displaystyle\int_0^{\frac{\pi}{4}} \cos\left(2x+\frac{\pi}{6}\right)\mathrm{d}x$ en posant $t=2x+\dfrac{\pi}{6}$.
\tcblower
Posons $t=2x+\dfrac{\pi}{6}$. Alors $\mathrm{d}t=2\,\mathrm{d}x$, donc $\mathrm{d}x=\dfrac{1}{2}\mathrm{d}t$.

\textbf{Changement des bornes :}
\begin{itemize}
\item quand $x=0$, $t=2\times 0+\dfrac{\pi}{6}=\dfrac{\pi}{6}$ ;
\item quand $x=\dfrac{\pi}{4}$, $t=2\times\dfrac{\pi}{4}+\dfrac{\pi}{6}=\dfrac{\pi}{2}+\dfrac{\pi}{6}=\dfrac{3\pi+\pi}{6}=\dfrac{2\pi}{3}$.
\end{itemize}

\textbf{Transformation de l'intégrale :} la fonction $t\mapsto \cos t$ est continue sur $\mathbb{R}$, donc sur $\left[\dfrac{\pi}{6}\,;\dfrac{2\pi}{3}\right]$, et :
\begin{align*}
M &= \int_{\frac{\pi}{6}}^{\frac{2\pi}{3}} \cos t \times \frac{1}{2}\,\mathrm{d}t = \frac{1}{2}\big[\sin t\big]_{\frac{\pi}{6}}^{\frac{2\pi}{3}} \\
&= \frac{1}{2}\left(\sin\frac{2\pi}{3}-\sin\frac{\pi}{6}\right) = \frac{1}{2}\left(\frac{\sqrt{3}}{2}-\frac{1}{2}\right).
\end{align*}
\textbf{Conclusion :} $M=\dfrac{\sqrt{3}-1}{4}\approx 0{,}18$.

\medskip
\emph{Vérification rapide :} une primitive directe de $x\mapsto\cos\left(2x+\frac{\pi}{6}\right)$ est $x\mapsto \dfrac{1}{2}\sin\left(2x+\frac{\pi}{6}\right)$, ce qui redonne le même résultat.
\end{exoresolu}

\courssub{Calculer l'aire entre une courbe et l'axe des abscisses, ou entre deux courbes}

\begin{methode}[Calcul d'aires]
Le plan est muni d'un repère orthogonal $(O\,;\vec{\imath},\vec{\jmath})$ ; on note $\|\vec{\imath}\|=a$ cm et $\|\vec{\jmath}\|=b$ cm, donc $1$ u.a. $=a\times b$ cm$^2$.
\begin{enumerate}
\item \textbf{Étudier la position} de la courbe par rapport à l'axe des abscisses (ou des deux courbes l'une par rapport à l'autre) sur l'intervalle considéré : résoudre $f(x)\geq 0$ ou $f(x)\geq g(x)$.
\item \textbf{Écrire la bonne intégrale} :
\begin{itemize}
\item si $f\geq 0$ sur $[a\,;b]$ : $\mathcal{A}=\displaystyle\int_a^b f(x)\,\mathrm{d}x$ (en u.a.) ;
\item si $f\leq 0$ sur $[a\,;b]$ : $\mathcal{A}=-\displaystyle\int_a^b f(x)\,\mathrm{d}x=\displaystyle\int_a^b |f(x)|\,\mathrm{d}x$ ;
\item si $f\geq g$ sur $[a\,;b]$ : $\mathcal{A}=\displaystyle\int_a^b \big(f(x)-g(x)\big)\mathrm{d}x$.
\end{itemize}
Si le signe change, \textbf{découper} avec la relation de Chasles aux points où $f$ s'annule.
\item \textbf{Calculer} puis \textbf{convertir} en cm$^2$ si l'énoncé le demande : multiplier par l'unité d'aire.
\item \textbf{Conclure} : une aire est toujours un nombre \emph{positif}.
\end{enumerate}
\end{methode}

\begin{exoresolu}[Aire entre deux courbes (type Bac)]
Dans le plan muni d'un repère orthonormé $(O\,;\vec{\imath},\vec{\jmath})$ d'unité $\SI{2}{cm}$, on considère les fonctions $f(x)=x^2$ et $g(x)=x+2$, de courbes $\mathcal{C}_f$ et $\mathcal{C}_g$.
\begin{enumerate}
\item Déterminer les abscisses des points d'intersection de $\mathcal{C}_f$ et $\mathcal{C}_g$.
\item Calculer, en unités d'aire puis en cm$^2$, l'aire du domaine délimité par $\mathcal{C}_f$ et $\mathcal{C}_g$.
\end{enumerate}
\tcblower
\textbf{1.} On résout $f(x)=g(x)$ :
\[ x^2=x+2 \iff x^2-x-2=0. \]
Le discriminant est $\Delta=(-1)^2-4\times 1\times(-2)=1+8=9>0$, d'où
\[ x_1=\frac{1-3}{2}=-1 \qquad\text{et}\qquad x_2=\frac{1+3}{2}=2. \]
Les courbes se coupent aux points d'abscisses $-1$ et $2$.

\textbf{2.} Sur $[-1\,;2]$, étudions le signe de $g(x)-f(x)=-x^2+x+2=-(x^2-x-2)$. Le trinôme $x^2-x-2$ est négatif entre ses racines $-1$ et $2$ (coefficient de $x^2$ positif), donc $g(x)-f(x)\geq 0$ sur $[-1\,;2]$ : la courbe $\mathcal{C}_g$ est au-dessus de $\mathcal{C}_f$. L'aire cherchée est :
\begin{align*}
\mathcal{A} &= \int_{-1}^{2}\big(g(x)-f(x)\big)\mathrm{d}x = \int_{-1}^{2}\left(-x^2+x+2\right)\mathrm{d}x \\
&= \Big[-\frac{x^3}{3}+\frac{x^2}{2}+2x\Big]_{-1}^{2} \\
&= \left(-\frac{8}{3}+2+4\right)-\left(\frac{1}{3}+\frac{1}{2}-2\right) \\
&= \frac{10}{3}-\left(-\frac{7}{6}\right) = \frac{20}{6}+\frac{7}{6} = \frac{27}{6} = \frac{9}{2}.
\end{align*}
Donc $\mathcal{A}=\dfrac{9}{2}$ u.a. Comme le repère est orthonormé d'unité $\SI{2}{cm}$, l'unité d'aire vaut $2\times 2=4$ cm$^2$, d'où :
\[ \mathcal{A}=\frac{9}{2}\times 4 = 18\ \text{cm}^2. \]
\textbf{Conclusion :} le domaine délimité par les deux courbes a pour aire $\dfrac{9}{2}$ u.a., soit $\SI{18}{cm^2}$.
\end{exoresolu}

\begin{popfigure}
\begin{center}
\begin{tikzpicture}[scale=1.1]
\draw[->,popmuted] (-1.8,0) -- (2.8,0) node[below left] {$x$};
\draw[->,popmuted] (0,-0.6) -- (0,4.6) node[below left] {$y$};
\draw[popblue,thick,domain=-1.6:2.1,samples=100] plot (\x,{\x*\x}) node[right] {$\mathcal{C}_f$};
\draw[poporange,thick,domain=-1.7:2.4] plot (\x,{\x+2}) node[above right] {$\mathcal{C}_g$};
\fill[popblue!15] plot[domain=-1:2,samples=100] (\x,{\x+2}) -- plot[domain=2:-1,samples=100] (\x,{\x*\x}) -- cycle;
\fill[popdark] (-1,1) circle (1.5pt) node[above left] {$(-1\,;1)$};
\fill[popdark] (2,4) circle (1.5pt) node[above right] {$(2\,;4)$};
\draw (0,0) node[below left] {$O$};
\end{tikzpicture}
\end{center}
\legende{Le domaine délimité par $\mathcal{C}_f$ et $\mathcal{C}_g$ (en couleur) : sur $[-1\,;2]$, la droite $\mathcal{C}_g$ est au-dessus de la parabole $\mathcal{C}_f$.}
\end{popfigure}

\courssub{Calculer un volume de révolution}

\begin{methode}[Volume de révolution autour de l'axe des abscisses]
Soit $f$ continue sur $[a\,;b]$. Le solide engendré par la rotation de la courbe de $f$ autour de l'axe $(Ox)$ a pour volume :
\[ \mathcal{V}=\pi\int_a^b \big[f(x)\big]^2\,\mathrm{d}x \quad\text{(en unités de volume, u.v.)}. \]
\begin{enumerate}
\item \textbf{Vérifier} que $f$ est continue sur $[a\,;b]$ (la continuité garantit l'existence de l'intégrale ; le carré rend l'intégrande positive, donc le résultat est toujours positif).
\item \textbf{Écrire} $\mathcal{V}=\pi\displaystyle\int_a^b f(x)^2\,\mathrm{d}x$ — attention : c'est le \emph{carré} de $f$ qu'on intègre, multiplié par $\pi$.
\item \textbf{Développer} $f(x)^2$ si nécessaire, primitiver, calculer.
\item \textbf{Convertir} en cm$^3$ si besoin : $1$ u.v. $=(\text{unité de longueur})^3$.
\end{enumerate}
\emph{Image mentale :} le solide est un empilement de disques de rayon $|f(x)|$, donc d'aire $\pi f(x)^2$.
\end{methode}

\begin{exoresolu}[Volume de révolution (type Bac)]
Une coopérative de potiers de la région de l'Ouest fabrique un vase dont le profil est modélisé, dans un repère orthonormé d'unité $\SI{1}{dm}$, par la courbe de la fonction $f(x)=\sqrt{x+1}$ pour $x\in[0\,;3]$, en rotation autour de l'axe des abscisses.
\begin{enumerate}
\item Calculer le volume $\mathcal{V}$ du vase en unités de volume.
\item En déduire la contenance du vase en litres.
\end{enumerate}
\tcblower
\textbf{1.} La fonction $f:x\mapsto\sqrt{x+1}$ est continue et positive sur $[0\,;3]$ (car $x+1>0$). Le volume du solide de révolution engendré par la rotation de sa courbe autour de $(Ox)$ est :
\begin{align*}
\mathcal{V} &= \pi\int_0^3 \big[f(x)\big]^2\,\mathrm{d}x = \pi\int_0^3 (x+1)\,\mathrm{d}x \\
&= \pi\Big[\frac{x^2}{2}+x\Big]_0^3 = \pi\left(\frac{9}{2}+3-0\right) = \frac{15\pi}{2}.
\end{align*}
\textbf{Conclusion :} $\mathcal{V}=\dfrac{15\pi}{2}$ u.v. $\approx 23{,}56$ u.v.

\textbf{2.} L'unité de longueur est $\SI{1}{dm}$, donc $1$ u.v. $=1$ dm$^3=1$ litre. La contenance du vase est donc :
\[ \mathcal{V}=\frac{15\pi}{2}\ \text{litres}\approx 23{,}6\ \text{litres}. \]
\textbf{Conclusion :} le vase peut contenir environ $\SI{23,6}{L}$.
\end{exoresolu}

\courssub{Encadrer une intégrale par la méthode des rectangles}

\begin{methode}[Méthode des rectangles]
Pour encadrer $I=\displaystyle\int_a^b f(x)\,\mathrm{d}x$ où $f$ est continue, positive et \textbf{monotone} sur $[a\,;b]$ :
\begin{enumerate}
\item \textbf{Partager} $[a\,;b]$ en $n$ intervalles de même amplitude $h=\dfrac{b-a}{n}$ ; les points de la subdivision sont $x_k=a+kh$.
\item \textbf{Construire} sur chaque intervalle $[x_k\,;x_{k+1}]$ le rectangle inférieur (hauteur = plus petite valeur de $f$) et le rectangle supérieur (hauteur = plus grande valeur de $f$).
\item Si $f$ est \textbf{croissante} :
\[ h\sum_{k=0}^{n-1} f(x_k)\ \leq\ \int_a^b f(x)\,\mathrm{d}x\ \leq\ h\sum_{k=0}^{n-1} f(x_{k+1}). \]
Si $f$ est \textbf{décroissante}, l'encadrement est inversé :
\[ h\sum_{k=0}^{n-1} f(x_{k+1})\ \leq\ \int_a^b f(x)\,\mathrm{d}x\ \leq\ h\sum_{k=0}^{n-1} f(x_k). \]
\item \textbf{Calculer} les deux sommes et conclure par un encadrement de $I$. L'amplitude de l'encadrement vaut $h\,|f(b)-f(a)|$ : elle diminue quand $n$ augmente.
\end{enumerate}
\end{methode}

\begin{popfigure}
\begin{center}
\begin{tikzpicture}[scale=3.2]
\draw[->,popmuted] (0.8,0) -- (2.3,0) node[below left] {$x$};
\draw[->,popmuted] (1,0) -- (1,1.25) node[below left] {$y$};
\foreach \x/\h in {1/1, 1.25/0.8, 1.5/0.6667, 1.75/0.5714}{
  \draw[fill=poporange!25,draw=poporange] (\x,0) rectangle (\x+0.25,\h);
}
\foreach \x/\h in {1.25/1, 1.5/0.8, 1.75/0.6667, 2/0.5714}{
  \draw[draw=popblue,dashed] (\x-0.25,\h) rectangle (\x,0);
}
\draw[popblue,thick,domain=0.95:2.15,samples=100] plot (\x,{1/\x}) node[right] {$y=\dfrac{1}{x}$};
\foreach \x in {1,1.25,1.5,1.75,2}{\draw (\x,0) node[below] {\scriptsize $\x$};}
\end{tikzpicture}
\end{center}
\legende{Méthode des rectangles pour $f(x)=\dfrac{1}{x}$ sur $[1\,;2]$ avec $n=4$ : les rectangles orange (hauteurs à droite) sont sous la courbe, les rectangles pointillés bleus (hauteurs à gauche) la recouvrent.}
\end{popfigure}

\begin{exoresolu}[Méthode des rectangles (type Bac)]
Soit $f(x)=\dfrac{1}{x}$ sur $[1\,;2]$. On partage $[1\,;2]$ en $4$ intervalles de même amplitude.
\begin{enumerate}
\item Justifier que la méthode des rectangles permet d'encadrer $I=\displaystyle\int_1^2 \dfrac{\mathrm{d}x}{x}$.
\item Donner un encadrement de $I$ à l'aide de cette subdivision, puis en déduire un encadrement de $\ln 2$.
\end{enumerate}
\tcblower
\textbf{1.} La fonction $f:x\mapsto\dfrac{1}{x}$ est continue, positive et strictement décroissante sur $[1\,;2]$ (sa dérivée $f'(x)=-\dfrac{1}{x^2}$ est strictement négative). La méthode des rectangles s'applique donc.

\textbf{2.} L'amplitude commune est $h=\dfrac{2-1}{4}=0{,}25$. La subdivision est :
\[ x_0=1,\quad x_1=1{,}25,\quad x_2=1{,}5,\quad x_3=1{,}75,\quad x_4=2. \]
Valeurs de $f$ :
\[ f(1)=1,\quad f(1{,}25)=0{,}8,\quad f(1{,}5)=\tfrac{2}{3}\approx 0{,}6667,\quad f(1{,}75)=\tfrac{4}{7}\approx 0{,}5714,\quad f(2)=0{,}5. \]
Comme $f$ est \textbf{décroissante}, sur chaque intervalle le rectangle de hauteur $f(x_{k+1})$ est sous la courbe et celui de hauteur $f(x_k)$ est au-dessus :
\[ h\sum_{k=0}^{3} f(x_{k+1}) \leq I \leq h\sum_{k=0}^{3} f(x_k). \]
Somme inférieure :
\[ 0{,}25\times(0{,}8+0{,}6667+0{,}5714+0{,}5) = 0{,}25\times 2{,}5381 \approx 0{,}6345. \]
Somme supérieure :
\[ 0{,}25\times(1+0{,}8+0{,}6667+0{,}5714) = 0{,}25\times 3{,}0381 \approx 0{,}7595. \]
Donc $0{,}6345 \leq I \leq 0{,}7595$.

Or une primitive de $x\mapsto\dfrac{1}{x}$ sur $[1\,;2]$ est $\ln$, donc $I=\big[\ln x\big]_1^2=\ln 2-\ln 1=\ln 2$. On obtient :
\[ 0{,}634 \leq \ln 2 \leq 0{,}760. \]
L'amplitude de cet encadrement est $0{,}7595-0{,}6345=0{,}125$ ; on retrouve l'amplitude théorique $h\,|f(2)-f(1)|=0{,}25\times 0{,}5=0{,}125$.

\textbf{Conclusion :} $0{,}634\leq\ln 2\leq 0{,}760$, ce qui est cohérent avec la valeur exacte $\ln 2\approx 0{,}693$.
\end{exoresolu}

\begin{attention}[Les 5 erreurs qui coûtent des points au Bac]
\begin{enumerate}
\item \textbf{Oublier de vérifier la continuité ou les conditions d'application.} Avant toute intégration par parties, écris que $u$ et $v$ sont de classe $\mathcal{C}^1$ sur l'intervalle ; avant tout calcul d'aire ou de volume, vérifie la continuité de $f$. Une formule appliquée sans ses hypothèses est sanctionnée.
\item \textbf{Confondre aire et intégrale.} Une intégrale peut être négative ou nulle, une aire jamais. Si $f$ change de signe sur $[a\,;b]$, il faut découper avec Chasles et intégrer $|f|$ : $\mathcal{A}=\displaystyle\int_a^b |f(x)|\,\mathrm{d}x$, et non $\left|\displaystyle\int_a^b f(x)\,\mathrm{d}x\right|$ en général.
\item \textbf{Oublier le facteur $\dfrac{1}{\alpha}$ dans une composition affine.} Une primitive de $x\mapsto \mathrm{e}^{2x}$ est $\dfrac{1}{2}\mathrm{e}^{2x}$, pas $\mathrm{e}^{2x}$. De même, dans un changement de variable, ne jamais conserver les anciennes bornes avec la nouvelle variable.
\item \textbf{Mal rédiger le volume de révolution.} C'est $\mathcal{V}=\pi\displaystyle\int_a^b f(x)^2\,\mathrm{d}x$ : le $\pi$ et le \emph{carré} sont tous deux obligatoires. Écrire $\pi\int f$ ou $\int f^2$ sans $\pi$ fait perdre tous les points de la question.
\item \textbf{Négliger les unités et la conversion finale.} Si le repère a pour unité $\SI{2}{cm}$, l'unité d'aire vaut $4$ cm$^2$ (et non $2$) ; l'unité de volume vaut $8$ cm$^3$. Termine toujours par la conversion demandée et une phrase de conclusion avec l'unité.
\end{enumerate}
\end{attention}

\newpage

\coursec{Exercices}

\coursec{Calcul intégral}

\courssub{Intégrale d'une fonction continue : définition et premières propriétés}

\begin{definition}[Intégrale d'une fonction continue positive]
Soit $f$ une fonction continue et positive sur un intervalle $[a\,;b]$ ($a<b$). On note $\displaystyle\int_a^b f(x)\,\mathrm{d}x$ l'aire, en unités d'aire, du domaine planaire borné par la courbe représentative de $f$, l'axe des abscisses et les droites d'équations $x=a$ et $x=b$.
\end{definition}

\begin{definition}[Intégrale d'une fonction continue quelconque]
Soit $f$ une fonction continue sur $[a\,;b]$. L'intégrale $\displaystyle\int_a^b f(x)\,\mathrm{d}x$ est le nombre réel défini par :
\[
\int_a^b f(x)\,\mathrm{d}x = \int_a^b f^+(x)\,\mathrm{d}x - \int_a^b f^-(x)\,\mathrm{d}x
\]
où $f^+(x)=\max(f(x),0)$ et $f^-(x)=\max(-f(x),0)$. Géométriquement, c'est l'aire algébrique : l'aire au-dessus de l'axe est comptée positivement, celle en dessous négativement.
\end{definition}

\begin{aretenir}[Notations et vocabulaire]
\begin{itemize}
\item $f$ est la fonction \textbf{intégrée} (ou intégrande).
\item $a$ et $b$ sont les \textbf{bornes} de l'intégrale ($a$ borne inférieure, $b$ borne supérieure).
\item $x$ est la \textbf{variable d'intégration} (variable muette : $\int_a^b f(x)\,\mathrm{d}x = \int_a^b f(t)\,\mathrm{d}t$).
\item $\mathrm{d}x$ rappelle la variable et l'infinitésimal de longueur (largeur des rectangles de Riemann).
\item Si $a>b$, on pose par convention $\displaystyle\int_a^b f(x)\,\mathrm{d}x = -\int_b^a f(x)\,\mathrm{d}x$.
\item $\displaystyle\int_a^a f(x)\,\mathrm{d}x = 0$.
\end{itemize}
\end{aretenir}

\begin{propriete}[Propriétés fondamentales de l'intégrale]
Soient $f$ et $g$ continues sur $[a\,;b]$ et $\lambda \in \mathbb{R}$.
\begin{enumerate}
\item \textbf{Linéarité :} $\displaystyle\int_a^b \big(\lambda f(x) + g(x)\big)\,\mathrm{d}x = \lambda\int_a^b f(x)\,\mathrm{d}x + \int_a^b g(x)\,\mathrm{d}x$.
\item \textbf{Positivité :} Si $f(x) \ge 0$ sur $[a\,;b]$, alors $\displaystyle\int_a^b f(x)\,\mathrm{d}x \ge 0$.
\item \textbf{Relation de Chasles :} Pour tout $c \in [a\,;b]$, $\displaystyle\int_a^b f(x)\,\mathrm{d}x = \int_a^c f(x)\,\mathrm{d}x + \int_c^b f(x)\,\mathrm{d}x$. (Valide pour tout ordre de $a,b,c$ avec la convention sur les bornes).
\item \textbf{Inégalité de la moyenne :} Si $m \le f(x) \le M$ sur $[a\,;b]$, alors $m(b-a) \le \displaystyle\int_a^b f(x)\,\mathrm{d}x \le M(b-a)$.
\end{enumerate}
\end{propriete}

\begin{exemplebox}[Calcul direct d'une intégrale par les propriétés]
Calculer $\displaystyle\int_0^2 (3x^2 - 4x + 1)\,\mathrm{d}x$.
\par On utilise la linéarité et les intégrales élémentaires ($\int x^n = \frac{x^{n+1}}{n+1}$) :
\[
\int_0^2 (3x^2 - 4x + 1)\,\mathrm{d}x = 3\int_0^2 x^2\,\mathrm{d}x - 4\int_0^2 x\,\mathrm{d}x + \int_0^2 1\,\mathrm{d}x = 3\left[\frac{x^3}{3}\right]_0^2 - 4\left[\frac{x^2}{2}\right]_0^2 + \big[x\big]_0^2 = 8 - 8 + 2 = 2.
\]
\end{exemplebox}

\courssub{Fonction définie par une intégrale et valeur moyenne}

\begin{definition}[Fonction intégrale à borne variable]
Soit $f$ continue sur un intervalle $I$ et $a \in I$. La fonction $F$ définie sur $I$ par $F(x) = \displaystyle\int_a^x f(t)\,\mathrm{d}t$ est appelée \textbf{fonction intégrale de $f$ à borne variable $x$ et origine $a$}.
\end{definition}

\begin{propriete}[Sens de variation de $F$]
Si $f$ est continue sur $I$ :
\begin{itemize}
\item Si $f(x) \ge 0$ sur $I$, alors $F$ est croissante sur $I$.
\item Si $f(x) \le 0$ sur $I$, alors $F$ est décroissante sur $I$.
\item Plus généralement, $F$ a le même sens de variation que $f$ sur les intervalles où $f$ garde un signe constant.
\end{itemize}
\emph{Remarque :} On n'utilise pas la notion de primitive pour justifier cela, mais la positivité de l'intégrale (aire algébrique) et la relation de Chasles : pour $x_1 < x_2$, $F(x_2)-F(x_1) = \int_{x_1}^{x_2} f(t)\,\mathrm{d}t$.
\end{propriete}

\begin{definition}[Valeur moyenne d'une fonction]
La \textbf{valeur moyenne} de $f$ (continue) sur $[a\,;b]$ ($a<b$) est le nombre réel :
\[
m_f = \frac{1}{b-a}\int_a^b f(x)\,\mathrm{d}x.
\]
Géométriquement, c'est la hauteur du rectangle de base $[a\,;b]$ ayant la même aire (algébrique) que le domaine sous la courbe.
\end{definition}

\begin{exemplebox}[Valeur moyenne]
Soit $f(x)=x^2$ sur $[0\,;2]$. $m_f = \frac{1}{2}\int_0^2 x^2\,\mathrm{d}x = \frac{1}{2}\left[\frac{x^3}{3}\right]_0^2 = \frac{4}{3}$.
Le rectangle de base $[0,2]$ et de hauteur $4/3$ a une aire de $8/3$, égale à $\int_0^2 x^2\,\mathrm{d}x$.
\end{exemplebox}

\courssub{Techniques de calcul d'intégrales}

\begin{methode}[Calcul par lecture directe d'une primitive]
Une \textbf{primitive} de $f$ sur $I$ est une fonction $F$ dérivable sur $I$ telle que $F'=f$. Si $F$ est une primitive de $f$ sur $[a\,;b]$, alors (Théorème fondamental de l'analyse) :
\[
\int_a^b f(x)\,\mathrm{d}x = \big[F(x)\big]_a^b = F(b) - F(a).
\]
\begin{center}
\begin{tabularx}{\linewidth}{|l|X|l|}
\hline
\rowcolor{popblueL}
$f(x)$ & Une primitive $F(x)$ sur l'intervalle convenable & Condition \\
\hline
$x^\alpha$ ($\alpha \neq -1$) & $\dfrac{x^{\alpha+1}}{\alpha+1}$ & $x>0$ si $\alpha\notin\mathbb{Z}$ \\
$1/x$ & $\ln|x|$ & $x \neq 0$ \\
$\mathrm{e}^x$ & $\mathrm{e}^x$ & $\mathbb{R}$ \\
$\cos x$ & $\sin x$ & $\mathbb{R}$ \\
$\sin x$ & $-\cos x$ & $\mathbb{R}$ \\
$\dfrac{1}{\sqrt{1-x^2}}$ & $\arcsin x$ & $]-1,1[$ \\
$\dfrac{1}{1+x^2}$ & $\arctan x$ & $\mathbb{R}$ \\
\hline
\end{tabularx}
\end{center}
\end{methode}

\begin{methode}[Intégration par parties]
Soient $u$ et $v$ deux fonctions de classe $\mathcal{C}^1$ sur $[a\,;b]$. Alors :
\[
\int_a^b u(x)v'(x)\,\mathrm{d}x = \big[u(x)v(x)\big]_a^b - \int_a^b u'(x)v(x)\,\mathrm{d}x.
\]
\textbf{Stratégie de choix (LIATE) :} On pose $u$ = Logarithme, Inverse trigonométrique, Algébrique (polynôme), Trigonométrique, Exponentielle (dans cet ordre de priorité pour $u$). $v'$ est le reste.
\end{methode}

\begin{methode}[Changement de variable affine]
Pour calculer $\displaystyle\int_a^b f(ax+b)\,\mathrm{d}x$ ($a \neq 0$), on pose $t = ax+b$.
\begin{enumerate}
\item $\mathrm{d}t = a\,\mathrm{d}x \iff \mathrm{d}x = \frac{1}{a}\,\mathrm{d}t$.
\item Nouvelles bornes : $t_a = aa+b$, $t_b = ab+b$.
\item $\displaystyle\int_a^b f(ax+b)\,\mathrm{d}x = \frac{1}{a}\int_{t_a}^{t_b} f(t)\,\mathrm{d}t$.
\end{enumerate}
\emph{Attention :} Ne pas oublier le facteur $1/a$ et le changement de bornes !
\end{methode}

\begin{methode}[Formes remarquables : $u'\,u^n$ et $u'/u$]
Soit $u$ une fonction de classe $\mathcal{C}^1$ et strictement positive (pour $u'/u$) sur $[a\,;b]$.
\begin{itemize}
\item $\displaystyle\int_a^b u'(x)\big(u(x)\big)^n\,\mathrm{d}x = \left[\frac{\big(u(x)\big)^{n+1}}{n+1}\right]_a^b$ ($n \neq -1$).
\item $\displaystyle\int_a^b \frac{u'(x)}{u(x)}\,\mathrm{d}x = \big[\ln u(x)\big]_a^b$.
\end{itemize}
\end{methode}

\begin{attention}[Pièges fréquents]
\begin{itemize}
\item Oublier de changer les bornes lors d'un changement de variable (ou oublier de revenir à la variable initiale avant d'appliquer les bornes initiales).
\item Oublier le facteur $1/a$ dans le changement de variable affine $t=ax+b$.
\item Confondre $\int f(g(x))\,\mathrm{d}x$ et $\int f(g(x))g'(x)\,\mathrm{d}x$ (seule la seconde se fait par changement de variable simple).
\item Écrire $\int_a^b f(x)g(x)\,\mathrm{d}x = \int_a^b f(x)\,\mathrm{d}x \times \int_a^b g(x)\,\mathrm{d}x$ (FAUX).
\item Ne pas vérifier la continuité de $f$ sur $[a\,;b]$ avant d'appliquer le théorème fondamental.
\end{itemize}
\end{attention}

\courssub{Intégrale d'une fonction valeur absolue}

\begin{methode}[Calcul de $\int_a^b |f(x)|\,\mathrm{d}x$]
\begin{enumerate}
\item Résoudre $f(x)=0$ sur $[a\,;b]$ pour trouver les zéros $x_1, x_2, \dots$
\item Découper $[a\,;b]$ en sous-intervalles où $f$ a un signe constant (tableau de signes).
\item Sur chaque sous-intervalle, remplacer $|f(x)|$ par $f(x)$ si $f\ge 0$ ou par $-f(x)$ si $f\le 0$.
\item Appliquer la relation de Chasles et calculer chaque intégrale.
\end{enumerate}
\end{methode}

\begin{exemplebox}[Intégrale avec valeur absolue]
Calculer $\displaystyle\int_{-1}^{3} |x-1|\,\mathrm{d}x$.
\par $x-1=0 \Leftrightarrow x=1$. Sur $[-1,1]$, $x-1\le 0 \Rightarrow |x-1|=1-x$. Sur $[1,3]$, $x-1\ge 0 \Rightarrow |x-1|=x-1$.
\[
\int_{-1}^{3} |x-1|\,\mathrm{d}x = \int_{-1}^{1} (1-x)\,\mathrm{d}x + \int_{1}^{3} (x-1)\,\mathrm{d}x = \big[x-\frac{x^2}{2}\big]_{-1}^{1} + \big[\frac{x^2}{2}-x\big]_{1}^{3} = 2 + 2 = 4.
\]
\end{exemplebox}

\courssub{Applications géométriques : Aires et Volumes de révolution}

\begin{propriete}[Aire entre une courbe et l'axe des abscisses]
Soit $f$ continue sur $[a\,;b]$. L'aire $\mathcal{A}$ (en unités d'aire) du domaine borné par la courbe $y=f(x)$, l'axe des abscisses et les droites $x=a$, $x=b$ est :
\[
\mathcal{A} = \int_a^b |f(x)|\,\mathrm{d}x.
\]
Si $f$ garde un signe constant, $\mathcal{A} = \left|\int_a^b f(x)\,\mathrm{d}x\right|$.
\end{propriete}

\begin{propriete}[Aire entre deux courbes]
Soient $f$ et $g$ continues sur $[a\,;b]$ avec $f(x) \le g(x)$. L'aire $\mathcal{A}$ du domaine borné par les courbes $y=f(x)$, $y=g(x)$ et les droites $x=a$, $x=b$ est :
\[
\mathcal{A} = \int_a^b \big(g(x)-f(x)\big)\,\mathrm{d}x.
\]
Si les courbes se coupent, on découpe l'intervalle aux points d'intersection.
\end{propriete}

\begin{propriete}[Volume de révolution autour de $(Ox)$]
Soit $f$ continue et positive sur $[a\,;b]$. Le volume $V$ (en unités de volume) du solide engendré par la rotation de la courbe $y=f(x)$ autour de l'axe $(Ox)$ est :
\[
V = \pi \int_a^b \big(f(x)\big)^2\,\mathrm{d}x.
\]
(Formule des disques : aire du disque $\pi R^2$ avec $R=f(x)$).
\end{propriete}

\begin{exemplebox}[Volume de révolution]
$f(x)=\sqrt{x}+1$ sur $[0,4]$. $V = \pi\int_0^4 (\sqrt{x}+1)^2\,\mathrm{d}x = \pi\int_0^4 (x+2\sqrt{x}+1)\,\mathrm{d}x = \pi\left[\frac{x^2}{2}+\frac{4}{3}x^{3/2}+x\right]_0^4 = \pi(8+\frac{32}{3}+4) = \frac{68\pi}{3}$ u.v.
\end{exemplebox}

\courssub{Approximation d'une intégrale : Méthode des rectangles}

\begin{definition}[Sommes de rectangles]
Soit $f$ continue et positive sur $[a\,;b]$. On partage $[a\,;b]$ en $n$ sous-intervalles de même amplitude $h = \frac{b-a}{n}$. Les points de subdivision sont $x_k = a+kh$ ($k=0,\dots,n$).
\begin{itemize}
\item \textbf{Rectangles à gauche :} $R_g = h\sum_{k=0}^{n-1} f(x_k)$. Hauteur = $f(x_k)$ sur $[x_k, x_{k+1}]$.
\item \textbf{Rectangles à droite :} $R_d = h\sum_{k=1}^{n} f(x_k)$. Hauteur = $f(x_k)$ sur $[x_{k-1}, x_k]$.
\end{itemize}
\end{definition}

\begin{propriete}[Encadrement pour une fonction monotone]
Si $f$ est \textbf{croissante} sur $[a\,;b]$ : $R_g \le \int_a^b f \le R_d$.
Si $f$ est \textbf{décroissante} sur $[a\,;b]$ : $R_d \le \int_a^b f \le R_g$.
L'erreur commise est majorée par $(b-a)\times$ (variation totale de $f$ sur l'intervalle) $/ n$.
\end{propriete}

\begin{savaistu}[Contexte camerounais : Gestion de l'eau]
La société CAMWATER (ou une société de distribution régionale) utilise des capteurs de débit pour modéliser la consommation. L'intégrale du débit (en $\mathrm{m}^3/\mathrm{h}$) sur une période donne le volume total distribué (en $\mathrm{m}^3$). La méthode des rectangles permet d'estimer ce volume à partir de relevés horaires discrets, avant d'avoir un modèle analytique précis.
\end{savaistu}

\begin{popfigure}
\begin{tikzpicture}[scale=0.7, >=Stealth]
\draw[->] (-0.5,0) -- (5,0) node[right] {$x$};
\draw[->] (0,-0.5) -- (0,5) node[above] {$y$};
\draw[domain=0:4, smooth, variable=\x, popcurve, thick] plot ({\x}, {0.2*\x*\x+0.5});
\foreach \k in {0,1,2,3} {
    \draw[popblue, fill=popblue!20] (\k,0) rectangle (\k+1, {0.2*\k*\k+0.5});
    \draw[dashed, popmuted] (\k,0) -- (\k, {0.2*\k*\k+0.5});
}
\draw[dashed, popmuted] (4,0) -- (4, {0.2*4*4+0.5});
\node at (0.5, -0.3) {$x_0$}; \node at (1.5, -0.3) {$x_1$}; \node at (2.5, -0.3) {$x_2$}; \node at (3.5, -0.3) {$x_3$}; \node at (4.5, -0.3) {$x_4$};
\node[popblue] at (2, 4.5) {Rectangles à gauche ($R_g$)};
\end{tikzpicture}
\legende{Méthode des rectangles à gauche pour $f(x)=0.2x^2+0.5$ sur $[0,4]$ avec $n=4$.}
\end{popfigure}

\courssub{Je vérifie mes connaissances}

\begin{exercice}[QCM : les bases de l'intégrale]
Pour chaque question, une seule réponse est exacte. Recopier la lettre correspondante en justifiant brièvement.
\begin{enumerate}
\item Si $f$ est continue et positive sur $[a\,;b]$, alors $\displaystyle\int_a^b f(x)\,\mathrm{d}x$ représente :
\begin{enumerate}
\item[a)] la dérivée de $f$ en $b$ ;
\item[b)] l'aire, en unités d'aire, du domaine compris entre la courbe de $f$, l'axe des abscisses et les droites $x=a$ et $x=b$ ;
\item[c)] la longueur de la courbe de $f$ entre $a$ et $b$.
\end{enumerate}
\item $\displaystyle\int_1^3 (2x+1)\,\mathrm{d}x$ est égale à :
\begin{enumerate}
\item[a)] $10$ ; \qquad b)] $12$ ; \qquad c)] $8$.
\end{enumerate}
\item Pour toute fonction $f$ continue sur $[a\,;b]$ et tout $c\in[a\,;b]$, la relation de Chasles s'écrit :
\begin{enumerate}
\item[a)] $\displaystyle\int_a^b f = \int_a^c f + \int_c^b f$ ;
\item[b)] $\displaystyle\int_a^b f = \int_a^c f \times \int_c^b f$ ;
\item[c)] $\displaystyle\int_a^b f = \int_b^c f - \int_a^c f$.
\end{enumerate}
\item La valeur moyenne de la fonction $f$ continue sur $[a\,;b]$ est :
\begin{enumerate}
\item[a)] $\dfrac{f(a)+f(b)}{2}$ ;
\item[b)] $\dfrac{1}{b-a}\displaystyle\int_a^b f(x)\,\mathrm{d}x$ ;
\item[c)] $\displaystyle\int_a^b f(x)\,\mathrm{d}x$.
\end{enumerate}
\end{enumerate}
\end{exercice}

\begin{exercice}[Vrai ou faux, justifié]
Pour chacune des affirmations suivantes, dire si elle est vraie ou fausse en justifiant la réponse (un contre-exemple suffit pour infirmer).
\begin{enumerate}
\item Si $\displaystyle\int_a^b f(x)\,\mathrm{d}x = 0$, alors $f$ est la fonction nulle sur $[a\,;b]$.
\item Pour toutes fonctions $f$ et $g$ continues sur $[a\,;b]$ : $\displaystyle\int_a^b f(x)g(x)\,\mathrm{d}x = \left(\int_a^b f(x)\,\mathrm{d}x\right)\left(\int_a^b g(x)\,\mathrm{d}x\right)$.
\item Si pour tout $x\in[a\,;b]$, $m\le f(x)\le M$, alors $m(b-a)\le \displaystyle\int_a^b f(x)\,\mathrm{d}x \le M(b-a)$.
\item $\displaystyle\int_{-2}^{2} x^3\,\mathrm{d}x = 0$.
\end{enumerate}
\end{exercice}

\begin{exercice}[Définitions et notations]
\begin{enumerate}
\item Énoncer la définition de l'intégrale d'une fonction continue et positive sur $[a\,;b]$ en termes d'aire, et préciser ce qu'est une unité d'aire dans un repère orthogonal $(O\,;\vec{\imath},\vec{\jmath})$.
\item Donner la définition de la fonction $F : x\longmapsto \displaystyle\int_a^x f(t)\,\mathrm{d}t$ et préciser son sens de variation lorsque $f$ est positive sur $[a\,;b]$.
\item Rappeler la formule du volume du solide engendré par la rotation de la courbe de $f$ (continue et positive) autour de l'axe $(Ox)$ sur $[a\,;b]$.
\end{enumerate}
\end{exercice}

\begin{exercice}[Calculs directs]
Calculer les intégrales suivantes par lecture directe d'une primitive :
\begin{enumerate}
\item $A=\displaystyle\int_0^2 (3x^2-4x+1)\,\mathrm{d}x$ ;
\item $B=\displaystyle\int_1^{4} \dfrac{1}{\sqrt{x}}\,\mathrm{d}x$ ;
\item $C=\displaystyle\int_0^{\pi} \sin x\,\mathrm{d}x$ ;
\item $D=\displaystyle\int_0^{1} \mathrm{e}^{3x}\,\mathrm{d}x$.
\end{enumerate}
\end{exercice}

\courssub{Je m'entraîne}

\begin{exercice}[Linéarité, Chasles et valeur absolue]
\begin{enumerate}
\item Sachant que $\displaystyle\int_0^1 f(x)\,\mathrm{d}x = 5$ et $\displaystyle\int_0^1 g(x)\,\mathrm{d}x = -2$, calculer $\displaystyle\int_0^1 \big(3f(x)-4g(x)\big)\,\mathrm{d}x$.
\item Sachant que $\displaystyle\int_1^5 f(x)\,\mathrm{d}x = 7$ et $\displaystyle\int_3^5 f(x)\,\mathrm{d}x = 2$, calculer $\displaystyle\int_1^3 f(x)\,\mathrm{d}x$.
\item Calculer $\displaystyle\int_{-1}^{3} |x-1|\,\mathrm{d}x$ en utilisant la relation de Chasles (on pourra s'aider d'une figure).
\end{enumerate}
\end{exercice}

\begin{exercice}[Intégration par parties guidée]
On veut calculer $I=\displaystyle\int_0^1 x\,\mathrm{e}^{x}\,\mathrm{d}x$.
On pose $u(x)=x$ et $v'(x)=\mathrm{e}^{x}$.
\begin{enumerate}
\item Déterminer $u'(x)$ et une fonction $v(x)$.
\item En utilisant la formule $\displaystyle\int_a^b u\,v' = \big[u\,v\big]_a^b - \int_a^b u'\,v$, calculer la valeur exacte de $I$.
\item Par la même méthode, calculer $J=\displaystyle\int_1^{\mathrm{e}} \ln x\,\mathrm{d}x$ (on posera $u(x)=\ln x$ et $v'(x)=1$).
\end{enumerate}
\end{exercice}

\begin{exercice}[Changement de variable affine et forme $u'u^n$]
\begin{enumerate}
\item En effectuant le changement de variable $t=2x+1$, calculer $\displaystyle\int_0^{1} (2x+1)^3\,\mathrm{d}x$.
\item Calculer $\displaystyle\int_0^{\pi/2} \cos(3x)\,\mathrm{d}x$ à l'aide du changement de variable $t=3x$.
\item En reconnaissant une forme $u'\,u$, calculer $\displaystyle\int_1^{\mathrm{e}} \dfrac{\ln x}{x}\,\mathrm{d}x$.
\end{enumerate}
\end{exercice}

\begin{exercice}[Aire et méthode des rectangles]
Soit $f$ la fonction définie sur $[0\,;2]$ par $f(x)=x^2+1$, de courbe représentative $\mathcal{C}_f$ dans un repère orthonormé d'unité \SI{1}{cm}.
\begin{enumerate}
\item Calculer l'aire $\mathcal{A}$, en unités d'aire puis en $\mathrm{cm}^2$, du domaine délimité par $\mathcal{C}_f$, l'axe des abscisses et les droites $x=0$ et $x=2$.
\item On partage $[0\,;2]$ en $4$ intervalles de même amplitude. Calculer la somme des aires des rectangles à gauche $R_g$ et celle des rectangles à droite $R_d$.
\item Justifier que $R_g \le \mathcal{A} \le R_d$ et en déduire un encadrement de $\mathcal{A}$. Comparer avec la valeur exacte.
\end{enumerate}
\end{exercice}

\courssub{Je me prépare au Bac}

\begin{exercice}[Type Bac — Trois techniques de calcul d'intégrales]
Le sujet propose de calculer trois intégrales par trois méthodes différentes.
\begin{enumerate}
\item \textbf{Par intégration par parties.} On considère $I=\displaystyle\int_0^1 (2x+3)\,\mathrm{e}^{x}\,\mathrm{d}x$.\\
On pose $u(x)=2x+3$ et $v'(x)=\mathrm{e}^{x}$.
\begin{enumerate}
\item Déterminer $u'(x)$ et une fonction $v(x)$.
\item À l'aide de la formule d'intégration par parties rappelée par l'énoncé, montrer que $I=3\mathrm{e}-1$.
\end{enumerate}
\item \textbf{Par lecture de la forme $u'\,u$.} On considère $J=\displaystyle\int_1^{\mathrm{e}} \dfrac{\ln x}{x}\,\mathrm{d}x$.\\
En remarquant que $\dfrac{\ln x}{x}$ est de la forme $u'(x)\,u(x)$ avec une fonction $u$ à préciser, montrer que $J=\dfrac{1}{2}$.
\item \textbf{Par changement de variable affine.} On considère $K=\displaystyle\int_0^{\pi/4} \cos(2x)\,\mathrm{d}x$.\\
On pose $t=2x$.
\begin{enumerate}
\item Déterminer les nouvelles bornes et exprimer $\mathrm{d}x$ en fonction de $\mathrm{d}t$.
\item En déduire la valeur exacte de $K$.
\end{enumerate}
\item \textbf{Application.} Un artisan de Bafoussam fabrique des tuiles dont le profil, en décimètres, est modélisé sur $[0\,;\pi/4]$ par la courbe d'équation $y=\cos(2x)$ dans un repère orthonormé. Calculer, en $\mathrm{dm}^2$, l'aire de la face d'une tuile délimitée par cette courbe, l'axe des abscisses et les droites $x=0$ et $x=\dfrac{\pi}{4}$.
\end{enumerate}
\end{exercice}

\begin{exercice}[Type Bac — Fonction définie par une intégrale et méthode des rectangles]
Une société de distribution d'eau de la région du Littoral mesure le débit instantané (en $\mathrm{m}^3/\mathrm{h}$) d'une conduite pendant $6$ heures. Ce débit est modélisé par une fonction $f$ continue et positive sur $[0\,;6]$, dont on donne le tableau de valeurs :
\begin{center}
\begin{tabularx}{\linewidth}{|l|X|X|X|X|X|X|X|}
\hline
\rowcolor{popblueL}
$x$ (heures) & $0$ & $1$ & $2$ & $3$ & $4$ & $5$ & $6$ \\
\hline
$f(x)$ ($\mathrm{m}^3/\mathrm{h}$) & $2$ & $4$ & $5$ & $4{,}5$ & $3$ & $2$ & $1$ \\
\hline
\end{tabularx}
\end{center}
On admet que $f$ est croissante sur $[0\,;2]$ et décroissante sur $[2\,;6]$. Pour tout $x\in[0\,;6]$, on pose $F(x)=\displaystyle\int_0^x f(t)\,\mathrm{d}t$, qui représente le volume total d'eau écoulé (en $\mathrm{m}^3$) entre l'instant $0$ et l'instant $x$.
\begin{enumerate}
\item Justifier, sans calculer de primitive, que $F$ est croissante sur $[0\,;6]$.
\item Donner la valeur de $F(0)$. Que représente $F(6)$ pour la société de distribution ?
\item \textbf{Méthode des rectangles.} On partage $[0\,;6]$ en $6$ intervalles de même amplitude $h=1$.
\begin{enumerate}
\item Calculer la somme $R_g$ des aires des rectangles à gauche et la somme $R_d$ des aires des rectangles à droite sur $[0\,;6]$.
\item En déduire un encadrement de $F(6)$, puis une valeur approchée du volume d'eau distribué en $6$ heures à $0{,}5\ \mathrm{m}^3$ près.
\end{enumerate}
\item On admet désormais que sur $[0\,;2]$, $f(t)=\dfrac{3}{2}t+2$.
\begin{enumerate}
\item Calculer la valeur exacte de $F(2)$.
\item En utilisant l'inégalité de la moyenne sur $[2\,;6]$ avec les valeurs du tableau, donner un encadrement de $\displaystyle\int_2^6 f(t)\,\mathrm{d}t$, puis un nouvel encadrement de $F(6)$. Comparer avec le résultat de la question 3.
\end{enumerate}
\end{enumerate}
\end{exercice}

\begin{exercice}[Type Bac — Aire entre deux courbes, volume de révolution et valeur moyenne]
Le plan est muni d'un repère orthogonal $(O\,;\vec{\imath},\vec{\jmath})$ d'unités graphiques : \SI{2}{cm} sur chaque axe. On considère les fonctions $f$ et $g$ définies sur $\mathbb{R}$ par $f(x)=x^2$ et $g(x)=2x+3$, de courbes respectives $\mathcal{P}$ et $\mathcal{D}$.

\textbf{Partie A : aire entre deux courbes}
\begin{enumerate}
\item Résoudre dans $\mathbb{R}$ l'équation $f(x)=g(x)$ et en déduire les coordonnées des points d'intersection de $\mathcal{P}$ et $\mathcal{D}$.
\item Étudier le signe de $g(x)-f(x)$ sur $[-1\,;3]$ et en déduire la position relative des deux courbes sur cet intervalle.
\item Montrer que l'aire $\mathcal{A}$ du domaine compris entre $\mathcal{P}$ et $\mathcal{D}$ est égale, en unités d'aire, à $\displaystyle\int_{-1}^{3}\big(g(x)-f(x)\big)\,\mathrm{d}x$, puis calculer sa valeur exacte.
\item En déduire l'aire de ce domaine en $\mathrm{cm}^2$.
\end{enumerate}

\textbf{Partie B : volume de révolution}\\
Un potier de Foumban modélise un vase par la rotation, autour de l'axe $(Ox)$, de la courbe d'équation $y=\sqrt{x}+1$ sur $[0\,;4]$, l'unité de longueur étant \SI{5}{cm}.
\begin{enumerate}
\item Rappeler la formule donnant le volume $V$ du solide de révolution engendré, en unités de volume.
\item Montrer que $V=\pi\displaystyle\int_0^4 \big(x+2\sqrt{x}+1\big)\,\mathrm{d}x$, puis calculer la valeur exacte de $V$ en unités de volume.
\item En déduire le volume du vase en $\mathrm{cm}^3$, puis sa contenance en litres (on rappelle que $1\ \mathrm{L}=1000\ \mathrm{cm}^3$ ; on prendra $\pi\approx 3{,}1416$ et on arrondira au centilitre).
\end{enumerate}

\textbf{Partie C : valeur absolue et valeur moyenne}
\begin{enumerate}
\item En utilisant la relation de Chasles et la périodicité de $|\cos|$, montrer que $\displaystyle\int_0^{2\pi} |\cos x|\,\mathrm{d}x = 4\int_0^{\pi/2}\cos x\,\mathrm{d}x$, puis calculer sa valeur exacte.
\item En déduire la valeur moyenne de la fonction $x\mapsto|\cos x|$ sur $[0\,;2\pi]$. Interpréter graphiquement ce résultat à l'aide d'un rectangle d'aire équivalente.
\end{enumerate}
\end{exercice}



\newpage

\coursec{Corrigés des exercices}

\begin{corrige}[Exercice 1 — QCM : les bases de l'intégrale]
\begin{enumerate}
\item \textbf{Réponse b).} Par définition, pour une fonction continue et positive sur $[a\,;b]$, l'intégrale $\displaystyle\int_a^b f(x)\,\mathrm{d}x$ est l'aire, en unités d'aire, du domaine délimité par la courbe de $f$, l'axe des abscisses et les droites d'équations $x=a$ et $x=b$. La réponse a) est fausse (la dérivée n'a rien à voir avec une aire) et la réponse c) est fausse (la longueur d'une courbe se calcule par une autre formule, hors programme).
\item \textbf{Réponse a).} Une primitive de $x\mapsto 2x+1$ est $x\mapsto x^2+x$. Donc :
\[
\int_1^3 (2x+1)\,\mathrm{d}x = \big[x^2+x\big]_1^3 = (9+3)-(1+1) = 12-2 = 10.
\]
\item \textbf{Réponse a).} La relation de Chasles traduit l'\cle{additivité} de l'aire algébrique : $\displaystyle\int_a^b f = \int_a^c f + \int_c^b f$. La réponse b) confond avec un produit (faux en général) et la réponse c) inverse les bornes.
\item \textbf{Réponse b).} Par définition, la valeur moyenne de $f$ sur $[a\,;b]$ est $m_f = \dfrac{1}{b-a}\displaystyle\int_a^b f(x)\,\mathrm{d}x$. La réponse a) n'est exacte que pour les fonctions affines ; la réponse c) oublie le facteur $\frac{1}{b-a}$.
\end{enumerate}
\textbf{Point de vigilance :} au Bac, une réponse de QCM sans justification peut ne pas rapporter tous les points ; cite toujours la définition ou fais le calcul.
\end{corrige}

\begin{corrige}[Exercice 2 — Vrai ou faux, justifié]
\begin{enumerate}
\item \textbf{Faux.} Contre-exemple : $f(x)=x$ sur $[-1\,;1]$. On a $\displaystyle\int_{-1}^{1} x\,\mathrm{d}x = \left[\dfrac{x^2}{2}\right]_{-1}^{1} = \dfrac{1}{2}-\dfrac{1}{2} = 0$, alors que $f$ n'est pas la fonction nulle. L'intégrale mesure une aire \emph{algébrique} : les aires au-dessus et en dessous de l'axe se compensent.
\emph{Remarque :} l'affirmation deviendrait vraie si l'on ajoutait l'hypothèse « $f$ continue et positive ».
\item \textbf{Faux.} L'intégrale d'un produit n'est pas le produit des intégrales. Contre-exemple : $f(x)=g(x)=x$ sur $[0\,;1]$. D'une part $\displaystyle\int_0^1 x^2\,\mathrm{d}x = \left[\dfrac{x^3}{3}\right]_0^1 = \dfrac{1}{3}$ ; d'autre part $\displaystyle\left(\int_0^1 x\,\mathrm{d}x\right)^2 = \left(\dfrac{1}{2}\right)^2 = \dfrac{1}{4} \neq \dfrac{1}{3}$.
\item \textbf{Vrai.} C'est l'\cle{inégalité de la moyenne}. Démonstration : pour tout $x\in[a\,;b]$, $f(x)-m\ge 0$, donc par positivité $\displaystyle\int_a^b (f(x)-m)\,\mathrm{d}x \ge 0$, soit par linéarité $\displaystyle\int_a^b f \ge m(b-a)$. De même, $M-f(x)\ge 0$ donne $\displaystyle\int_a^b f \le M(b-a)$.
\item \textbf{Vrai.} Calcul direct : $\displaystyle\int_{-2}^{2} x^3\,\mathrm{d}x = \left[\dfrac{x^4}{4}\right]_{-2}^{2} = \dfrac{16}{4}-\dfrac{16}{4} = 0$. Justification géométrique : la fonction $x\mapsto x^3$ est impaire et l'intervalle $[-2\,;2]$ est centré en $0$, donc les aires algébriques se compensent.
\end{enumerate}
\end{corrige}

\begin{corrige}[Exercice 3 — Définitions et notations]
\begin{enumerate}
\item Soit $f$ continue et positive sur $[a\,;b]$. L'intégrale $\displaystyle\int_a^b f(x)\,\mathrm{d}x$ est l'aire, exprimée en unités d'aire, du domaine plan borné par la courbe représentative de $f$, l'axe des abscisses et les droites d'équations $x=a$ et $x=b$. Dans un repère orthogonal $(O\,;\vec{\imath},\vec{\jmath})$, l'\cle{unité d'aire} est l'aire du rectangle construit sur les vecteurs $\vec{\imath}$ et $\vec{\jmath}$ : $1\ \text{u.a.} = \|\vec{\imath}\|\times\|\vec{\jmath}\|$. (Dans un repère orthonormé d'unité \SI{1}{cm}, $1\ \text{u.a.} = 1\ \mathrm{cm}^2$.)
\item Pour $f$ continue sur un intervalle $I$ et $a\in I$, on définit sur $I$ la fonction $F : x\longmapsto \displaystyle\int_a^x f(t)\,\mathrm{d}t$, appelée fonction intégrale de $f$ d'origine $a$. Si $f$ est positive sur $I$, alors pour $x_1 < x_2$ dans $I$, la relation de Chasles donne :
\[
F(x_2)-F(x_1) = \int_{x_1}^{x_2} f(t)\,\mathrm{d}t \ge 0 \quad (\text{positivité de l'intégrale}),
\]
donc $F$ est \textbf{croissante} sur $I$.
\item Le volume du solide engendré par la rotation de la courbe de $f$ (continue et positive) autour de l'axe $(Ox)$ entre $a$ et $b$ est :
\[
V = \pi\int_a^b \big(f(x)\big)^2\,\mathrm{d}x \quad \text{(en unités de volume).}
\]
\end{enumerate}
\textbf{Point de vigilance :} dans la question 2, le programme exige la justification \emph{sans} utiliser la notion de primitive : on s'appuie uniquement sur Chasles et la positivité.
\end{corrige}

\begin{corrige}[Exercice 4 — Calculs directs par primitives]
\begin{enumerate}
\item Une primitive de $3x^2-4x+1$ est $x^3-2x^2+x$. Donc :
\[
A = \big[x^3-2x^2+x\big]_0^2 = (8-8+2)-0 = \boxed{2}.
\]
\item $\dfrac{1}{\sqrt{x}} = x^{-1/2}$, de primitive $2\sqrt{x}$ sur $]0\,;+\infty[$. Donc :
\[
B = \big[2\sqrt{x}\big]_1^4 = 2(2-1) = \boxed{2}.
\]
\item Une primitive de $\sin$ est $-\cos$. Donc :
\[
C = \big[-\cos x\big]_0^{\pi} = -\cos\pi + \cos 0 = -(-1)+1 = \boxed{2}.
\]
\item Une primitive de $\mathrm{e}^{3x}$ est $\dfrac{1}{3}\mathrm{e}^{3x}$ (car $\big(\frac{1}{3}\mathrm{e}^{3x}\big)' = \mathrm{e}^{3x}$). Donc :
\[
D = \left[\dfrac{1}{3}\mathrm{e}^{3x}\right]_0^1 = \dfrac{\mathrm{e}^3-1}{3} = \boxed{\dfrac{\mathrm{e}^3-1}{3}}.
\]
\end{enumerate}
\textbf{Point de vigilance :} pour $D$, ne pas oublier le facteur $\frac{1}{3}$ : la primitive de $\mathrm{e}^{ax}$ est $\frac{1}{a}\mathrm{e}^{ax}$, pas $\mathrm{e}^{ax}$.
\end{corrige}

\begin{corrige}[Exercice 5 — Linéarité, Chasles et valeur absolue]
\begin{enumerate}
\item Par linéarité de l'intégrale :
\[
\int_0^1 \big(3f(x)-4g(x)\big)\,\mathrm{d}x = 3\int_0^1 f(x)\,\mathrm{d}x - 4\int_0^1 g(x)\,\mathrm{d}x = 3\times 5 - 4\times(-2) = 15+8 = \boxed{23}.
\]
\item Par la relation de Chasles : $\displaystyle\int_1^5 f = \int_1^3 f + \int_3^5 f$, donc $7 = \displaystyle\int_1^3 f + 2$, d'où $\displaystyle\int_1^3 f(x)\,\mathrm{d}x = \boxed{5}$.
\item $x-1 = 0 \Leftrightarrow x=1$. Tableau de signes : $x-1\le 0$ sur $[-1\,;1]$ et $x-1\ge 0$ sur $[1\,;3]$. Par Chasles :
\begin{align*}
\int_{-1}^{3} |x-1|\,\mathrm{d}x &= \int_{-1}^{1} (1-x)\,\mathrm{d}x + \int_{1}^{3} (x-1)\,\mathrm{d}x\\
&= \left[x-\dfrac{x^2}{2}\right]_{-1}^{1} + \left[\dfrac{x^2}{2}-x\right]_{1}^{3}\\
&= \left(\dfrac{1}{2}-\left(-\dfrac{3}{2}\right)\right) + \left(\dfrac{3}{2}-\left(-\dfrac{1}{2}\right)\right) = 2+2 = \boxed{4}.
\end{align*}
\emph{Vérification géométrique :} le domaine est réunion de deux triangles rectangles de côté $2$, chacun d'aire $\frac{2\times 2}{2}=2$ ; total $4$. Cohérent.
\end{enumerate}
\end{corrige}

\begin{corrige}[Exercice 6 — Intégration par parties guidée]
\begin{enumerate}
\item $u(x)=x$ donne $u'(x)=1$ ; $v'(x)=\mathrm{e}^x$ donne on choisit $v(x)=\mathrm{e}^x$.
\item Les fonctions $u$ et $v$ sont de classe $\mathcal{C}^1$ sur $[0\,;1]$, la formule s'applique :
\[
I = \big[x\,\mathrm{e}^x\big]_0^1 - \int_0^1 \mathrm{e}^x\,\mathrm{d}x = (\mathrm{e}-0) - \big[\mathrm{e}^x\big]_0^1 = \mathrm{e}-(\mathrm{e}-1) = \boxed{1}.
\]
\item On pose $u(x)=\ln x$ (donc $u'(x)=\dfrac{1}{x}$) et $v'(x)=1$ (donc $v(x)=x$). Ces fonctions sont de classe $\mathcal{C}^1$ sur $[1\,;\mathrm{e}]$ :
\[
J = \big[x\ln x\big]_1^{\mathrm{e}} - \int_1^{\mathrm{e}} x\times\dfrac{1}{x}\,\mathrm{d}x = (\mathrm{e}\times 1 - 1\times 0) - \int_1^{\mathrm{e}} 1\,\mathrm{d}x = \mathrm{e} - (\mathrm{e}-1) = \boxed{1}.
\]
\end{enumerate}
\textbf{Points de vigilance :} (i) citer l'hypothèse « $u$ et $v$ de classe $\mathcal{C}^1$ » avant d'appliquer la formule ; (ii) pour $J$, l'astuce consiste à voir $1$ comme $v'(x)$ ; (iii) $\ln 1 = 0$ et $\ln\mathrm{e}=1$ doivent être écrits proprement.
\end{corrige}

\begin{corrige}[Exercice 7 — Changement de variable affine et forme $u'u^n$]
\begin{enumerate}
\item Posons $t=2x+1$. Alors $\mathrm{d}t = 2\,\mathrm{d}x$, soit $\mathrm{d}x = \dfrac{1}{2}\mathrm{d}t$. Bornes : $x=0 \Rightarrow t=1$ ; $x=1 \Rightarrow t=3$. Donc :
\[
\int_0^1 (2x+1)^3\,\mathrm{d}x = \dfrac{1}{2}\int_1^3 t^3\,\mathrm{d}t = \dfrac{1}{2}\left[\dfrac{t^4}{4}\right]_1^3 = \dfrac{1}{8}(81-1) = \boxed{10}.
\]
\emph{Vérification :} une primitive directe de $(2x+1)^3$ est $\frac{(2x+1)^4}{8}$ ; $\frac{3^4-1^4}{8} = \frac{80}{8} = 10$. Cohérent.
\item Posons $t=3x$. Alors $\mathrm{d}x = \dfrac{1}{3}\mathrm{d}t$. Bornes : $x=0 \Rightarrow t=0$ ; $x=\dfrac{\pi}{2} \Rightarrow t=\dfrac{3\pi}{2}$. Donc :
\[
\int_0^{\pi/2} \cos(3x)\,\mathrm{d}x = \dfrac{1}{3}\int_0^{3\pi/2} \cos t\,\mathrm{d}t = \dfrac{1}{3}\big[\sin t\big]_0^{3\pi/2} = \dfrac{1}{3}(-1-0) = \boxed{-\dfrac{1}{3}}.
\]
Le résultat négatif est normal : $\cos(3x)$ devient négatif sur $]\pi/6\,;\pi/2]$.
\item On reconnaît la forme $u'(x)\,u(x)$ avec $u(x)=\ln x$ et $u'(x)=\dfrac{1}{x}$ ($u$ est de classe $\mathcal{C}^1$ sur $[1\,;\mathrm{e}]$). Une primitive est $\dfrac{1}{2}\big(\ln x\big)^2$ :
\[
\int_1^{\mathrm{e}} \dfrac{\ln x}{x}\,\mathrm{d}x = \left[\dfrac{(\ln x)^2}{2}\right]_1^{\mathrm{e}} = \dfrac{1^2-0^2}{2} = \boxed{\dfrac{1}{2}}.
\]
\end{enumerate}
\textbf{Point de vigilance :} dans un changement de variable affine $t=ax+b$, les trois gestes obligatoires sont : exprimer $\mathrm{d}x$, transformer les bornes, ne pas oublier le facteur $\frac{1}{a}$.
\end{corrige}

\begin{corrige}[Exercice 8 — Aire et méthode des rectangles]
\begin{enumerate}
\item $f(x)=x^2+1 > 0$ sur $[0\,;2]$, donc l'aire est :
\[
\mathcal{A} = \int_0^2 (x^2+1)\,\mathrm{d}x = \left[\dfrac{x^3}{3}+x\right]_0^2 = \dfrac{8}{3}+2 = \dfrac{14}{3}\ \text{u.a.}
\]
Le repère est orthonormé d'unité \SI{1}{cm}, donc $1\ \text{u.a.} = 1\ \mathrm{cm}^2$ et $\mathcal{A} = \boxed{\dfrac{14}{3}\ \mathrm{cm}^2 \approx 4{,}67\ \mathrm{cm}^2}$.
\item Amplitude : $h = \dfrac{2-0}{4} = 0{,}5$. Subdivision : $x_0=0$, $x_1=0{,}5$, $x_2=1$, $x_3=1{,}5$, $x_4=2$.
Valeurs : $f(0)=1$ ; $f(0{,}5)=1{,}25$ ; $f(1)=2$ ; $f(1{,}5)=3{,}25$ ; $f(2)=5$.
\[
R_g = 0{,}5\,(1+1{,}25+2+3{,}25) = 0{,}5\times 7{,}5 = \boxed{3{,}75} \ ;\qquad
R_d = 0{,}5\,(1{,}25+2+3{,}25+5) = 0{,}5\times 11{,}5 = \boxed{5{,}75}.
\]
\item $f'(x)=2x \ge 0$ sur $[0\,;2]$ : $f$ est croissante. Sur chaque sous-intervalle $[x_k\,;x_{k+1}]$, on a $f(x_k)\le f(x)\le f(x_{k+1})$, donc le rectangle à gauche est sous la courbe et celui à droite la recouvre. En sommant : $R_g \le \mathcal{A} \le R_d$, soit :
\[
\boxed{3{,}75 \le \mathcal{A} \le 5{,}75}.
\]
La valeur exacte $\dfrac{14}{3}\approx 4{,}67$ appartient bien à cet encadrement. L'amplitude de l'encadrement est $R_d - R_g = h\big(f(2)-f(0)\big) = 0{,}5\times 4 = 2$ : en augmentant $n$, on réduit l'erreur.
\end{enumerate}
\end{corrige}

\begin{corrige}[Exercice 9 — Type Bac : trois techniques de calcul d'intégrales]
\emph{Barème indicatif (sur 6 points) : 1.a) 0,5 pt ; 1.b) 1 pt ; 2) 1,5 pt ; 3.a) 1 pt ; 3.b) 1 pt ; 4) 1 pt.}
\begin{enumerate}
\item \begin{enumerate}
\item $u(x)=2x+3$ donne $u'(x)=2$ ; $v'(x)=\mathrm{e}^x$ donne on choisit $v(x)=\mathrm{e}^x$.
\item $u$ et $v$ sont de classe $\mathcal{C}^1$ sur $[0\,;1]$. La formule d'intégration par parties donne :
\begin{align*}
I &= \big[(2x+3)\mathrm{e}^x\big]_0^1 - \int_0^1 2\,\mathrm{e}^x\,\mathrm{d}x = \big(5\mathrm{e} - 3\big) - 2\big[\mathrm{e}^x\big]_0^1\\
&= 5\mathrm{e}-3 - 2(\mathrm{e}-1) = 5\mathrm{e}-3-2\mathrm{e}+2 = 3\mathrm{e}-1.
\end{align*}
On a bien démontré que $I = 3\mathrm{e}-1$.
\end{enumerate}
\item Posons $u(x)=\ln x$ ; alors $u'(x)=\dfrac{1}{x}$ et l'intégrande s'écrit $u'(x)\,u(x)$, de primitive $\dfrac{1}{2}\big(u(x)\big)^2 = \dfrac{(\ln x)^2}{2}$. Comme $\ln$ est de classe $\mathcal{C}^1$ sur $[1\,;\mathrm{e}]$ :
\[
J = \left[\dfrac{(\ln x)^2}{2}\right]_1^{\mathrm{e}} = \dfrac{(\ln\mathrm{e})^2 - (\ln 1)^2}{2} = \dfrac{1-0}{2} = \dfrac{1}{2}.
\]
\item \begin{enumerate}
\item $t=2x$ donne $\mathrm{d}t = 2\,\mathrm{d}x$, soit $\mathrm{d}x = \dfrac{1}{2}\mathrm{d}t$. Bornes : $x=0 \Rightarrow t=0$ ; $x=\dfrac{\pi}{4} \Rightarrow t=\dfrac{\pi}{2}$.
\item $K = \dfrac{1}{2}\displaystyle\int_0^{\pi/2} \cos t\,\mathrm{d}t = \dfrac{1}{2}\big[\sin t\big]_0^{\pi/2} = \dfrac{1}{2}(1-0) = \boxed{\dfrac{1}{2}}$.
\end{enumerate}
\item Sur $[0\,;\pi/4]$, $2x\in[0\,;\pi/2]$ donc $\cos(2x)\ge 0$ : la courbe est au-dessus de l'axe. L'aire de la face de la tuile est donc :
\[
\mathcal{A} = \int_0^{\pi/4} \cos(2x)\,\mathrm{d}x = K = \dfrac{1}{2}\ \text{u.a.}
\]
L'unité de longueur est le décimètre, donc $1\ \text{u.a.} = 1\ \mathrm{dm}^2$ : l'aire vaut $\boxed{0{,}5\ \mathrm{dm}^2}$.
\end{enumerate}
\textbf{Points de vigilance (Bac) :} citer l'hypothèse de classe $\mathcal{C}^1$ pour l'intégration par parties ; écrire explicitement le changement de bornes ; justifier le signe de la fonction avant d'identifier intégrale et aire ; convertir correctement l'unité d'aire.
\end{corrige}

\begin{corrige}[Exercice 10 — Type Bac : fonction définie par une intégrale et méthode des rectangles]
\emph{Barème indicatif (sur 8 points) : 1) 1 pt ; 2) 1 pt ; 3.a) 1,5 pt ; 3.b) 1,5 pt ; 4.a) 1 pt ; 4.b) 2 pts.}
\begin{enumerate}
\item $f$ est continue et positive sur $[0\,;6]$. Pour $x_1 < x_2$ dans $[0\,;6]$, la relation de Chasles donne :
\[
F(x_2) - F(x_1) = \int_0^{x_2} f - \int_0^{x_1} f = \int_{x_1}^{x_2} f(t)\,\mathrm{d}t \ge 0
\]
par positivité de l'intégrale (intégrale d'une fonction positive sur $[x_1\,;x_2]$ avec $x_1<x_2$). Donc $F$ est \textbf{croissante} sur $[0\,;6]$. \emph{Conformément au programme, cette justification n'utilise ni primitive ni dérivée.}
\item $F(0) = \displaystyle\int_0^0 f(t)\,\mathrm{d}t = \boxed{0}$ (intégrale sur un intervalle réduit à un point). $F(6)$ représente le \textbf{volume total d'eau écoulé dans la conduite entre l'instant $0$ et l'instant $6$}, exprimé en $\mathrm{m}^3$.
\item \begin{enumerate}
\item $h=1$. Rectangles à gauche (hauteurs $f(0),\dots,f(5)$) :
\[
R_g = 1\times(2+4+5+4{,}5+3+2) = \boxed{20{,}5\ \mathrm{m}^3}.
\]
Rectangles à droite (hauteurs $f(1),\dots,f(6)$) :
\[
R_d = 1\times(4+5+4{,}5+3+2+1) = \boxed{19{,}5\ \mathrm{m}^3}.
\]
\item $f$ n'est pas monotone sur $[0\,;6]$ entier, donc on ne peut pas ordonner directement $R_g$, $F(6)$ et $R_d$ sur tout l'intervalle. On découpe avec Chasles :
\begin{itemize}
\item Sur $[0\,;2]$, $f$ croissante : $R_g^{[0,2]} = 2+4 = 6 \le \displaystyle\int_0^2 f \le 4+5 = 9 = R_d^{[0,2]}$.
\item Sur $[2\,;6]$, $f$ décroissante : $R_d^{[2,6]} = 4{,}5+3+2+1 = 10{,}5 \le \displaystyle\int_2^6 f \le 5+4{,}5+3+2 = 14{,}5 = R_g^{[2,6]}$.
\end{itemize}
En additionnant : $6+10{,}5 \le F(6) \le 9+14{,}5$, soit $\boxed{16{,}5 \le F(6) \le 23{,}5}$.
Une valeur approchée naturelle est la moyenne $\dfrac{R_g+R_d}{2} = \dfrac{20{,}5+19{,}5}{2} = 20$, donc le volume distribué est estimé à $\boxed{20{,}0\ \mathrm{m}^3}$ avec une incertitude de $\pm 3{,}5\ \mathrm{m}^3$ (encadrement $[16{,}5\,;\,23{,}5]$).
\end{enumerate}
\item \begin{enumerate}
\item Sur $[0\,;2]$, $f(t)=\dfrac{3}{2}t+2$ :
\[
F(2) = \int_0^2 \left(\dfrac{3}{2}t+2\right)\mathrm{d}t = \left[\dfrac{3}{4}t^2+2t\right]_0^2 = 3+4 = \boxed{7\ \mathrm{m}^3}.
\]
\emph{Cohérence :} $7\in[6\,;9]$, l'encadrement de la question 3.b).
\item Sur $[2\,;6]$, $f$ est décroissante, donc pour tout $t\in[2\,;6]$ : $f(6)\le f(t)\le f(2)$, c'est-à-dire $1\le f(t)\le 5$. L'inégalité de la moyenne (avec $b-a=4$) donne :
\[
1\times 4 \le \int_2^6 f(t)\,\mathrm{d}t \le 5\times 4 \quad\Longleftrightarrow\quad 4 \le \int_2^6 f \le 20.
\]
Comme $F(6) = F(2) + \displaystyle\int_2^6 f$ (Chasles), on obtient :
\[
\boxed{11 \le F(6) \le 27}.
\]
Cet encadrement est plus large que celui de la question 3.b) ($[16{,}5\,;23{,}5]$), car l'inégalité de la moyenne n'utilise que les valeurs extrêmes et ignore la décroissance progressive du débit ; il reste néanmoins compatible avec l'estimation $F(6)\approx 20\ \mathrm{m}^3$.
\end{enumerate}
\end{enumerate}
\textbf{Points de vigilance (Bac) :} pour la question 1, le correcteur attend explicitement Chasles + positivité, sans primitive ; pour 3.b), il faut remarquer que $f$ change de monotonie et découper l'intervalle — appliquer brutalement « $R_g \le \int \le R_d$ » sur $[0\,;6]$ est une erreur ; bien préciser les unités ($\mathrm{m}^3$) dans les réponses concrètes.
\end{corrige}

\begin{corrige}[Exercice 11 — Type Bac : aire entre deux courbes, volume de révolution et valeur moyenne]
\emph{Barème indicatif (sur 12 points) : Partie A : 4 pts ; Partie B : 4 pts ; Partie C : 4 pts.}

\textbf{Partie A : aire entre deux courbes}
\begin{enumerate}
\item $f(x)=g(x) \Leftrightarrow x^2 = 2x+3 \Leftrightarrow x^2-2x-3=0$. Discriminant : $\Delta = 4+12 = 16$, racines $x=\dfrac{2\pm 4}{2}$, soit $x=-1$ et $x=3$.
Ordonnées : $f(-1)=1$ et $f(3)=9$. Points d'intersection : $\boxed{A(-1\,;1)}$ et $\boxed{B(3\,;9)}$.
\item $g(x)-f(x) = -x^2+2x+3 = -(x-3)(x+1)$. Sur $[-1\,;3]$ : $x+1\ge 0$ et $x-3\le 0$, donc $(x-3)(x+1)\le 0$, d'où $g(x)-f(x)\ge 0$, c'est-à-dire $g(x)\ge f(x)$ : \textbf{la droite $\mathcal{D}$ est au-dessus de la parabole $\mathcal{P}$} sur $[-1\,;3]$ (avec égalité seulement aux bornes).
\item Puisque $g\ge f$ sur $[-1\,;3]$, l'aire du domaine compris entre les deux courbes est :
\begin{align*}
\mathcal{A} &= \int_{-1}^{3}\big(g(x)-f(x)\big)\,\mathrm{d}x = \int_{-1}^{3}\big(-x^2+2x+3\big)\,\mathrm{d}x = \left[-\dfrac{x^3}{3}+x^2+3x\right]_{-1}^{3}\\
&= \left(-9+9+9\right) - \left(\dfrac{1}{3}+1-3\right) = 9 - \left(-\dfrac{5}{3}\right) = \boxed{\dfrac{32}{3}\ \text{u.a.}}
\end{align*}
\item Unités graphiques : \SI{2}{cm} par axe, donc $1\ \text{u.a.} = 2\times 2 = 4\ \mathrm{cm}^2$. D'où :
\[
\mathcal{A} = \dfrac{32}{3}\times 4 = \boxed{\dfrac{128}{3}\ \mathrm{cm}^2 \approx 42{,}67\ \mathrm{cm}^2}.
\]
\end{enumerate}

\textbf{Partie B : volume de révolution}
\begin{enumerate}
\item Pour $f$ continue et positive sur $[a\,;b]$, le volume du solide engendré par la rotation de la courbe autour de $(Ox)$ est $V = \pi\displaystyle\int_a^b \big(f(x)\big)^2\,\mathrm{d}x$ (formule des disques).
\item Ici $f(x)=\sqrt{x}+1 \ge 0$ sur $[0\,;4]$ et $\big(f(x)\big)^2 = x+2\sqrt{x}+1$. Donc :
\begin{align*}
V &= \pi\int_0^4 \big(x+2\sqrt{x}+1\big)\,\mathrm{d}x = \pi\left[\dfrac{x^2}{2} + \dfrac{4}{3}x^{3/2} + x\right]_0^4\\
&= \pi\left(8 + \dfrac{4}{3}\times 8 + 4\right) = \pi\left(12+\dfrac{32}{3}\right) = \boxed{\dfrac{68\pi}{3}\ \text{u.v.}}
\end{align*}
(car $4^{3/2} = \big(\sqrt{4}\big)^3 = 8$).
\item $1\ \text{u.l.} = \SI{5}{cm}$, donc $1\ \text{u.v.} = 5^3 = 125\ \mathrm{cm}^3$ :
\[
V = \dfrac{68\pi}{3}\times 125 = \dfrac{8500\pi}{3}\ \mathrm{cm}^3 \approx \dfrac{8500\times 3{,}1416}{3} \approx 8901{,}2\ \mathrm{cm}^3.
\]
Contenance : $\dfrac{8901{,}2}{1000} \approx 8{,}901\ \mathrm{L}$, soit $\boxed{\approx 8{,}90\ \mathrm{L}}$ au centilitre près.
\end{enumerate}

\textbf{Partie C : valeur absolue et valeur moyenne}
\begin{enumerate}
\item La fonction $x\mapsto|\cos x|$ est $\pi$-périodique, donc $\displaystyle\int_0^{2\pi}|\cos x|\,\mathrm{d}x = 2\int_0^{\pi}|\cos x|\,\mathrm{d}x$. Sur $[0\,;\pi]$ : $\cos x \ge 0$ sur $[0\,;\pi/2]$ et $\cos x \le 0$ sur $[\pi/2\,;\pi]$. Par Chasles :
\[
\int_0^{\pi}|\cos x|\,\mathrm{d}x = \int_0^{\pi/2}\cos x\,\mathrm{d}x + \int_{\pi/2}^{\pi}(-\cos x)\,\mathrm{d}x = \big[\sin x\big]_0^{\pi/2} - \big[\sin x\big]_{\pi/2}^{\pi} = 1 - (0-1) = 2.
\]
Par symétrie de la courbe de $\cos$ sur $[0\,;\pi/2]$ et $[\pi/2\,;\pi]$, on a aussi $\displaystyle\int_{\pi/2}^{\pi}|\cos x|\,\mathrm{d}x = \int_0^{\pi/2}\cos x\,\mathrm{d}x$, donc $\displaystyle\int_0^{\pi}|\cos x| = 2\int_0^{\pi/2}\cos x$, et finalement :
\[
\int_0^{2\pi}|\cos x|\,\mathrm{d}x = 4\int_0^{\pi/2}\cos x\,\mathrm{d}x = 4\big[\sin x\big]_0^{\pi/2} = \boxed{4}.
\]
\item Valeur moyenne : $m = \dfrac{1}{2\pi}\displaystyle\int_0^{2\pi}|\cos x|\,\mathrm{d}x = \dfrac{4}{2\pi} = \boxed{\dfrac{2}{\pi} \approx 0{,}637}$.
\emph{Interprétation graphique :} le rectangle de base $[0\,;2\pi]$ et de hauteur $\dfrac{2}{\pi}$ a pour aire $2\pi\times\dfrac{2}{\pi} = 4$, exactement l'aire totale sous la courbe de $x\mapsto|\cos x|$ sur $[0\,;2\pi]$ : la valeur moyenne est la « hauteur équivalente » constante.
\end{enumerate}
\textbf{Points de vigilance (Bac) :} (i) pour l'aire entre deux courbes, justifier la position relative \emph{avant} d'intégrer $g-f$ ; (ii) ne pas confondre les conversions : l'aire se multiplie par le produit des unités ($4\ \mathrm{cm}^2$ ici) et le volume par le \emph{cube} de l'unité de longueur ($125\ \mathrm{cm}^3$) ; (iii) pour la valeur absolue, le découpage de Chasles aux points d'annulation est indispensable — intégrer $\cos x$ directement sur $[0\,;2\pi]$ donnerait $0$, ce qui serait faux.
\end{corrige}

\newpage

\coursec{Fiche bilan}

\coursec{Calcul intégral}

\courssub{Intégrale d'une fonction continue : définition et intuition}

\begin{savaistu}[Naissance du calcul intégral]
Le problème du calcul des aires de domaines curvilignes remonte à l'Antiquité (Archimède, méthode d'épuisement). La formalisation moderne (intégrale de Riemann) date du XIX\textsuperscript{e} siècle. Au Cameroun, on retrouve cette idée dans le calcul de surfaces agricoles (parcelles irrégulières), de volumes de retenues d'eau (barrages de Lagdo, Memve'ele) ou de débits de flux (trafic MTN/Orange, débit CAMTEL).
\end{savaistu}

\begin{experience}[Approche par les rectangles : de l'approximation à la définition]
Soit $f$ une fonction continue et \textbf{positive} sur $[a;b]$ ($a<b$). On veut mesurer l'aire $\mathcal{A}$ du domaine $\mathcal{D}$ limité par $\mathcal{C}_f$, l'axe $(Ox)$ et les droites $x=a$, $x=b$.
\begin{enumerate}
  \item On partage $[a;b]$ en $n$ sous-intervalles de même pas $h_n=\frac{b-a}{n}$ : $x_k=a+kh_n$.
  \item Sur chaque $[x_{k-1};x_k]$, on construit deux rectangles :
  \begin{itemize}
    \item Rectangle \textbf{inscrit} (somme inférieure) : hauteur $m_k = \min_{[x_{k-1},x_k]} f$ (existe car $f$ continue).
    \item Rectangle \textbf{circonscrit} (somme supérieure) : hauteur $M_k = \max_{[x_{k-1},x_k]} f$.
  \end{itemize}
  \item Aires totales : $s_n = h_n\sum_{k=1}^n m_k$ (somme inférieure), $S_n = h_n\sum_{k=1}^n M_k$ (somme supérieure).
  \item On a toujours $s_n \le \mathcal{A} \le S_n$. Quand $n\to +\infty$, $h_n\to 0$ et $S_n-s_n \to 0$ (continuité uniforme de $f$ sur le segment). Les deux suites convergent vers une même limite.
\end{enumerate}
Cette limite commune \textbf{est} l'intégrale de $f$ sur $[a;b]$.
\end{experience}

\begin{definition}[Intégrale d'une fonction continue]
Soit $f$ une fonction continue sur un intervalle $[a;b]$ ($a,b\in\mathbb{R}$, $a\le b$).
\begin{itemize}
  \item Si $a<b$ et $f\ge 0$ sur $[a;b]$, $\displaystyle\int_a^b f(x)\,\mathrm{d}x$ est l'\cle{aire} (en unités d'aire) du domaine borné par $\mathcal{C}_f$, l'axe des abscisses et les droites $x=a$, $x=b$.
  \item Si $f$ est de signe quelconque, on pose $f^+=\max(f,0)$ et $f^-=\max(-f,0)$ (parties positive et négative). Alors :
  \[
  \int_a^b f(x)\,\mathrm{d}x = \int_a^b f^+(x)\,\mathrm{d}x - \int_a^b f^-(x)\,\mathrm{d}x.
  \]
  C'est l'\cle{aire algébrique} (aire au-dessus de l'axe $-$ aire en-dessous).
  \item Par convention : $\displaystyle\int_a^a f(x)\,\mathrm{d}x = 0$ et $\displaystyle\int_b^a f(x)\,\mathrm{d}x = -\int_a^b f(x)\,\mathrm{d}x$ pour $a<b$.
\end{itemize}
La notation $\mathrm{d}x$ rappelle la variable d'intégration ; $x$ est une \cle{variable muette} (on peut la remplacer par $t$, $u$, etc.).
\end{definition}

\begin{propriete}[Existence de l'intégrale]
Toute fonction continue sur un segment $[a;b]$ est \cle{intégrable} sur $[a;b]$ (l'intégrale existe et est un nombre réel fini).
\end{propriete}

\begin{exemplebox}[Aires élémentaires et intégrales]
\begin{enumerate}
  \item $f(x)=c$ (constante) sur $[a;b]$ : $\mathcal{A}=c(b-a)$ $\Rightarrow$ $\int_a^b c\,\mathrm{d}x = c(b-a)$.
  \item $f(x)=x$ sur $[0;1]$ : triangle rectangle d'aire $1/2$ $\Rightarrow$ $\int_0^1 x\,\mathrm{d}x = \frac{1}{2}$.
  \item $f(x)=\sqrt{1-x^2}$ sur $[-1;1]$ : demi-disque unité $\Rightarrow$ $\int_{-1}^1 \sqrt{1-x^2}\,\mathrm{d}x = \frac{\pi}{2}$.
\end{enumerate}
\end{exemplebox}

\courssub{Propriétés algébriques fondamentales}

\begin{propriete}[Linéarité de l'intégrale]
Pour $f,g$ continues sur $[a;b]$ et $\alpha,\beta\in\mathbb{R}$ :
\[
\int_a^b \big(\alpha f(x)+\beta g(x)\big)\,\mathrm{d}x = \alpha\int_a^b f(x)\,\mathrm{d}x + \beta\int_a^b g(x)\,\mathrm{d}x.
\]
\emph{Preuve (idée) :} Vraie pour les sommes de Riemann (linéarité de la somme), passe à la limite.
\end{propriete}

\begin{propriete}[Relation de Chasles]
Pour $f$ continue sur un intervalle $I$ et $a,b,c\in I$ (dans n'importe quel ordre) :
\[
\int_a^b f(x)\,\mathrm{d}x = \int_a^c f(x)\,\mathrm{d}x + \int_c^b f(x)\,\mathrm{d}x.
\]
\emph{Interprétation géométrique :} L'aire totale est la somme des aires partielles (valable pour les aires algébriques grâce aux conventions sur le sens des bornes).
\end{propriete}

\begin{propriete}[Positivité, comparaison, inégalité de la moyenne, valeur absolue]
Soient $f,g$ continues sur $[a;b]$ avec $a<b$.
\begin{enumerate}
  \item \textbf{Positivité :} Si $f(x)\ge 0$ sur $[a;b]$, alors $\displaystyle\int_a^b f(x)\,\mathrm{d}x \ge 0$.
  \item \textbf{Comparaison :} Si $f(x)\le g(x)$ sur $[a;b]$, alors $\displaystyle\int_a^b f(x)\,\mathrm{d}x \le \int_a^b g(x)\,\mathrm{d}x$.
  \item \textbf{Inégalité de la moyenne (encadrement) :} Si $m\le f(x)\le M$ sur $[a;b]$, alors
  \[
  m(b-a) \le \int_a^b f(x)\,\mathrm{d}x \le M(b-a).
  \]
  \item \textbf{Valeur absolue :} $\displaystyle\left|\int_a^b f(x)\,\mathrm{d}x\right| \le \int_a^b |f(x)|\,\mathrm{d}x$.
\end{enumerate}
\emph{Preuves :} Découlent de la positivité appliquée à $g-f$, $f-m$, $M-f$, et $|f|\pm f$.
\end{propriete}

\begin{exemplebox}[Encadrement d'une intégrale sans la calculer]
Encadrer $I=\int_0^1 \frac{\mathrm{d}x}{1+x^2}$.
Sur $[0;1]$, $1\le 1+x^2 \le 2 \Rightarrow \frac{1}{2} \le \frac{1}{1+x^2} \le 1$.
Par inégalité de la moyenne : $\frac{1}{2}(1-0) \le I \le 1(1-0) \Rightarrow \boxed{\frac{1}{2} \le I \le 1}$.
(On sait que $I=\frac{\pi}{4}\approx 0,785$, l'encadrement est correct).
\end{exemplebox}

\courssub{Fonction définie par une intégrale et valeur moyenne}

\begin{definition}[Fonction intégrale]
Soit $f$ continue sur un intervalle $I$ et $a\in I$ fixé. La fonction $F$ définie sur $I$ par
\[
F(x) = \int_a^x f(t)\,\mathrm{d}t
\]
est appelée \cle{fonction intégrale} de $f$ d'origine $a$.
\end{definition}

\begin{propriete}[Sens de variation de la fonction intégrale]
Soit $F(x)=\int_a^x f(t)\,\mathrm{d}t$ avec $f$ continue sur $I$.
\begin{itemize}
  \item Si $f(x)\ge 0$ sur $I$, alors $F$ est \textbf{croissante} sur $I$.
  \item Si $f(x)\le 0$ sur $I$, alors $F$ est \textbf{décroissante} sur $I$.
  \item Plus généralement, pour $x_1<x_2$ dans $I$ : $F(x_2)-F(x_1) = \int_{x_1}^{x_2} f(t)\,\mathrm{d}t$ a le signe de $f$ sur $[x_1;x_2]$.
\end{itemize}
\emph{Justification (sans primitive) :} Pour $x_2>x_1$, $F(x_2)-F(x_1)=\int_{x_1}^{x_2} f$. Si $f\ge 0$, l'intégrale est $\ge 0$ (positivité), donc $F(x_2)\ge F(x_1)$.
\end{propriete}

\begin{exemplebox}[Variation de $F(x)=\int_0^x (t^2-1)\,\mathrm{d}t$]
$f(t)=t^2-1$. $f(t)\le 0$ sur $[-1;1]$, $f(t)\ge 0$ sur $]-\infty;-1]$ et $[1;+\infty[$.
Donc $F$ décroît sur $[-1;1]$ et croît sur $]-\infty;-1]$ et $[1;+\infty[$.
$F$ a un minimum en $x=1$ (puisque $F$ passe de décroissante à croissante).
\end{exemplebox}

\begin{definition}[Valeur moyenne d'une fonction]
La \cle{valeur moyenne} de $f$ (continue) sur $[a;b]$ ($a<b$) est le nombre réel :
\[
\mu = \frac{1}{b-a}\int_a^b f(x)\,\mathrm{d}x.
\]
\emph{Interprétation :} C'est la hauteur du rectangle de base $[a;b]$ ayant la même aire (algébrique) que le domaine sous $\mathcal{C}_f$.
\end{definition}

\begin{propriete}[Théorème de la valeur moyenne intégrale -- Culture Bac]
Si $f$ est continue sur $[a;b]$, il existe au moins un $c\in[a;b]$ tel que $f(c)=\mu = \frac{1}{b-a}\int_a^b f$.
\emph{Preuve :} $m\le f\le M \Rightarrow m\le \mu \le M$ (inégalité de la moyenne). Par le théorème des valeurs intermédiaires (TVI), $\exists c, f(c)=\mu$.
\end{propriete}

\courssub{Techniques de calcul d'intégrales}

\begin{methode}[Calcul par primitive directe (Théorème fondamental de l'analyse)]
\begin{enumerate}
  \item Trouver une \cle{primitive} $F$ de $f$ sur $[a;b]$ (i.e. $F'=f$ sur $[a;b]$).
  \item Appliquer la formule : $\displaystyle\int_a^b f(x)\,\mathrm{d}x = F(b)-F(a) = \big[F(x)\big]_a^b$.
  \item \textbf{Vérifier systématiquement} en dérivant $F$ qu'on retrouve bien $f$.
\end{enumerate}
\end{methode}

\begin{aretenir}[Tableau des primitives usuelles (à connaître par cœur)]
\begin{center}
\begin{tabular}{|c|c|c|}
\hline
\rowcolor{popblueL} \textbf{Fonction $f(x)$} & \textbf{Primitive $F(x)$ sur $I$} & \textbf{Intervalle $I$} \\
\hline
$x^\alpha$ ($\alpha\in\mathbb{Q}$, $\alpha\ne -1$) & $\dfrac{x^{\alpha+1}}{\alpha+1}$ & $]0;+\infty[$ (ou $\mathbb{R}$ si $\alpha\in\mathbb{N}$) \\
\hline
$\dfrac{1}{x}$ & $\ln|x|$ & $]-\infty;0[$ ou $]0;+\infty[$ \\
\hline
$\mathrm{e}^x$ & $\mathrm{e}^x$ & $\mathbb{R}$ \\
\hline
$\cos x$ & $\sin x$ & $\mathbb{R}$ \\
\hline
$\sin x$ & $-\cos x$ & $\mathbb{R}$ \\
\hline
$\dfrac{1}{\cos^2 x}$ & $\tan x$ & $]-\frac{\pi}{2}+k\pi;\frac{\pi}{2}+k\pi[$ \\
\hline
$\dfrac{1}{1+x^2}$ & $\arctan x$ & $\mathbb{R}$ \\
\hline
$\dfrac{1}{\sqrt{1-x^2}}$ & $\arcsin x$ & $]-1;1[$ \\
\hline
\end{tabular}
\end{center}
\emph{Rappel :} Toute primitive diffère d'une constante. Pour l'intégrale définie, la constante s'annule.
\end{aretenir}

\begin{exemplebox}[Calculs directs]
\begin{enumerate}
  \item $\int_0^1 (3x^2+2x+1)\,\mathrm{d}x = \big[x^3+x^2+x\big]_0^1 = 3$.
  \item $\int_1^{\mathrm{e}} \frac{\mathrm{d}x}{x} = \big[\ln x\big]_1^{\mathrm{e}} = 1$.
  \item $\int_0^{\pi/4} \frac{\mathrm{d}x}{\cos^2 x} = \big[\tan x\big]_0^{\pi/4} = 1$.
  \item $\int_0^1 \frac{\mathrm{d}x}{1+x^2} = \big[\arctan x\big]_0^1 = \frac{\pi}{4}$.
\end{enumerate}
\end{exemplebox}

\begin{experience}[Démonstration de la formule d'intégration par parties]
Soient $u,v$ de classe $\mathcal{C}^1$ sur $[a;b]$. On dérive le produit : $(uv)' = u'v + uv'$.
On intègre entre $a$ et $b$ : $\int_a^b (uv)' = \int_a^b u'v + \int_a^b uv'$.
Or $\int_a^b (uv)' = \big[uv\big]_a^b$ (primitive directe).
D'où $\big[uv\big]_a^b = \int_a^b u'v + \int_a^b uv'$, soit la formule.
\end{experience}

\begin{propriete}[Intégration par parties]
Pour $u,v$ de classe $\mathcal{C}^1$ sur $[a;b]$ :
\[
\int_a^b u(x)v'(x)\,\mathrm{d}x = \big[u(x)v(x)\big]_a^b - \int_a^b u'(x)v(x)\,\mathrm{d}x.
\]
\end{propriete}

\begin{methode}[Étapes de l'intégration par parties]
\begin{enumerate}
  \item Identifier le produit $u \cdot v'$ dans l'intégrande (souvent indiqué par l'énoncé).
  \item Choisir $u$ (à dériver) et $v'$ (à intégrer). \textbf{Règle pratique :} $u$ = polynôme, $\ln$, $\arctan$ (se simplifient en dérivant) ; $v'$ = $\mathrm{e}^x$, $\cos$, $\sin$, $1$ (s'intègrent facilement).
  \item Calculer $u'$ et $v$ (une primitive de $v'$, sans constante).
  \item Appliquer la formule : $\int u v' = [uv] - \int u' v$.
  \item Calculer la nouvelle intégrale (souvent plus simple) et conclure.
\end{enumerate}
\end{methode}

\begin{exemplebox}[Intégration par parties : exemples classiques]
\begin{enumerate}
  \item $I=\int_0^1 x\mathrm{e}^x\,\mathrm{d}x$. $u=x\Rightarrow u'=1$ ; $v'=\mathrm{e}^x\Rightarrow v=\mathrm{e}^x$.
  $I = \big[x\mathrm{e}^x\big]_0^1 - \int_0^1 \mathrm{e}^x\,\mathrm{d}x = \mathrm{e} - \big[\mathrm{e}^x\big]_0^1 = \mathrm{e} - (\mathrm{e}-1) = 1$.
  \item $J=\int_1^{\mathrm{e}} \ln x\,\mathrm{d}x$. $u=\ln x\Rightarrow u'=1/x$ ; $v'=1\Rightarrow v=x$.
  $J = \big[x\ln x\big]_1^{\mathrm{e}} - \int_1^{\mathrm{e}} x\cdot\frac{1}{x}\,\mathrm{d}x = \mathrm{e} - \big[x\big]_1^{\mathrm{e}} = \mathrm{e} - (\mathrm{e}-1) = 1$.
  \item $K=\int_0^{\pi/2} x\cos x\,\mathrm{d}x$. $u=x, v'=\cos x \Rightarrow u'=1, v=\sin x$.
  $K = \big[x\sin x\big]_0^{\pi/2} - \int_0^{\pi/2} \sin x\,\mathrm{d}x = \frac{\pi}{2} - \big[-\cos x\big]_0^{\pi/2} = \frac{\pi}{2} - 1$.
\end{enumerate}
\end{exemplebox}

\begin{attention}[Pièges fréquents : Intégration par parties]
\begin{itemize}
  \item Oublier d'évaluer le terme $[uv]_a^b$ aux bornes (ne pas laisser $x$ dedans).
  \item Choisir $u$ et $v'$ à l'envers (intégral qui se complique au lieu de se simplifier).
  \item Oublier le signe $-$ devant la nouvelle intégrale.
  \item Ne pas vérifier que $u,v$ sont $\mathcal{C}^1$ sur l'intervalle (souvent implicite au Bac).
\end{itemize}
\end{attention}

\begin{experience}[Démonstration du changement de variable affine]
Soit $I=\int_a^b f(\alpha x+\beta)\,\mathrm{d}x$ avec $\alpha\ne 0$. Posons $u=\alpha x+\beta$. Alors $\mathrm{d}u = \alpha\,\mathrm{d}x \iff \mathrm{d}x = \frac{\mathrm{d}u}{\alpha}$.
Bornes : $x=a \Rightarrow u=\alpha a+\beta$ ; $x=b \Rightarrow u=\alpha b+\beta$.
Si $\alpha>0$, l'ordre des bornes est conservé : $I = \int_{\alpha a+\beta}^{\alpha b+\beta} f(u)\frac{\mathrm{d}u}{\alpha} = \frac{1}{\alpha}\int_{\alpha a+\beta}^{\alpha b+\beta} f(u)\,\mathrm{d}u$.
Si $\alpha<0$, les bornes s'inversent, le signe $-$ du $\mathrm{d}x$ compense l'inversion : la formule $\frac{1}{\alpha}[F(\alpha x+\beta)]_a^b$ reste valide avec $F'=f$.
\end{experience}

\begin{propriete}[Changement de variable affine]
Pour $f$ continue et $\alpha\ne 0$, $\beta\in\mathbb{R}$ :
\[
\int_a^b f(\alpha x+\beta)\,\mathrm{d}x = \frac{1}{\alpha}\Big[F(\alpha x+\beta)\Big]_a^b = \frac{1}{\alpha}\int_{\alpha a+\beta}^{\alpha b+\beta} f(u)\,\mathrm{d}u
\]
où $F$ est une primitive de $f$ (ou directement sur l'intégrale avec changement de bornes).
\end{propriete}

\begin{methode}[Changement de variable affine : mode opératoire]
\begin{enumerate}
  \item Repérer la forme $f(ax+b)$.
  \item Poser $u = ax+b$, exprimer $\mathrm{d}x = \frac{\mathrm{d}u}{a}$.
  \item Modifier les bornes : $x=a \to u=aa+b$, $x=b \to u=ab+b$.
  \item Intégrer en $u$ (primitive de $f$) et appliquer les nouvelles bornes, \textbf{ou} revenir en $x$ avec le facteur $\frac{1}{a}$.
\end{enumerate}
\end{methode}

\begin{exemplebox}[Changement de variable affine]
\begin{enumerate}
  \item $\int_0^1 (2x+1)^3\,\mathrm{d}x$. $u=2x+1$, $\mathrm{d}x=\mathrm{d}u/2$. Bornes : $0\to 1$, $1\to 3$.
  $I = \int_1^3 u^3 \frac{\mathrm{d}u}{2} = \frac{1}{2}\big[\frac{u^4}{4}\big]_1^3 = \frac{1}{8}(81-1)=10$.
  \item $\int_0^{\pi/2} \cos(2x+\pi/3)\,\mathrm{d}x$. $u=2x+\pi/3$, $\mathrm{d}x=\mathrm{d}u/2$. Bornes : $0\to \pi/3$, $\pi/2\to 4\pi/3$.
  $I = \frac{1}{2}\int_{\pi/3}^{4\pi/3} \cos u\,\mathrm{d}u = \frac{1}{2}\big[\sin u\big]_{\pi/3}^{4\pi/3} = \frac{1}{2}(-\frac{\sqrt{3}}{2} - \frac{\sqrt{3}}{2}) = -\frac{\sqrt{3}}{2}$.
\end{enumerate}
\end{exemplebox}

\begin{attention}[Pièges : Changement affine]
\begin{itemize}
  \item Oublier le facteur $\frac{1}{\alpha}$ (erreur n°1 au Bac).
  \item Ne pas changer les bornes (si on travaille directement en $u$).
  \item Confondre $\alpha$ (coefficient de $x$) et $\beta$ (constante).
\end{itemize}
\end{attention}

\courssub{Applications géométriques : aires et volumes de révolution}

\begin{methode}[Aire entre une courbe et l'axe des abscisses]
Pour $f$ continue sur $[a;b]$ ($a<b$) :
\begin{enumerate}
  \item Étudier le signe de $f$ sur $[a;b]$ (tableau de signes).
  \item Si $f$ ne change pas de signe : $\mathcal{A} = \left|\int_a^b f(x)\,\mathrm{d}x\right|$.
  \item Si $f$ change de signe en $c_1,\dots,c_k$ : découper $[a;b]$ avec Chasles.
  $\mathcal{A} = \sum \left|\int_{\text{sous-intervalle}} f(x)\,\mathrm{d}x\right|$.
\end{enumerate}
\end{methode}

\begin{exemplebox}[Aire avec changement de signe]
Aire entre $\mathcal{C}_f : y=\sin x$ et $(Ox)$ sur $[0;2\pi]$.
$\sin x \ge 0$ sur $[0;\pi]$, $\sin x \le 0$ sur $[\pi;2\pi]$.
$\mathcal{A} = \int_0^\pi \sin x\,\mathrm{d}x - \int_\pi^{2\pi} \sin x\,\mathrm{d}x = \big[-\cos x\big]_0^\pi - \big[-\cos x\big]_\pi^{2\pi} = 2 - (-2) = 4$.
\end{exemplebox}

\begin{methode}[Aire entre deux courbes]
Soient $f,g$ continues sur $[a;b]$ avec $f(x)\le g(x)$ (courbe $\mathcal{C}_g$ \textbf{au-dessus} de $\mathcal{C}_f$).
L'aire du domaine borné par $\mathcal{C}_f$, $\mathcal{C}_g$, $x=a$, $x=b$ est :
\[
\mathcal{A} = \int_a^b \big(g(x)-f(x)\big)\,\mathrm{d}x.
\]
Si les courbes se coupent, déterminer les intersections (bornes) et le sens ($g-f$ ou $f-g$) sur chaque sous-intervalle.
\end{methode}

\begin{exemplebox}[Aire entre deux courbes]
Aire entre $y=x^2$ et $y=x$ (intersections $x=0,1$). Sur $[0;1]$, $x\ge x^2$.
$\mathcal{A} = \int_0^1 (x-x^2)\,\mathrm{d}x = \big[\frac{x^2}{2}-\frac{x^3}{3}\big]_0^1 = \frac{1}{2}-\frac{1}{3}=\frac{1}{6}$.
\end{exemplebox}

\begin{propriete}[Volume de révolution autour de $(Ox)$]
Soit $f$ continue et \textbf{positive} sur $[a;b]$ ($a<b$). Le volume du solide engendré par la rotation du domaine sous $\mathcal{C}_f$ autour de l'axe des abscisses est :
\[
V = \pi\int_a^b \big[f(x)\big]^2\,\mathrm{d}x.
\]
\emph{Preuve (méthode des disques) :} Tranche d'épaisseur $\mathrm{d}x$ $\approx$ disque de rayon $f(x)$, volume $\pi f(x)^2\,\mathrm{d}x$. Somme = intégrale.
\end{propriete}

\begin{exemplebox}[Volumes de révolution]
\begin{enumerate}
  \item \textbf{Cône} : $f(x)=\frac{R}{h}x$ sur $[0;h]$. $V=\pi\int_0^h \frac{R^2}{h^2}x^2\,\mathrm{d}x = \pi\frac{R^2}{h^2}\frac{h^3}{3} = \frac{1}{3}\pi R^2 h$.
  \item \textbf{Sphère} (demi-cercle $y=\sqrt{R^2-x^2}$ sur $[-R;R]$) : $V=\pi\int_{-R}^R (R^2-x^2)\,\mathrm{d}x = \pi\big[R^2x-\frac{x^3}{3}\big]_{-R}^R = \frac{4}{3}\pi R^3$.
  \item \textbf{Tore} (cercle $(x-R)^2+y^2=r^2$, $R>r$, rotation autour de $Oy$) : utilisation de la méthode des anneaux (cylindres) ou théorème de Guldin ($V=2\pi^2 R r^2$).
\end{enumerate}
\end{exemplebox}

\begin{attention}[Pièges : Volumes]
\begin{itemize}
  \item Oublier le $\pi$ ou le carré sur $f(x)$.
  \item Appliquer la formule $V=\pi\int f^2$ si $f$ n'est pas positive (prend la valeur absolue ou découpe).
  \item Confondre rotation autour de $(Ox)$ (disques $\perp Ox$) et $(Oy)$ (anneaux $\perp Oy$, formule $V=2\pi\int x f(x)\,\mathrm{d}x$ ou changement de variable).
\end{itemize}
\end{attention}

\courssub{Approximation d'une intégrale : méthode des rectangles}

\begin{definition}[Subdivision et sommes de Riemann]
Une \cle{subdivision} de $[a;b]$ ($a<b$) est une suite $a=x_0<x_1<\dots<x_n=b$. Le \cle{pas} est $h=\max(x_k-x_{k-1})$. Pour une subdivision \textbf{régulière} : $x_k=a+kh$, $h=\frac{b-a}{n}$.
Les \cle{sommes de Riemann} (rectangles) sont :
\begin{itemize}
  \item Somme inférieure (rectangles inscrits) : $s_n = \sum_{k=1}^n f(m_k)(x_k-x_{k-1})$ avec $m_k\in[x_{k-1},x_k]$ tel que $f(m_k)=\min f$.
  \item Somme supérieure (rectangles circonscrits) : $S_n = \sum_{k=1}^n f(M_k)(x_k-x_{k-1})$ avec $M_k$ tel que $f(M_k)=\max f$.
\end{itemize}
Pour $f$ continue, $\lim_{h\to 0} s_n = \lim_{h\to 0} S_n = \int_a^b f$.
\end{definition}

\begin{propriete}[Encadrement par rectangles pour $f$ monotone]
Si $f$ est continue et \textbf{monotone} sur $[a;b]$, avec une subdivision régulière de pas $h=\frac{b-a}{n}$ :
\begin{itemize}
  \item $f$ \textbf{croissante} : $\displaystyle h\sum_{k=0}^{n-1} f(x_k) \le \int_a^b f \le h\sum_{k=1}^{n} f(x_k)$ (somme gauche $\le$ intégrale $\le$ somme droite).
  \item $f$ \textbf{décroissante} : $\displaystyle h\sum_{k=1}^{n} f(x_k) \le \int_a^b f \le h\sum_{k=0}^{n-1} f(x_k)$ (somme droite $\le$ intégrale $\le$ somme gauche).
\end{itemize}
L'erreur commise est majorée par $\frac{(b-a)}{n}|f(b)-f(a)|$.
\end{propriete}

\begin{methode}[Encadrer une intégrale par la méthode des rectangles]
\begin{enumerate}
  \item Vérifier que $f$ est continue et monotone sur $[a;b]$. Déterminer son sens de variation.
  \item Choisir $n$ (donné ou à déterminer pour une précision $\varepsilon$).
  \item Calculer le pas $h=\frac{b-a}{n}$ et les abscisses $x_k=a+kh$.
  \item Selon la monotonie, écrire l'encadrement avec les sommes gauche/droite.
  \item Calculer numériquement les bornes (souvent somme de suites géométriques/arithmétiques ou calculatrice).
\end{enumerate}
\end{methode}

\begin{exemplebox}[Encadrement de $\ln 2 = \int_1^2 \frac{\mathrm{d}x}{x}$]
$f(x)=1/x$ est continue et \textbf{décroissante} sur $[1;2]$.
Pour $n=4$, $h=0,25$. $x_0=1, x_1=1,25, x_2=1,5, x_3=1,75, x_4=2$.
Somme droite (inférieure) : $S_d = 0,25(1/1,25 + 1/1,5 + 1/1,75 + 1/2) \approx 0,25(0,8+0,667+0,571+0,5) \approx 0,634$.
Somme gauche (supérieure) : $S_g = 0,25(1/1 + 1/1,25 + 1/1,5 + 1/1,75) \approx 0,25(1+0,8+0,667+0,571) \approx 0,759$.
Encadrement : $\boxed{0,634 \le \ln 2 \le 0,759}$ (valeur exacte $\approx 0,693$).
\end{exemplebox}

\begin{savaistu}[Au-delà des rectangles]
La méthode des rectangles converge lentement (erreur $\sim 1/n$). La \cle{méthode des trapèzes} (moyenne gauche/droite) converge en $1/n^2$. La \cle{méthode de Simpson} (paraboles) converge en $1/n^4$. Ces méthodes sont utilisées en informatique (simulation numérique, CAO, traitement de signal MTN/Orange) et en ingénierie (calcul de structures, hydraulique).
\end{savaistu}

\courssub{Synthèse visuelle et récapitulatif}

\begin{popfigure}
\begin{tikzpicture}[mindmap, grow cyclic, every node/.style={concept, font=\small, execute at begin node=\hskip0pt},
  root concept/.append style={concept color=popblue, font=\bfseries, minimum size=3.5cm},
  level 1/.append style={level distance=5cm, sibling angle=60, font=\scriptsize},
  level 2/.append style={level distance=3cm, sibling angle=45, font=\tiny}]
\node[root concept]{Calcul intégral}
  child[concept color=poporange]{ node {1. Définition}
    child { node {Aire algébrique} }
    child { node {$f$ continue $\Rightarrow$ intégrable} }
    child { node {Notation $\int_a^b f$} }
  }
  child[concept color=popgreen]{ node {2. Propriétés}
    child { node {Linéarité} }
    child { node {Chasles} }
    child { node {Positivité / Comparaison} }
    child { node {Inégalité moyenne} }
    child { node {$|\int f|\le \int|f|$} }
  }
  child[concept color=poppurple]{ node {3. Fonction $F(x)=\int_a^x f$}
    child { node {Variation = signe de $f$} }
    child { node {Valeur moyenne $\mu$} }
    child { node {TVM intégrale : $\exists c, f(c)=\mu$} }
  }
  child[concept color=poppink]{ node {4. Techniques}
    child { node {Primitive directe} }
    child { node {Int. par parties} }
    child { node {Chgt variable affine} }
  }
  child[concept color=popgold]{ node {5. Géométrie}
    child { node {Aire courbe/axe} }
    child { node {Aire 2 courbes} }
    child { node {Volume révolution $V=\pi\int f^2$} }
  }
  child[concept color=popmuted]{ node {6. Approximation}
    child { node {Rectangles (gauche/droite)} }
    child { node {Encadrement monotone} }
    child { node {Erreur $\sim 1/n$} }
  };
\end{tikzpicture}
\legende{Carte mentale structurée de la leçon M08 : les 6 piliers du programme officiel.}
\end{popfigure}

\begin{propriete}[Tableau récapitulatif des formules du Bac]
\begin{center}
\begin{tabularx}{\linewidth}{|l|X|}
\hline
\rowcolor{popblueL} \textbf{Propriété / Formule} & \textbf{Énoncé / Expression} \\
\hline
Linéarité & $\int_a^b(\alpha f+\beta g)=\alpha\int_a^b f+\beta\int_a^b g$ \\
\hline
Chasles & $\int_a^b f=\int_a^c f+\int_c^b f$ ($\forall a,b,c\in I$) \\
\hline
Positivité & $f\ge 0 \Rightarrow \int_a^b f\ge 0$ ($a<b$) \\
\hline
Comparaison & $f\le g \Rightarrow \int_a^b f\le\int_a^b g$ \\
\hline
Inégalité moyenne & $m\le f\le M \Rightarrow m(b-a)\le\int_a^b f\le M(b-a)$ \\
\hline
Valeur absolue & $\left|\int_a^b f\right|\le \int_a^b |f|$ \\
\hline
Valeur moyenne & $\mu=\frac{1}{b-a}\int_a^b f$ \\
\hline
Int. par parties & $\int_a^b u v' = [uv]_a^b - \int_a^b u' v$ \\
\hline
Chgt variable affine & $\int_a^b f(\alpha x+\beta)\mathrm{d}x = \frac{1}{\alpha}\int_{\alpha a+\beta}^{\alpha b+\beta} f(u)\mathrm{d}u$ \\
\hline
Aire (courbe/axe) & $\mathcal{A}=\int_a^b |f(x)|\mathrm{d}x$ (découper si signe change) \\
\hline
Aire (2 courbes) & $\mathcal{A}=\int_a^b |g(x)-f(x)|\mathrm{d}x$ (haut $-$ bas) \\
\hline
Volume (Ox) & $V=\pi\int_a^b [f(x)]^2\mathrm{d}x$ ($f\ge 0$) \\
\hline
Rectangles (croiss.) & $h\sum_{0}^{n-1}f(x_k) \le \int_a^b f \le h\sum_{1}^{n}f(x_k)$ \\
\hline
\end{tabularx}
\end{center}
\end{propriete}

\begin{methode}[Check-list de rédaction pour le Bac]
\begin{enumerate}
  \item \textbf{Définition :} Préciser « $f$ continue sur $[a;b]$ », donner le sens géométrique (aire algébrique).
  \item \textbf{Propriétés :} Citer le nom (Linéarité, Chasles, Positivité) avant d'appliquer.
  \item \textbf{Primitive :} Écrire « Une primitive de $f$ sur $I$ est $F(x)=...$ car $F'(x)=...=f(x)$ ». Puis $\int_a^b f = [F]_a^b = F(b)-F(a)$.
  \item \textbf{Parties :} Écrire explicitement : « Posons $u=...$, $v'=...$ donc $u'=...$, $v=...$ ». Puis formule.
  \item \textbf{Changement affine :} Écrire « Posons $u=\alpha x+\beta$, $\mathrm{d}x=\mathrm{d}u/\alpha$ ». Changer les bornes OU garder $x$ avec facteur $1/\alpha$.
  \item \textbf{Aire :} \textbf{Obligatoire} : tableau de signes de $f$ (ou $g-f$). Découper avec Chasles. Intégrer la valeur absolue.
  \item \textbf{Volume :} Vérifier $f\ge 0$. Écrire $V=\pi\int_a^b f(x)^2\mathrm{d}x$. Calculer.
  \item \textbf{Rectangles :} Dire « $f$ continue et monotone sur $[a;b]$ ». Donner $h$, $x_k$. Écrire l'encadrement selon la variation. Conclure.
\end{enumerate}
\end{methode}

\begin{aretenir}[Checklist finale avant le Baccalauréat]
\begin{itemize}
  \item[$\square$] Je définis l'intégrale comme une aire (positive) ou aire algébrique (signe quelconque) pour $f$ continue.
  \item[$\square$] J'applique sans hésiter : Linéarité, Chasles, Positivité, Comparaison, Inégalité de la moyenne, $|\int f|\le\int|f|$.
  \item[$\square$] Je connais par cœur le tableau des primitives (puissances, $1/x$, $\mathrm{e}^x$, $\cos$, $\sin$, $1/\cos^2$, $1/(1+x^2)$, $1/\sqrt{1-x^2}$).
  \item[$\square$] Je vérifie \textbf{systématiquement} ma primitive en la dérivant.
  \item[$\square$] Je maîtrise l'intégration par parties : choix $u/v'$, calcul $u'/v$, formule, évaluation bornes.
  \item[$\square$] Je gère le changement de variable affine : facteur $1/\alpha$, changement de bornes (ou retour en $x$).
  \item[$\square$] J'étudie le sens de variation de $x\mapsto\int_a^x f(t)\mathrm{d}t$ via le signe de $f$.
  \item[$\square$] Je calcule une valeur moyenne $\mu$ et j'applique l'inégalité de la moyenne pour encadrer.
  \item[$\square$] Pour une aire : j'étudie le signe, je découpe (Chasles), j'intègre $|f|$ ou $|g-f|$.
  \item[$\square$] Pour un volume : je vérifie $f\ge 0$, j'applique $V=\pi\int f^2$, je ne oublie pas $\pi$ et le carré.
  \item[$\square$] Pour les rectangles : je vérifie la monotonie, je donne $h$, $x_k$, j'écris l'encadrement correct (gauche/droite selon croissante/décroissante).
\end{itemize}
\end{aretenir}

\findecours

\end{document}
